\documentclass[12pt,a4paper]{article}

\usepackage{mathtools}
\usepackage{mathrsfs}
\usepackage{float}
\usepackage{caption}
\usepackage{tikz}
\usepackage{tikz-cd}
\usetikzlibrary{arrows,positioning,calc,decorations.pathmorphing}

%% Current project v3.070 source-reference bindings; previous macro bodies unchanged.
% Current II-D reference root for PapersMinimal v3.063; mathematical macros otherwise unchanged.
% Macros-v2_188: re-based current paper identities and A--G family; no change to geometry.
% Macros-v2_184: active Paper III-E v1.200/III-F v1.267 binding for framework v3.031; mathematical macros unchanged.
% Macros-v2_183: updated shared cross-reference bindings for Paper1B/III-E/III-F v3.030; mathematical macros unchanged.
% Macros-v2_182: active Paper II-F/V-A cross-references for framework v3.029; definitions unchanged.
% Paper-series v2.381: new I-B torsion/modular results, bibliography and current source roots.
% Paper-series v2.380: new exceptional bibliography and synchronized paper references.
% Paper-series v2.379 current 31-paper macro/aux registry.
% Paper-series v2.378: current 31-paper AUX/bibliography registry updated for Paper I-B v1.2 literature additions.
% Framework v2.371: current-source reference registry refreshed; old theorem labels retained.
% Framework v2.370: current 30-paper source/AUX root rebinding; native mathematical macros unchanged.
% Framework v2.368: synchronized merged 30-root registry; mathematical notation unchanged.
% Framework v2.358: active 30-paper source registry synchronized; existing mathematical macros unchanged.
% Framework v2.357: selective new source-owned certificates; active reference roots updated for five revised papers.
% Framework v2.356: active 30-root registry synchronized to the 30 papers in this release.
% Framework v2.275: exact positive dense-twin word cubature, finite repeatability and eighth-moment classification, nested sampling frames, joint action and POVM-jet exports; preserves v2.274.
% Framework v2.274: filename and root-source synchronization.
% Framework v2.273: restored missing Problem 3 theorems from v2.268; exact exceptional frames, twin geometry, convolution repeatability and projector/occurrence firewalls; merged on v2.272 without overwriting independent branches.
% Framework v2.272: merged v2.270 covariant native-readout mathematics with v2.271 Problem 4 finite-action Hopf hypercharge/holonomy corrections; no universal bulk-gauge uniqueness claim.
% Independent branches: v2.270 readout/guard exports and v2.270 K16 action-consumed charge closure; all prior scoped assumptions retained.
% Framework v2.269: synchronized 30 active AUX roots; mathematics and bibliography preserved from Macros-v2_161.
% Framework v2.266: Paper II-C converse-derived outcome resources; II-E, III-C, IV-A/B, V-A and VI-B downstream mathematical closures; 30-root registry synchronized.
% Framework v2.265: II-C coordinate-stage constitution and exact first-symbol host inheritance; prior mathematical domains retained.
% Framework v2.264: paper II-B reverse/audit and IV-A observable-minimal exchange; synchronized external references to 30 active roots.
% Framework v2.263: Paper I-C adds bridge-controlled root-carry parity and source-relative old-cap cancellation; other root mathematical bodies retained; active 30-root references synchronized.
% Framework v2.262: integrates Paper II-F radius, Casimir and readout results; active 30-root references synchronized.
% Framework v2.261: merged Paper II-B GP2-level/geometric repairs; preserves v2.260 II-D repairs; active references synchronized.
% Framework v2.260: Paper II-D signed RS, finite-symbol typing, and conditional curvature/Jordan repairs; references synchronized.
% Framework v2.259: conditional native SU(2) atlas and explicit later discharge; v2.258 mathematical bodies retained elsewhere.
% Framework v2.258: preserves parallel II-A/dynamics update and adds complete readout/frequency closures.
% Framework v2.257: exact minimal complete position readout and arbitrary-preparation frequency law.
% Framework v2.256: exact native inertia and sharp bounded ordered-source closure; preserves concurrent v2.255 work.
% Framework v2.255: preserves parallel detector integration and adds exact finite readout closures.
% Framework v2.254: Pasterski bibliography and synchronized 30-paper external-reference registry.
% Framework v2.253: typed raw coordinate defect in I-A and synchronized VI-A import; active references synchronized.
% Framework v2.252: synchronized source registry; optional III-G shell-frame calculation.
% Framework v2.251: forward common-action execution, joint source jets and scoped coupled control.
% Framework v2.251: exact full-heat variations, fermion source jets and explicit physical completion frontier.
% Framework v2.251: certified interacting source jets, pole stability and scoped completion proofs.
% Framework v2.251: synchronized active source registry; mathematical notation unchanged.
% Framework v2.251: synchronized active source registry; mathematical notation unchanged.
% Framework v2.251: synchronized active source registry; mathematical notation unchanged.
% Framework v2.251: synchronized source registry; no mathematical notation changes.
% Framework v2.244: synchronized 30-paper AUX registry and rooted phase-source action-order tests.
% Framework v2.242: signed rooted phase lift, complete transverse return and exact family cubic.
% Framework v2.241: proved rooted conormal geometry, conditional cap realization and scoped action tests.
% Framework v2.239: all 30 active source roots synchronized; no mathematical macro changes.
% Framework v2.238: Paper IV A/B/C split and 30-root reference registry.
% Framework v2.237: Paper VII-D technical split and 28-root reference registry.
% Framework v2.236: Paper VI-C technical-volume reference added; all current root bindings advanced.
% Framework v2.235: 26-root AUX and bibliography bindings synchronized; mathematical macros unchanged.
% BEGIN INLINE SOURCE: Paper00-TeXBuildCompatibility-v1_0.tex
% Rendering-only repairs for inherited long appendices and bitmap text fonts.
% Integer counters through 26 retain the original labels; later values use AA,
% AB, ... . This file changes no action coefficient, state or theorem premise.
\RequirePackage{alphalph}
\makeatletter
\let\BHPrimitiveAlph\@Alph
\newalphalph{\BHExtendedAlph}{\BHPrimitiveAlph}{26}
\let\@Alph\BHExtendedAlph
\AtBeginDocument{%
  \@ifpackageloaded{microtype}{\microtypesetup{expansion=false}}{}%
}
\makeatother

% END INLINE SOURCE: Paper00-TeXBuildCompatibility-v1_0.tex

% Framework v2.164: current reference/bibliography profile; mathematical definitions retained.
% Framework v2.162: finite metric, functional Hilbert and scoped Ward integration.
% Framework v2.161: inherits Macros-v2_57.tex exactly except the current external-reference binding.
% Macros v2.57 / framework v2.156: current reference registry only; mathematical definitions unchanged.
% Framework v2.143 / Macros v2.55: current external-reference bindings; mathematical notation unchanged.
% Framework v2.110 / Macros v2.54: current external-reference registry and bibliography; mathematical notation inherited.
% Framework v2.67 / Macros-v2_53: adds Born-action-cell and joint-EH readout symbols; absolute unit realization is metrological; MATCH_{J/scal}^{EH} remains the sole G24 normalization bridge.
% Framework v2.35 / Macros-v2_52: grouped paper numbering, active AUX registry, and Born-v6_20 bibliography binding synchronized; notation unchanged.
% Framework v2.26: active-source AUX and bibliography bindings synchronized; notation unchanged.
% Framework v2.09: reviewed early-reciprocity integration; scoped physical imports.
% Framework v2.08: Born--Euler mathematics and conditional complete-action potential criterion.
% Framework v2.07: active AUX registry includes Paper I-B v1.2; shared bibliography Born-v6_18.
% Macros-v2_46 / framework v2.04: mathematical atlas and downstream realization correction.
% Macros-v2_46 / framework v2.03: scoped structure integration and current dependency bindings.
% Macros-v2_46 / framework v2.01: shared source bindings; prior revision comments are historical.
% v2.42 / framework v1.99: Synchronized RG/AS owner and bibliography bindings.
% v2.41 / framework v1.99: active manuscript AUX bindings and geometric-closure bibliography synchronized.
% Release-v1.80 centralized source synchronization for scalar-circle/comparison-band integration.
% Release v1.79: Hurwitz hierarchy and kinetic-normalized action response integration; physical mass matching remains conditional.
% Release v1.78: approved M8 closure-batch integration; physical and analytic domains retained.
% Current v1.78 registry: all active TeX roots are synchronized below.
% Earlier release comments are historical scope records, not current status.
% Release-v1.77 synchronization: native same-owner Noether closure, generator--flow/history typing, analytic-certification firewalls, scalaron/K16, LF/QCD, quantum-Ward, and RAR boundary-owner transport are registered; shared mathematical notation is unchanged.
% Release-v1.76 synchronization: registry points to the v1.76 active paper roots.
% Release-v1.76 synchronization: registers the local A+Phi+g Noether/constraint closure and dynamics-frontier propagation; shared mathematical notation is unchanged.
% Release-v1.75 synchronization: registry points to the v1.75 active paper roots.
% Release-v1.75 synchronization: integrates the joint-Dirac Einstein bulk real form and the local odd-scalar heat-coefficient action on explicit analytic first-jet/interior domains; v1.74 Yang--Mills closure is retained. Fermionic reality, pure-metric curvature-squared/exact-nonlocal terms, noncanonical boundaries, physical vacuum selection, and actual solution-history tokens remain open.
% Release-v1.74 synchronization: registry points to the v1.74 active paper roots and retains the v1.73 theorem labels.
% Release-v1.74 synchronization: canonical Macros-v2_34.tex; generated local nonlinear Yang--Mills action, constrained connection-field phase space, source-free Gauss momentum map, and Legendre closure are integrated on their proved domains. The remaining non-gauge ACT_cl sectors and actual solution-history selection remain open.
% Release-v1.71 synchronization: closes the former JOINT_x/JOINT_5
% connection-occurrence gap by canonical structure-group extension of the
% nonflat Born connection, faithful associated representation, and the local
% cotangent--Clifford operator constructor.  Actual fluctuation members,
% actions, vacua, solutions, calibrations, and analytic QFT domains remain
% separately typed.  The active-AUX registry is advanced once.
% Release-v1.71 canonical macro and active-AUX registry.
% ======================================================================
% Macros-v2_39.tex -- canonical publication macro file
% Release v1.71 retains the v1.70 null-ray/QCD boundaries and additionally
% removes the duplicate structural connection actuality premise.  It does not
% claim global bundle topology, select a gauge/Higgs member, or promote the
% finite construction to an interacting continuum theory.
% Release v1.68 constructs the finite light-front total-fraction operators and
% parton number/momentum operator measures, proves their simplex and Hausdorff-limit
% theorems, and separates positive cut measures from scheme-renormalized
% light-ray distributions.  It also separates finite real Schur poles, stable
% norm-resolvent/Riesz pole-residue limits, and continued resonance
% pole/residue limits.  The Paper-VII-C scattering package is claim-indexed:
% local QFT and spectral claims import no S-matrix or infrared premise, while
% asymptotic amplitudes, LSZ, S-matrix poles, completeness, and cross sections
% import only their named massive or massless/long-range clauses.
% Release v1.67 factorizes BND48_H into a same-action Hessian/Schur
% descendant and still-open typed boundary-trace, critical-lift, and
% refinement-limit conditions.  It proves the finite-collar Dirichlet-to-
% Neumann law sqrt(y)tanh(L sqrt(y)), its nonuniform critical limit, and the
% finite-return ultra-infrared turnover.  The Paper-I-A Hardy constant mode is
% an exact geometric near-hit, not yet the physical metric boundary mode.
% The current local f(R) Hessian remains massive and therefore still refutes
% a RAR prediction on the minimal action branch.
% Release v1.66 derives the transverse depth-six normalizer bridge from the
% already generated C6 Reynolds image and Paper-I-A binary-dihedral normalizer.
% The bridge is the unique least 2D3-stable scalar-containing realization of
% the faithful quotient 2D3/C6=C2, has trace defect six and hence K=48, but is
% transverse to the nested active carrier.  Exact K=44 and K=52 contraction,
% semigroup, and energy-midpoint identities are registered as structural
% readouts rather than projector events.  The active E48 carrier gap remains
% unchanged.  A same-action boundary-symbol/Calderon intertwiner is isolated
% as the necessary and still unproved owner bridge to the static DtN sort;
% representation equivalence alone neither derives RAR nor creates a K=48
% action junction, pole, occupation, or cosmological transition.
% Release v1.65 derives the native E_48-centered carrier mass gap from the
% GP1--GP3-generated active spectrum, distinguishes its carrier-free interval
% from a particle located at E_48, and withdraws the unsupported promotion to
% a positive continuum Yang--Mills mass gap.  The latter remains a valid
% conditional transfer theorem on an enlarged continuum branch, but it is not
% an active framework prediction or an unfinished native-gap proof.
% Release v1.64 synchronizes Derivation Batches 14--22.  It registers the
% GP1--GP3 heat-profile orbit, same-Hessian electroweak response and pole
% ownership, K16/K48 action-response corrections, perturbation state and
% non-Gaussianity typing, finite-vacuum/gap/QCD limit repairs, factorized
% position-area ownership, the actual static metric response, and the ordered
% Wilson/string-tension limits.  Strict continuum convergence, positive
% physical gap, surviving higher cumulants, and positive tension remain open.
% Release v1.61 closes TIMEWIDTH16 as an independent direct-production datum.
% JUNC_{2|7} continues to own the finite intrinsic microscopic profile/width,
% while the actual same-action background/production history determines its
% physical-time pullback, support, and duration.  The current Paper-I-A--VII-C AUX
% registry is advanced once; shared mathematical notation is unchanged.
% Release v1.60 derives the final-octet primitive-remnant depletion, canonical
% Hopf-current representative, exact Hopf descent to dOmega/(4pi), and
% closed-FLRW a^{-3} neutral-current shape.  It retypes K=48 as the onset
% of the final one-unit geometric transfer and records cDM/RAR as distinct
% physical projections of that common neutral parent while preserving
% occurrence, Hilbert-owner, stress-matching, and static-kernel premises.
% Release v1.59 transports the v1.58 exceptional-shell/K16 closure through
% Papers III-F and VI-B--VII-C.  It replaces Paper VI-B's obsolete mixed-Hessian
% output typing by the third-variation pair vertex and Bogoliubov occupation
% support, derives the post-seed Higgs-pair feedback vertex and its exact
% alpha--delta kernel factors, propagates the unique K16/K48 action-junction
% split, identifies the QCD delta count with the upstream 2O shell atom only
% as a common-scalar numerical compatibility, and narrows THERM16 to dynamical
% approach/beta selection because the finite Gibbs/KMS target is already
% unique on a supplied simple observable algebra once its Hamiltonian and
% inverse temperature are supplied; direct-sum central weights remain typed.
% Release v1.58 registers the exceptional nearest-shell closure, the derived
% transverse K=16 action event, the unique interior action junction, the
% type-correct ordered pair vertex, the thin-layer squeeze law, and the
% radiation-like/matter-like kinematic split.  REG_0, transfer magnitude,
% physical layer width, and thermalization remain explicitly conditional.
% Release v1.57 registers the active-depth-six obstruction, the native
% K=48 carrier desert, the typed 7|9 owner reduction, and the dynamical
% reheating handoff across Papers III-B--VI-B.
% Release v1.56 registers the least-support n_2 inflation dressing of the
% persistent n_7 Laurent response, the analytic global-convexity and positive
% endpoint-measure proofs, the derived inflation core and K=16 coefficient
% freeze, and the no-particle dark-sector split into static RAR plus the
% conditional homogeneous effective-curvature response at reflected K=48.
% There is no n_6 action, scalaron-DM mass, or second Einstein--Hilbert term.
% Release v1.55 registers the protected-octet resolved/residual reflection
% q_sp+Z_W=8, its j<->8-j and K<->64-K scale covariance, and the derived
% rank-one-exit quadratic grade C_1^2/U=C_2 in Papers III-C, III-E, and VIII.
% The reflection is a typed trace/scale theorem, not a time reversal or a
% physical event-occurrence theorem; reheating transfer remains conditional.
% Release v1.54 registers the exact positive-part Laurent indices, the
% active-domain/terminal-endpoint action split, the no-second-EH firewall, and
% the q_DM=6/K_DM=48 scalaron readout typing across Papers III-D/E/F, VI,
% VIII, and IX.
% Release v1.53 registers the common nonnegative Laurent obstacle projection
% for n_1^+ and n_7^+, discharges the three former exit/freeze assumptions,
% and advances the Paper-III-D/E/F, VIII, and IX consumers.  Shared notation
% and bibliography ownership are unchanged.
% Release v1.51 registers the PQTI derivation batch and its updated
% Paper-III-A/C/D/E/F, VIII, and IX consumers.  Shared notation and
% bibliography ownership are unchanged.
% Release v1.50 advances the Paper-III-B AUX target to v1.40 after the
% rigorous audit repair; shared notation and bibliography ownership are unchanged.
% Release v1.49 advances the Paper-II-E/F AUX target to v3.14 after the
% minimal audit repair; shared notation and bibliography ownership are unchanged.
% Release v1.46 synchronizes Paper II-A's minimal frozen constructor-origin
% theorem and its five dependent AUX targets.  Shared notation is unchanged.
% Release v1.45 updates the Paper-II-B AUX target to v2.23 after the
% GP3-rigidity and appendix-placement repair; shared notation is unchanged.
% Release v1.44 advances the centralized Paper-I-A--VII-C AUX registry after the
% Paper-I-A v5.16 protected-completion scope repair.  Every current
% paper imports this one file; no paper-level opt-in or \externaldocument list
% is permitted.  The bibliography remains centralized here through Born-v6_13.
% Shared bibliography, theorem infrastructure, colors, notation alignment,
% and current Paper-I-A--VII-C external-reference registration.
% This file centralizes common imports, bibliography sorting, colors, theorem
% infrastructure, and the notation shared by the complete series.
% v2.9 (release v1.43) made the Paper-I-A--VII-C AUX registry automatic and
% centralized the bibliography and all external-reference ownership.
% v2.8 release-registry refresh (v1.33) registers the rooted generated-parent
% cascade and the POL/REG/free-law/11+4/Born-group closure revisions.  No
% shared mathematical notation or Paper-II-A fundamental principle is changed.
% v2.8 release-registry refresh (v1.32) registers the universal generated-
% continuation revisions in Papers II-A, III-A, III-C, IV, and VIII.  No
% shared mathematical notation or Paper-II-A fundamental principle is changed.
% v2.8 release-registry refresh (v1.30) registers the CRP4/projective-state
% source revisions and their direct registry-only consumers.  No shared
% mathematical notation or Paper-II-A axiom is changed.
% v2.7 release-registry refresh (v1.27) updates only the current external
% Paper-I-A--VII-C basenames after the placement/assumption synchronization; no
% shared mathematical notation or Paper-II-A axiom is changed.
% v2.7 synchronizes the geometric defect-repair revision: Paper-II-A v3.42,
% Paper-III-C v1.26, Paper-III-D v1.27, Paper-III-E v1.32, and Paper-III-F
% v1.34.  It registers the common trace-defect/spectral-filtration theorem,
% its owner-typed response pullback, and the reduced GRED and TERM interfaces.
% No shared mathematical notation is changed.
% v2.6 synchronizes Paper-III-E v1.31's cone-type repair, minimal-response
% classification, and explicit classical action-role boundary, together with
% Paper-III-F v1.32's downstream response-atlas synchronization.  No shared
% mathematical notation is changed.
% v2.5 synchronizes the Paper-III-D v1.26 faithful-trace-duality proof,
% exact ordered-owner joint locus, and typed alternative-presentation
% bookkeeping repair, together with its Paper-III-E and closure-frontier
% references.  No shared mathematical notation is changed.
% v2.4 synchronizes the Paper-III-C v1.25 orbit/process/owner repair, its
% downstream Papers III-B--F, V, VI, VIII, and IX revisions, and Born-v6_13.
% v2.3 synchronizes the Paper-II-A v3.41 witness/support/principal inheritance
% layer, the Paper-III-A v1.24 joint-realizability/readout firewall, all dependent
% host-rate revisions, and the current Paper-I-A--VII-B source registry.
% v2.2 synchronizes the generated active-realization, determinant-depth,
% dimension-weighted hierarchy, and typed cosmological-event interfaces.
% v2.1 synchronizes the exceptional-support interface revisions: Paper I-A
% v5.9, Paper II-A v3.30, Paper II-E/F v3.7, and Paper IV v4.16.
% v2.0 introduces the stage-safe G0 occurrence notation, synchronizes the
% active Paper-I-A--IV source registry (current rigorous Paper-III closure:
% III-A v1.17, III-B v1.22, III-C--E v1.17, III-F v1.14, IV v4.14), and makes
% external-reference self-skip depend on logical paper role rather than a
% versioned job name.
% ======================================================================

% ======================================================================
% macros.tex — Shared macros for the Born Geometry series
% ======================================================================

% ======================================================================
% 1. PACKAGES, BIBLIOGRAPHY, COLORS, AND UNIT DEFINITIONS
% ======================================================================
% Core packages shared across the complete Paper-I-A--VII-C series.  Loading them here makes
% the macro file a genuine series preamble while remaining compatible with existing
% paper preambles that load the same packages before \input.
\usepackage{amsmath}
\usepackage{amssymb}
\usepackage{amsthm}
\usepackage{mathrsfs}
\usepackage{array}
\usepackage{booktabs}
\usepackage{xcolor}
\makeatletter
\@ifpackageloaded{geometry}{}{\usepackage{geometry}\geometry{margin=1in}}
\makeatother
\usepackage{mdframed}

% Shared support packages required by the macros below and by most papers.
\usepackage{ragged2e}
\usepackage{longtable}
\usepackage{tabularx}
\usepackage{enumitem}
\usepackage[most]{tcolorbox}

% Series-wide numeric bibliography convention using the portable BibTeX path.
% All Papers I-A--VII-C use the ordinary \cite interface.
\usepackage[numbers,sort&compress]{natbib}
\providecommand{\printbibliography}{%
  \bibliographystyle{plainnat}%
  \bibliography{Born-v6_59}%
}

% Shared colors.  \providecolor avoids clobbering a paper-specific override.
\providecolor{postulatecolor}{RGB}{230,242,255}
\providecolor{compare}{RGB}{230,245,230}
\providecolor{definitioncolor}{RGB}{230,242,255}

\providecolor{openriskcolor}{RGB}{255,235,235}
\providecolor{bridgecolor}{RGB}{240,240,255}
\providecolor{gold}{RGB}{255,215,0}

% ======================================================================
% 2. CUSTOM ENVIRONMENTS
% ======================================================================
% Common table column types used in all papers.  Guarded so that a paper may
% define its own versions before loading this file.
\makeatletter
\@ifundefined{NC@find@L}{%
  \newcolumntype{L}[1]{>{\RaggedRight\arraybackslash}p{#1}}%
}{}
\@ifundefined{NC@find@C}{%
  \newcolumntype{C}[1]{>{\Centering\arraybackslash}p{#1}}%
}{}
\@ifundefined{NC@find@Y}{%
  \newcolumntype{Y}{>{\RaggedRight\arraybackslash}X}%
}{}
\makeatother

% ======================================================================
% 3. THEOREM STYLES AND ENVIRONMENTS
% ======================================================================
% Custom underlined theorem style (useful for remarks or notes)
\newtheoremstyle{underlined}
  {3pt}   % Space above
  {3pt}   % Space below
  {}      % Body font (empty means normal/roman text, typical for remarks)
  {}      % Indent amount
  {}      % Theorem head font (can be \itshape, \bfseries, etc.)
  {.}     % Punctuation after theorem head
  { }     % Space after theorem head
  {\underline{\thmname{#1}\thmnumber{ #2}}\thmnote{ (#3)}} % Theorem head formatting

% Standard plain theorem style environments.  Each environment is guarded so
% the shared macro file can be loaded after a paper-specific theorem setup.
\theoremstyle{plain}
\makeatletter
\@ifundefined{theorem}{\newtheorem{theorem}{Theorem}[section]}{}
\@ifundefined{lemma}{\newtheorem{lemma}{Lemma}[section]}{}
\@ifundefined{proposition}{\newtheorem{proposition}{Proposition}[section]}{}
\@ifundefined{corollary}{\newtheorem{corollary}{Corollary}[section]}{}
\makeatother

% Definition style environments (typically non-italicized body)
\theoremstyle{definition}
\makeatletter
\@ifundefined{definition}{\newtheorem{definition}{Definition}[section]}{}
\@ifundefined{postulate}{\newtheorem{postulate}{Postulate}[section]}{}
\@ifundefined{closure}{\newtheorem{closure}{Closure}[section]}{}
\makeatother

% Applying the custom underlined style
\theoremstyle{underlined}
\makeatletter
\@ifundefined{remark}{\newtheorem{remark}{Remark}[section]}{}
\@ifundefined{convention}{\newtheorem{convention}{Convention}[section]}{}
\@ifundefined{principle}{\newtheorem{principle}{Principle}[section]}{}
\makeatother

% Framed proof environment.
% Every \begin{proof} ... \end{proof} is automatically framed/shaded.
% This avoids mdframed/trivlist nesting, which can drop the top border after
% theorem-like environments.
\makeatletter
\renewenvironment{proof}[1][\proofname]{%
  \par\addvspace{6pt}%
  \noindent
  \begin{tcolorbox}[%
    enhanced,
    breakable,
    colback=black!4,
    colframe=black!35,
    boxrule=0.35pt,
    arc=1pt,
    left=6pt,
    right=6pt,
    top=4pt,
    bottom=4pt,
    boxsep=0pt,
    before skip=0pt,
    after skip=6pt,
    parbox=false
  ]%
  \pushQED{\qed}%
  \normalfont\itshape #1\@addpunct{.}\enspace
  \normalfont\ignorespaces
}{%
  \unskip\nobreak\hfill\popQED
  \end{tcolorbox}%
  \par\addvspace{2pt}%
}
\makeatother

\mdfdefinestyle{postulatebox}{
  backgroundcolor=yellow!10,
  linecolor=black,
  linewidth=0.4pt,
  skipabove=\baselineskip,
  skipbelow=\baselineskip,
  nobreak=true
}

% ======================================================================
% PROOF-READABILITY HELPERS
% ======================================================================

\newcommand{\ProofPlan}[1]{%
  \par\smallskip
  \noindent\textbf{Proof plan. }#1%
  \par\smallskip
}

\newlist{proofsteps}{enumerate}{2}
\setlist[proofsteps,1]{
  label=\textbf{(\alph*)},
  ref=\alph*,
  leftmargin=2.2em,
  itemsep=0.45\baselineskip,
  topsep=0.35\baselineskip
}
\setlist[proofsteps,2]{
  label=\textbf{(\roman*)},
  ref=\roman*,
  leftmargin=2.0em,
  itemsep=0.25\baselineskip,
  topsep=0.2\baselineskip
}


% ======================================================================
% 4. NUMBER SYSTEMS AND ALPHABETS
% ======================================================================
\providecommand{\RR}{\mathbb{R}} % Real numbers
\providecommand{\PP}{\mathbb{P}} % Projective space / Probability
\providecommand{\CC}{\mathbb{C}} % Complex numbers
\providecommand{\ZZ}{\mathbb{Z}} % Integers
\providecommand{\NN}{\mathbb{N}} % Natural numbers
\providecommand{\HH}{\mathbb{H}} % Quaternions / Hyperbolic space
\providecommand{\LFRayGeom}{\mathsf{LF\!\mbox{-}RAY\!\mbox{-}GEOM}}
\providecommand{\LFRep}{\mathsf{LF\!\mbox{-}REP}}
\providecommand{\LFState}{\mathsf{LF\!\mbox{-}STATE}}
\providecommand{\LFLim}{\mathsf{LF\!\mbox{-}LIM}}
\providecommand{\LFScat}{\mathsf{LF\!\mbox{-}SCAT}}
% Number systems / projective spaces
\providecommand{\CP}{\mathbb{CP}}
\providecommand{\RP}{\mathbb{RP}}

% ======================================================================
% 5. GROUPS AND LIE THEORY
% ======================================================================
\providecommand{\SO}{\ensuremath{\mathrm{SO}}}     % Special Orthogonal group
\providecommand{\SU}{\ensuremath{\mathrm{SU}}}     % Special Unitary group
\providecommand{\Spin}{\ensuremath{\mathrm{Spin}}} % Spin group
\providecommand{\CS}{\ensuremath{\mathrm{CS}}}     % Chern-Simons (capitalized)
\providecommand{\cs}{\ensuremath{\mathrm{cs}}}     % Chern-Simons (lowercase)

% ======================================================================
% 6. COMMON MANIFOLDS AND TOPOLOGY
% ======================================================================
\providecommand{\Sone}{S^1}                 % Standard circle

% ======================================================================
% 7. MATH OPERATORS (Declared via \DeclareMathOperator)
% ======================================================================
\DeclareMathOperator{\End}{End}   % Endomorphisms
\DeclareMathOperator{\Aut}{Aut}   % Automorphisms
\DeclareMathOperator{\spec}{spec} % Lowercase spectrum
\DeclareMathOperator{\tr}{tr}     % Lowercase trace (if you already have \Tr)
\DeclareMathOperator{\supp}{supp} % Support
\DeclareMathOperator{\TrA}{Tr_{A}} % Partial trace over subsystem A
\DeclareMathOperator{\TrB}{Tr_{B}} % Partial trace over subsystem B
% Lie/algebra operators
\DeclareMathOperator{\Ad}{Ad}
\DeclareMathOperator{\Sym}{Sym}
\DeclareMathOperator{\su}{su}
% ======================================================================
% 8. STANDARD OPERATOR-LIKE SYMBOLS (Safe definitions via \providecommand)
% ======================================================================
\providecommand{\vol}{\operatorname{vol}}   % Volume form
\providecommand{\Vol}{\operatorname{Vol}}   % Total volume
\providecommand{\Ric}{\operatorname{Ric}}   % Ricci tensor
\providecommand{\Rm}{\operatorname{Rm}}     % Riemann tensor
\providecommand{\Tr}{\operatorname{Tr}}     % Uppercase trace
\providecommand{\Spec}{\operatorname{Spec}} % Uppercase spectrum
\providecommand{\diag}{\operatorname{diag}} % Diagonal matrix
\providecommand{\rank}{\operatorname{rank}} % Rank
\providecommand{\scal}{\mathcal R} % Scalar curvature
\providecommand{\Cl}{\operatorname{Cl}}     % Clifford algebra
\providecommand{\diam}{\operatorname{diam}} % Diameter

% ======================================================================
% 9. FUNCTIONS, TRIGONOMETRY, AND CALCULUS
% ======================================================================

% Codifferential naming conventions
% (Avoid \ddag: commonly defined by physics.sty and symbol packages)
\providecommand{\codiff}{\mathrm{d}^{\dagger}} % Codifferential operator

% Quantum mechanics bra-ket notation (needed for basis tables)
\providecommand{\ket}[1]{\lvert #1 \rangle}
\providecommand{\bra}[1]{\langle #1 \rvert}
\providecommand{\braket}[2]{\langle #1 \mid #2 \rangle}

% ======================================================================
% 10. PHYSICAL SCALES AND CONSTANTS (Paper 1c / Born Geometry)
% ======================================================================
\providecommand{\Rp}{\ensuremath{R_{\mathrm P}}}     % matched common EH curvature scale (not defined from observed G)
\providecommand{\hBorn}{\ensuremath{h_{\mathrm B}}} % physical representative 2\pi\hbar of one primitive Born scalar action cell
\providecommand{\UEHjoint}{\ensuremath{U_{J}}} % joint-Dirac EH curvature-coordinate scale 8 kappa_J
\providecommand{\Gread}{\ensuremath{G_{\rm scal}}} % scalar-action coordinate c_op^3/(h_B U); physical matching explicit
\providecommand{\GreadJoint}{\ensuremath{G_{J,\mathrm{raw}}}} % raw ordinary-trace coordinate = G_scal/96
\providecommand{\Rds}{R_{\mathrm{dS}}}    % de Sitter Curvature
\providecommand{\Mpl}{M_{\mathrm{Pl}}}    % Planck mass
\providecommand{\lp}{\ell_{\mathrm{Pl}}}  % Planck length
\providecommand{\Pmax}{P_{\mathrm{max}}}  % Maximum momentum
\providecommand{\Mplbar}{\bar{M}_{\mathrm{Pl}}} % Reduced Planck mass
\providecommand{\eff}{\mathrm{eff}}       % Subscript for effective quantities
\providecommand{\Cosmo}{\ensuremath{\Lambda}}                % Cosmological constant (use \Cosmo in source)
\providecommand{\LCDM}{\ensuremath{\Cosmo\mathrm{CDM}}}  % \LambdaCDM (typesets with \Cosmo)
\providecommand{\Lcut}{\ensuremath{\Lambda_{\mathrm{cut}}}}  % Cutoff scale (brace when used as script)
\providecommand{\Rcs}{\ensuremath{R_{\mathrm{cs}}}}      % Configuration-space radius
\providecommand{\CosmoEff}{\ensuremath{\Cosmo_{\eff}}}   % Effective cosmological constant

% Disambiguation: \Rmom is the momentum-space (S^3) radius. 
% Reserve plain R for the 4D curvature scalar in f(R).
\providecommand{\Rmom}{\ensuremath{R_{\mathrm{mom}}}} 
\providecommand{\Punit}{P_{0}} % Momentum unit set by \Rmom (= \hbar/\Rmom)
\providecommand{\Rvec}{\ensuremath{R_{\mathrm{vec}}}} % invariant vector-space radius in the reduced round sector
% Specific numeric values and expressions
\providecommand{\Vlogcap}{L_{\mathrm{Hopf}}^{(1)}}
\providecommand{\Hopf}{\operatorname{Hopf}}
% Scalar phase (S^1) sector + universal cover notation (Paper I-A)
\providecommand{\tper}{T_\star}      % Time period/scale
\providecommand{\Eper}{E_\star}      % Energy period/scale
\providecommand{\ttilde}{\tilde t}   % Universal cover time
\providecommand{\Etilde}{\tilde E}   % Universal cover energy

% ======================================================================
% 10A. SERIES-WIDE LOGICAL / NOTATIONAL DISAMBIGUATIONS
% ======================================================================
% Use these symbols to avoid collisions among: (i) the Paper--III
% no-insertion/geometric-closure class, (ii) Paper--I-A trimer diagonals,
% (iii) Paper--V-A/V-C family-space diagonals, (iv) scalar-index running,
% and (v) protected shell-clock labels.
\providecommand{\Rprim}{\ensuremath{\mathsf R_{\mathrm{prim}}^{\Sigma}}} % primitive finite structural readouts
\providecommand{\Rloop}{\ensuremath{\mathsf R_{\circ}^{\Sigma}}} % closed-loop readout role
\providecommand{\Rbdry}{\ensuremath{\mathsf R_{\partial}^{\Sigma}}} % boundary-conditioned readout role
\providecommand{\Rgen}{\ensuremath{\mathsf R_{\mathrm{fin}}^{\Sigma}}} % finite generated closure of primitive roles
\providecommand{\PhysOcc}{\operatorname{PhysOcc}}
\providecommand{\Actual}{\operatorname{Actual}}
\providecommand{\GeomLic}{\operatorname{GeomLic}}
\providecommand{\InvDist}{\operatorname{InvDist}}
\providecommand{\Dcl}{\ensuremath{\mathsf{D}_1}}      % unmarked positive Dirac-even closed spectral-readout form class
\providecommand{\Dheat}{\ensuremath{\mathsf D_{\mathrm{heat}}}} % one-parameter heat-semigroup subclass
\providecommand{\DtriI}{\ensuremath{D^{\mathrm{tri}}_1}} % Paper--I-A cyclic-trimer diagonal
\providecommand{\DtriII}{\ensuremath{D^{\mathrm{tri}}_2}}
\providecommand{\DfamI}{\ensuremath{D^{\mathrm{fam}}_1}} % Family-space diagonal selector
\providecommand{\DfamII}{\ensuremath{D^{\mathrm{fam}}_2}}
\providecommand{\alpharun}{\ensuremath{\alpha_{\mathrm{run}}}} % scalar spectral-index running
\providecommand{\Ksh}{\ensuremath{K_{\mathrm{sh}}}} % protected shell-clock coordinate
\providecommand{\ddef}{\ensuremath{\delta_{\mathrm{def}}}} % universal defect scalar after Paper--III realization

% ======================================================================
% 10B. SERIES-WIDE LOGICAL STATUS TAGS
% ======================================================================
% Upright small-caps wrapper for audit/status labels.  It prevents LaTeX from
% trying to use unavailable small-caps italic font shapes inside theorem bodies
% such as TU/lmr/m/scit.
\providecommand{\AuditLabel}[1]{}
\renewcommand{\AuditLabel}[1]{{\normalfont\scshape #1}}

% These tags are bookkeeping labels for consistency audits. They do not
% create new mathematical hypotheses; they record how a datum is being used
% in the proof chain.
\providecommand{\StatusPrimitive}{}
\renewcommand{\StatusPrimitive}{\AuditLabel{Primitive}}
\providecommand{\StatusDerived}{}
\renewcommand{\StatusDerived}{\AuditLabel{Derived}}
\providecommand{\StatusImported}{}
\renewcommand{\StatusImported}{\AuditLabel{Imported}}
\providecommand{\StatusConditional}{}
\renewcommand{\StatusConditional}{\AuditLabel{Conditional}}
\providecommand{\StatusBridge}{}
\renewcommand{\StatusBridge}{\AuditLabel{Bridge-derived}}
\providecommand{\StatusComparison}{}
\renewcommand{\StatusComparison}{\AuditLabel{Comparison-only}}

% ======================================================================
% 10C. SERIES-WIDE CLAIM-WORDING AUDIT LABELS
% ======================================================================
% These are textual audit labels used to keep modal language consistent
% across the series.  They introduce no new mathematical assumptions.
\providecommand{\ClaimForced}{}
\renewcommand{\ClaimForced}{\AuditLabel{forced}}
\providecommand{\ClaimUnique}{}
\renewcommand{\ClaimUnique}{\AuditLabel{unique}}
\providecommand{\ClaimCanonical}{}
\renewcommand{\ClaimCanonical}{\AuditLabel{canonical}}
\providecommand{\ClaimMinimal}{}
\renewcommand{\ClaimMinimal}{\AuditLabel{minimal}}
\providecommand{\ClaimClosed}{}
\renewcommand{\ClaimClosed}{\AuditLabel{closed}}
\providecommand{\ClaimPrediction}{}
\renewcommand{\ClaimPrediction}{\AuditLabel{prediction}}

% ======================================================================
% 10D. SERIES-WIDE THEOREM-CHAIN AUDIT LABELS
% ======================================================================
% These labels mark the direction of load-bearing theorem chains. They are
% audit annotations only; they introduce no new mathematical assumptions.
\providecommand{\ChainInput}{}
\renewcommand{\ChainInput}{\AuditLabel{Chain input}}
\providecommand{\ChainOutput}{}
\renewcommand{\ChainOutput}{\AuditLabel{Chain output}}
\providecommand{\ChainGuard}{}
\renewcommand{\ChainGuard}{\AuditLabel{Chain guard}}
\providecommand{\ChainNoBackflow}{}
\renewcommand{\ChainNoBackflow}{\AuditLabel{No backflow}}

% ======================================================================
% 10E. SERIES-WIDE LOCAL THEOREM-AUDIT LABELS
% ======================================================================
% These labels mark a statement-level dependency audit around individual
% theorem/definition nodes. They are bookkeeping labels only.
\providecommand{\LocalAudit}{}
\renewcommand{\LocalAudit}{\AuditLabel{Local audit}}
\providecommand{\UsesOnly}{}
\renewcommand{\UsesOnly}{\AuditLabel{Uses only}}
\providecommand{\Exports}{}
\renewcommand{\Exports}{\AuditLabel{Exports}}
\providecommand{\DoesNotUse}{}
\renewcommand{\DoesNotUse}{\AuditLabel{Does not use}}

% ======================================================================
% 10F. SERIES-WIDE NUMERICAL-ROW AUDIT LABELS
% ======================================================================
% These labels classify individual rows in numerical, benchmark, status,
% and comparison tables.  They are audit annotations only and do not change
% any theorem, value, formula, or comparison number.
\providecommand{\RowAudit}{}
\renewcommand{\RowAudit}{\AuditLabel{Numerical-row audit}}
\providecommand{\RowTheorem}{}
\renewcommand{\RowTheorem}{\AuditLabel{Theorem-fixed row}}
\providecommand{\RowBridge}{}
\renewcommand{\RowBridge}{\AuditLabel{Bridge-fixed row}}
\providecommand{\RowRepresentative}{}
\renewcommand{\RowRepresentative}{\AuditLabel{Representative row}}
\providecommand{\RowComparison}{}
\renewcommand{\RowComparison}{\AuditLabel{Comparison row}}
\providecommand{\RowCalibrated}{}
\renewcommand{\RowCalibrated}{\AuditLabel{Calibrated row}}
\providecommand{\RowEstimate}{}
\renewcommand{\RowEstimate}{\AuditLabel{Estimate row}}
\providecommand{\RowOpen}{}
\renewcommand{\RowOpen}{\AuditLabel{Open row}}
\providecommand{\RowNoSelector}{}
\renewcommand{\RowNoSelector}{\AuditLabel{No selector}}

% ======================================================================
% 10G. TABLE HELPERS
% ======================================================================
% Table helpers used by Paper V-C.  The unused readout-label preset has been
% removed; only active table helpers remain.
\providecommand{\gpsection}[1]{%
  \addlinespace[0.55em]
  \multicolumn{3}{@{}l}{\textsc{#1}}\\[-0.25em]
  \midrule
}
\providecommand{\mathcell}[1]{%
  \(\begin{gathered}#1\end{gathered}\)%
}

\newcommand{\DerBy}[1]{%
  \begin{tabular}[t]{@{}l@{}}#1\end{tabular}%
}

\newcommand{\TableNotes}[1]{%
  \par\medskip
  \begin{minipage}{0.94\linewidth}
  \footnotesize\emph{Notes.} #1
  \end{minipage}\par
}
% ======================================================================
% 10H. PAPER VI-A / QCD NOTATION
% ======================================================================
% Shared symbols used by Paper VI-A and downstream QCD/Yang--Mills handoffs.
\providecommand{\Pqcd}{\mathfrak P_{\rm VIII}^{\rm QCD}}
\providecommand{\Iqcd}{\mathfrak I_{\rm QCD}^{(7)}}
\providecommand{\muQCD}{\mu_{\rm QCD}^{\rm int}}
\providecommand{\LambdaQCD}{\Lambda_{\rm QCD}}
\providecommand{\LambdaBH}{\Lambda_{\rm BH}^{\rm QCD}}
\providecommand{\HQCD}{\mathcal H_{\rm QCD}^{\rm BH}}
\providecommand{\HQCDfin}{H_{\rm QCD}^{\rm fin}}
\providecommand{\Hraw}{\mathcal H_{\rm raw}}
\providecommand{\HphysQCD}{\mathcal H_{\rm phys}^{\rm QCD}}
\providecommand{\Vcol}{V_C}
\providecommand{\Ps}{P_{\rm sing}}
\providecommand{\PiOct}{\Pi_8}
\providecommand{\MSbar}{\overline{\rm MS}}
\providecommand{\YM}{\mathrm{YM}}
\providecommand{\Ktri}{K_{\triangle}}
\providecommand{\KYM}{K_{\YM}^{\rm BH}}
\providecommand{\aBH}{a_{\rm BH}}
\providecommand{\epsBH}{\varepsilon_{\rm BH}}
\newcommand{\RowConstructive}{\textnormal{\textsc{Constructive}}}
\newcommand{\RowObligation}{\textnormal{\textsc{Bridge obligation}}}
\newcommand{\RowTarget}{\textnormal{\textsc{Target row}}}
\newcommand{\RowValidation}{\textnormal{\textsc{Validation row}}}
\providecommand{\slashed}[1]{\not{#1}}

% ======================================================================
% 10I. PAPER VI-B / CONTINUOUS-QFT NOTATION
% ======================================================================
% Shared symbols used by Paper VI-B and any downstream continuous-QFT or
% Yang--Mills handoff paper.  All definitions are guarded so individual
% papers may override them locally before loading this file if needed.
\providecommand{\A}{\mathfrak A}
\providecommand{\Hphys}{\mathcal H_{\rm phys}}
\providecommand{\QCD}{\mathrm{QCD}}
\providecommand{\BH}{\mathrm{BH}}
\providecommand{\loc}{\mathrm{loc}}

\providecommand{\RBH}{\mathcal R_{\rm BH}}
\providecommand{\DG}{\mathfrak D_G}
\providecommand{\PW}{\mathcal{P\!W}} % Peter--Weyl spectral-symbol subspace
\providecommand{\Toe}{\operatorname{Toe}} % Toeplitz compression/representative
\providecommand{\Hol}{\operatorname{Hol}} % Holonomy

% ======================================================================
% 11. LATE-LOADED PACKAGES (Must be at the end of the preamble)
% ======================================================================
% Load late packages only if a paper has not already loaded them, avoiding
% option clashes in paper-specific preambles.
\makeatletter
\@ifpackageloaded{hyperref}{}{\usepackage{hyperref}}
\makeatother
% --- Series-wide disambiguations (Pass 3) ---
\providecommand{\diracdelta}{\delta_{\rm D}} % Dirac delta distribution (delta reserved for hierarchy)
\providecommand{\Hg}{H_{\mathrm g}} % Hodge harmonic projector (keep H for Hubble)
\providecommand{\Gr}{G_{\mathrm r}} % Green operator (keep G for Newton)
\providecommand{\FD}{\mathrm{D}} % functional/variational derivative symbol (delta reserved)
\providecommand{\Id}{\mathbf{1}} % identity/Kronecker tensor (avoid \delta_{ij})


% BEGIN INLINE SOURCE: Paper00-ActiveExternalReferences-v2_130.tex
% Framework v2.191: current reference binding; existing definitions retained.
% Framework v2.190: current22-root registry with native clock-evolution proofs.
% Framework v2.189: exact current22-root external registry, including retained v2.188 hardening owners.
% Framework v2.187: current registry synchronization only; source release v2.186 preserved.
% Framework v2.186: authoritative current 22-paper registry.
% Framework v2.182: all current 22 logical paper AUX bindings.
% Framework v2.170: current I-B v1.43 external AUX binding.
% Framework v2.167: exact current 22-root registry after CP--BL2 integration.
% Framework v2.166: exact current 22-root compact-release AUX registry.
% Framework v2.165: authoritative current root bindings.
% Framework v2.164: exact current 22 manuscript AUX bindings.
% Framework v2.162: finite metric, functional Hilbert and scoped Ward integration.
% Framework v2.161: current 22 manuscript AUX bindings; source only.
% Framework v2.156 / registry v2.113: synchronized current 22 manuscript roots; source-only registry.
% Framework v2.143 / registry v2.111: active manuscript AUX basenames; no generated AUX or PDF included.
% ======================================================================
% Central external-reference imports, numbers only, no PDF jumping
% ======================================================================

% xr-hyper is loaded centrally so that papers may retain their existing
% \usepackage{xr-hyper} lines without redefining \externaldocument after this
% registry is installed.  If plain xr was loaded earlier, retain it.
\makeatletter
\@ifpackageloaded{xr-hyper}{}{%
  \@ifpackageloaded{xr}{}{\usepackage{xr-hyper}}%
}
\makeatother

% Make external-document registration idempotent.  The registry identifies a
% logical paper by its Paper<number><optional letter> prefix.  Prefix-qualified
% imports remain independent.  Current paper roots contain no local
% \externaldocument calls; this guard also makes accidental duplicates harmless.
\ExplSyntaxOn
\prop_new:N \g_bh_external_refs_prop
\cs_new_eq:NN \BHOriginalExternalDocument \externaldocument

\cs_new_protected:Npn \bh_external_document:nnnn #1#2#3#4
 {
  \tl_set:Nn \l_tmpa_tl {#3}
  \regex_replace_once:nnN
   { \A (?: xref/ )? (Paper[0-9]+[A-G]?) .* \Z }
   { \1 }
   \l_tmpa_tl
  \tl_set:Nx \l_tmpb_tl { #1 | \tl_use:N \l_tmpa_tl }
  \prop_if_in:NVTF \g_bh_external_refs_prop \l_tmpb_tl
   {
    \PackageInfo{BornHopf}
     {Skipping~duplicate~external~refs~for~#3}
   }
   {
    \prop_gput:NVn \g_bh_external_refs_prop \l_tmpb_tl {#3}
    \IfNoValueTF{#4}
     { \BHOriginalExternalDocument[#1][#2]{#3} }
     { \BHOriginalExternalDocument[#1][#2]{#3}[#4] }
   }
 }

\RenewDocumentCommand{\externaldocument}{O{} O{} m o}
 { \bh_external_document:nnnn {#1}{#2}{#3}{#4} }

\NewDocumentCommand{\BHExternalRefs}{m}
 {
  % Do not import the aux file of the logical paper currently being compiled,
  % even when the job and registry use different version suffixes.
  \tl_set:Nx \l_tmpa_tl { \c_sys_jobname_str }
  \regex_replace_once:nnN
   { \A (Paper[0-9]+[A-G]?) .* \Z }
   { \1 }
   \l_tmpa_tl
  \tl_set:Nn \l_tmpb_tl { #1 }
  \regex_replace_once:nnN
   { \A (Paper[0-9]+[A-G]?) .* \Z }
   { \1 }
   \l_tmpb_tl
  \str_if_eq:eeTF { \tl_use:N \l_tmpa_tl } { \tl_use:N \l_tmpb_tl }
   {
    \PackageInfo{BornHopf}
     {Skipping~external~refs~for~#1~because~this~is~the~current~document}
   }
   {
    \IfFileExists{xref/#1.aux}
     { \externaldocument[][nocite]{xref/#1}[#1.pdf] }
     {
      \IfFileExists{#1.aux}
       { \externaldocument[][nocite]{#1} }
       {
        \PackageInfo{BornHopf}
         {External~aux~file~#1.aux~not~found;~no~labels~imported}
       }
     }
   }
 }
\ExplSyntaxOff

% External references are automatic.  This file is the sole owner of both the
% bibliography and the current Paper-I-A--VII-D AUX registry (30 roots including I-C).  Missing auxiliaries
% are informational on the first pass and are populated by the source-only
% two-pass series build.  Basenames deliberately omit temporary download/copy
% suffixes such as ``(1)'' or ``-1''.
  % Paper I-A: round S3 compendium.
  \BHExternalRefs{Paper1A-Geometry-v5_181}

  % Paper I-B: complete exceptional geometry and finite spectral structures.
  \BHExternalRefs{Paper1B-ExceptionalGeometry-v1_25}
  % Paper I-D: technical companion of I-A through I-C.
  \BHExternalRefs{Paper1D-TechnicalMaterial-v1_69}
  % Paper I-C: reciprocal M8 compendium.
  \BHExternalRefs{Paper1C-M8Geometry-v1_146}

\ifdefined\BHMathOnly
  % Papers I-A and I-B import only the mathematical atlas labels.
\else
  % Paper II: five physical-mathematical main papers plus one technical proof volume.
  \BHExternalRefs{Paper2A-GeometricConstitution-v3_227}
  \BHExternalRefs{Paper2B-MomentumCarrier-v2_147}
  \BHExternalRefs{Paper2C-ReciprocalLocalization-v2_176}
  \BHExternalRefs{Paper2D-SymmetrySpectra-v2_145}
  % The former II-F export labels remain in II-E; current II-F is the technical volume.
  \BHExternalRefs{Paper2E-KinematicExports-v3_140}
  % Paper II-F: complete technical proof and calculation owner for II-A--II-E.
  \BHExternalRefs{Paper2F-TechnicalMaterial-v1_71}

  % Paper III: Dynamics and Intrinsic Geometric Constitution (six components).
  \BHExternalRefs{Paper3A-GeneratedDynamics-v1_181}
  \BHExternalRefs{Paper3B-SpectralAction-v1_192}
  \BHExternalRefs{Paper3C-ProtectedHierarchy-v1_195}
  \BHExternalRefs{Paper3D-RunningResponse-v1_167}
  \BHExternalRefs{Paper3E-GravityAction-v1_203}
  \BHExternalRefs{Paper3F-CompleteDynamics-v1_279}
  % Paper III-G: complete preserved technical proof owner.
  \BHExternalRefs{Paper3G-TechnicalMaterial-v1_57}

  % Paper IV-A: generated finite quantum mechanics.
  \BHExternalRefs{Paper4A-FiniteQuantumMechanics-v4_217}
  % Paper IV-B: complete finite quantum processes and physical response.
  \BHExternalRefs{Paper4B-QuantumResponse-v1_55}
  % Paper IV-C: extended proofs and finite technical calculations.
  \BHExternalRefs{Paper4C-TechnicalMaterial-v1_50}

  % Paper V-A: The Standard Model.
  \BHExternalRefs{Paper5A-FiniteMatterGeometry-v3_218}

  % Paper V-B: Finite Spectrum and Flavor Closure.
  \BHExternalRefs{Paper5B-TerminalSpectrum-v5_202}

  % Paper V-C: Finite Interaction Boundaries of the Standard Model.
  \BHExternalRefs{Paper5C-InteractionResponse-v2_184}
  % Paper V-D: preserved technical proof and calculation owner.
  \BHExternalRefs{Paper5D-TechnicalMaterial-v1_58}

  % Paper VI-A: Gravity and Closed Cosmology.
  \BHExternalRefs{Paper6A-GravityCosmology-v2_268}

  % Paper VI-B: Geometric Dark Sector and Cosmological Readouts.
  \BHExternalRefs{Paper6B-CosmologyReadouts-v1_215}

  % Paper VI-C: complete technical proof and comparison owner.
  \BHExternalRefs{Paper6C-TechnicalMaterial-v1_53}

  % Paper VII-A: Finite QCD Kernel and Nonperturbative Closure.
\BHExternalRefs{Paper7A-QCDKernel-v1_183}

  % Paper VII-B: Hadronic, Thermal, and Nuclear QCD Readouts.
  \BHExternalRefs{Paper7B-QCDReadouts-v1_167}

  % Paper VII-C: Continuous Quantum Field Theory.
  \BHExternalRefs{Paper7C-ContinuousQFT-v2_208}

  % Paper VII-D: complete technical material for VII-A through VII-C.
  \BHExternalRefs{Paper7D-TechnicalMaterial-v1_60}

\fi

% END INLINE SOURCE: Paper00-ActiveExternalReferences-v2_130.tex


\hypersetup{hidelinks}


\usepackage{xr-hyper}

\title{\textbf{Born--Hopf Geometry III-B: Spectral Operators and the Common Action Class}\\[0.5em]
\large Cartan geometry, spectral triples, and heat-coefficient ancestry}

\author{Peter Kugelmann\thanks{\texttt{pkugelmann@proton.me}}}
\date{20 November 2026}
\setlength{\emergencystretch}{3em}

\begin{document}
\maketitle
\begin{abstract}
Paper III-B treats the generated self-adjoint spectral owner, the
Cartan and finite-module constructions, the spectral triple and the
loop-generated additive profile class used in a common action. Heat
coefficients and normalization statements are distinguished from the
physical appearance of a selected local source action. Optional spectra,
operator calculations, and countermodels are preserved in Paper~III-G.
\end{abstract}
\tableofcontents
\clearpage

\section{Constituted operators and compact spectral input}
\label{paper3b:sec:organized-scope}
\subsection{Introduction and Paper-III series allocation}
\label{p3b:sec:series-allocation}
Paper~III-B begins after the compact structural momentum-role seed has already
been constituted.  The imported theorem statement is
\[
 \boxed{\mathcal H_p\cong\mathcal H_1,\qquad m_p=1},
\]
from Paper~II-B.  Paper~II-C intrinsic \(E_1\)-evaluation together with the
Paper-II-A GP2 theorem supplies the Born/GP2-generated reciprocal round role, the
universal relation-free joint product, and the Born exchange involution.
Paper~III-A reconstructs the same certification statuses on the optional
full-profile metric-curry domain.  The tags \(\mathrm{BR}_{\rm ind}\) and
\(\mathrm{IND}_{xp}\) therefore denote derived certifications and are not
additional hypotheses of Paper~III-B.  The derived mathematical reciprocal
position family is
\[
 \mathsf K(m_x,1),
\]
where the primitive Paper-II-C frontier has \(m_x=1\).  A later occurring
frontier history \((\mathrm{HIST}_x^{\rm fin})\) derives
\(m_x(\widetilde t)\); it is not a free integer trajectory.  On a dynamical
branch that retains the frozen momentum carrier, persistence of the already
derived value is the conclusion \((\mathrm{FRZ}_p)\) of Paper~II-A's
chronological-persistence theorem, not a separate interface premise.

Whenever a dimensionful configuration radius is written, this paper inherits
Paper~III-A's reciprocal scale calibration
\[
 (\mathrm{CAL}_{xp}):\qquad \Rcs=\alpha_{xp}\Rmom.
\]
Whenever the scalar Poincar\'e class is represented by the dimensionful phase
cell \(2\pi\hbar\), this paper inherits
\[
 (\mathrm{CAL}_{\hbar}):\qquad
 \left[\frac{\omega_{\rm sc}}{2\pi\hbar}\right]=c_{\rm sc}.
\]
Whenever the primitive scalar turn is given a physical duration, this paper
also inherits Paper~III-A's clock calibration
\[
 (\mathrm{CAL}_{t}):\qquad
 T_t>0
 =\text{physical duration of one positive primitive turn}.
\]
None of these calibrations is generated by Born reciprocity itself.  In
particular, the GP2 derivation of the Born product and involution fixes
reciprocal role type and exchange structure but determines neither the
dimensionful ratio \(\alpha_{xp}\), the physical phase unit, nor the clock
duration \(T_t\).

\subparagraph{Imported G0 actualization boundary.}
All physical interfaces below specialize Paper~II-A's stage-safe G0.  Each
asserted occurrence requires a contingent token of a geometrically licensed,
claim-relative invariant distinction with result-independent provenance; its
tag alone supplies no witness, and no instance reopens frozen upstream
geometry.  A new role or standard enters only through a declared downstream
frontier.  No independent G0-MOM witness exists; later momentum language uses
the co-realized carrier, while a particular process, test, or readout remains
a token-level G0 claim on its stated occurrence domain.
The frontier occurrences \((\mathrm{HIST}_x^{\rm fin})\) and
\(\mathsf{SF}_\Lambda\), and the
actual closed-process member represented in the derived form class \(\Dcl\), are G0
occurrence certificates.  The interfaces \((\mathrm{CAL}_{xp})\),
\((\mathrm{CAL}_{\hbar})\), and \((\mathrm{CAL}_{t})\) are G0
calibrations: they assign realized standards to already existing invariants
and do not add structural content.  G0 selects neither the process law nor a
calibration value and does not supply heat-kernel regularity; those remain
occurrence/calibration data and a separate analytic hypothesis.

The task of Paper~III-B is therefore dynamical rather than constitutional:
first classify a reciprocal off-shell response around the fixed seed, then
consume the exhaustive loop role to derive \(\Dcl\), and finally classify its
possible spectral process members and generated decorations.  The generic word/tower family remains auxiliary mathematics.  Absolute
Paper-I-A coordinate values and completion-dependent commutator scales still
specialize the unresolved completion family to \(\mathsf C_{\rm PI}\),
whereas the normalized protected even algebra is inherited already on
\((\mathrm{ND}_{\rm sp})\).  Hurwitz/trimer results require an actualized
\(\mathsf{Supp}_{2T}\) core-support role.

\subparagraph{Completion and support notation.}
Throughout Paper~III-B, ``assume \(\mathsf C_{\rm PI}\)'' means choose the
Paper-I-A point of
\(\mathfrak M_{\rm coord}/\!{\sim}\cong\mathbb R_{\ge0}^3\); it is retained
only for conclusions that use absolute completion data.  For the normalized
protected Clifford chain we instead write
\[
 (\mathrm{ND}_{\rm sp}):\qquad \varrho_{\rm sp}>0,
\]
importing Paper~II-C's nondegenerate protected-coordinate stratum.  On this
whole stratum the normalized even generators satisfy exactly
\(Y_iY_4=C_iC_4\), so the represented protected even algebra is the same
\(\Cl(0,3)\) and no Paper-I-A completion representative is needed for that
algebraic role.  Thus \([\mathsf C]\in\mathbb R_{\ge0}^3\) is an unresolved
realization selector only for claims that use raw completion-dependent
coordinate, radius, commutator, or ambient-lift data; \((\mathrm{ND}_{\rm sp})\)
alone suffices for the normalized protected even algebra.  Likewise, ``assume
\(\mathsf{Supp}_{2T}\)'' means that
either the Paper~I-B exceptional-core constructor has an actual G0 application
or a generated bridge connects an occurring metric point-resolution support
to its domain.  The output is then co-realized without a second witness, and
all labelled Hurwitz vertices, diagonals, and trimers are representatives of
the retained unmarked orbit.  Point resolution alone proves finite metric
support, not this exceptional-domain identification.

\subparagraph{GP allocation.}
Throughout this paper GP2 means Exact/Faithful constitution of generated
content; GP3 means Irreducible/Least reduction of residual structured
redundancy.  GP3 is never used here to minimize a truncation label or to
supply an ungenerated numerical parameter.  Every later least, leading-order,
or endpoint statement is instead typed by the casewise GP2/GP3 theorem and
must prove its own structural comparison, vanishing, variational, or
exact-fidelity antecedent.

\subparagraph{Nonoperative selector alternatives.}
The \(\mathsf P_{\rm sli}\), coordinate-ceiling, and fixed-action
routes to \(m=1\) are technical diagnostic alternatives documented in
Paper~III-G, Section~\ref{p3b:app:legacy-selector}.  They are not antecedents of the
GP2 momentum-seed derivation or of a physical momentum-carrier occurrence.

\subparagraph{Dependency stratification.}
The native III-B constructions use the constituted geometry of Papers~I-A--I-C,
the Paper-II carrier and reciprocal interfaces, and the dynamics of Paper~III-A.
Results requiring later hierarchy, index-locking, active-density, or gravitational
synthesis are forward inputs and are not used to establish the native action theorems.

\subparagraph{Selector-status stratification and later discharge.}
For series bookkeeping, III-B separates four logical classes: generated or
co-realized structure; G0 occurrence tokens such as an actual frontier or
closed process; calibrations \((\mathrm{CAL}_{xp})\),
\((\mathrm{CAL}_{\hbar})\), and \((\mathrm{CAL}_{t})\); and unresolved
realization/bridge data such as a raw coordinate completion, a generic active
density, or a same-process/reduction comparison.  An unresolved III-B family
may be discharged later only by independently constituted data that genuinely
distinguish its members.  In that case its later status is
\(\mathrm{DERIVED\text{-}LATER}\), while its status within the III-B dependency prefix remains
unresolved:
\[
 \boxed{
 \text{unresolved at III-B}
 \;\longrightarrow\;
 \text{independent later discriminator}
 \;\longrightarrow\;
 \mathrm{DERIVED\text{-}LATER}.}
\]
No later success is read backwards as an upstream selector.

\label{p3:part:spectral-action-closure}

\subsection{Fixed Momentum Seed and Reduced Spectral Stability}
\label{p3:sec:m_selection}

Table~\ref{tab:p3-cell-momentum-ledger} records the current logical order.
The compact momentum carrier is imported; the Born-even trace below is an
off-shell stability calculation on that already constituted datum.

\begin{table}[H]
\centering
\small
\begingroup
\setlength{\tabcolsep}{3pt}
\renewcommand{\arraystretch}{1.15}
\begin{tabular}{@{}L{3.0cm}L{8.0cm}L{\dimexpr\linewidth-11.0cm-4\tabcolsep\relax}@{}}
\toprule
\textbf{Ledger row}&\textbf{Claim}&\textbf{Discharged by}\\
\midrule
Physical momentum seed
& \(\mathcal H_p\cong\mathcal H_1\) and \(m_p=1\) are already fixed before III-B.
& Paper~II-B, Theorem~\ref{p2:thm:momentum-host-H1},
  Theorem~\ref{p2b:thm:momentum-carrier-rigidity},
  Corollary~\ref{p2:thm:m1-component}\\
Position finite family
& The reciprocal mathematical family is derived; an occurring frontier
history supplies its least stage \(m_x(\widetilde t)\), with primitive value
one.
& Paper~II-C
Corollary~\ref{p2c:cor:derived-mathematical-position-tower} and
\((\mathrm{HIST}_x^{\rm fin})\)\\
Reciprocal response
& The additive Born-even central trace has the unique normalized response \(\tfrac12(\chi+\chi^{-1})\) inside its declared trace class.
& Theorem~\ref{p3:thm:slack-from-spectral-action}\\
Physical basepoint
& The actual momentum spectrum gives \(\chi_{\rm phys}=1\); the off-shell response is strictly convex there.
& Corollary~\ref{p3:cor:selfdual-implies-sharp-saturation}, Lemma~\ref{p3:lem:slack-positive-convex}\\
No momentum-level flow
& A fixed physical \(\mathcal H_1\) cannot change type-support without a newly constituted inter-carrier operation.
& Theorem~\ref{p3:thm:no-momentum-m-flow}\\
Position unfolding
& An occurring Paper-II-C frontier history derives
\(m_x(\widetilde t)\); changing that position statistic does not change the
fixed momentum seed.
& Definition~\ref{p3:def:position-unfolding-level}, Proposition~\ref{p3:prop:position-unfolding-no-insertion}\\
\bottomrule
\end{tabular}
\caption[Fixed-seed reciprocal ledger]{Fixed-seed reciprocal ledger.}
\label{tab:p3-cell-momentum-ledger}
\TableNotes{The current momentum-seed theorem is upstream.  No row in this table uses \(\mathsf P_{\rm sli}\), a level-one ceiling, or GP3 minimization of \(m\).}
\endgroup
\end{table}

\begin{definition}[Realized compact-cell ratio]
\label{p3:def:cell-ratio}
For the physical momentum carrier \(\mathcal H_p\cong\mathcal H_1\), let
\(\Pmax\) denote its actual component spectral maximum and set
\[
\chi \, := \, \frac{\Pmax}{\Punit}
\,=\,
\Lambda_{\rm mom}\,\Rmom,
\qquad
\Punit:=\frac{\hbar}{\Rmom},
\qquad
\Lambda_{\rm mom}:=\frac{\Pmax}{\hbar}.
\]
By Paper~II, Proposition~\ref{p2:prop:component_spectrum}, the Paper-I-A
auxiliary module \(\mathcal H_m\) has
\[
        \Pmax^{(m)}=m\Punit.
\]
Paper~II-B gives \(\Pmax=\Punit\) on the physical seed, hence
\(\chi_{\rm phys}=1\).  For the off-shell reciprocal-response calculation we
extend \(\chi\) to a positive central response coordinate and write
\[
 q:=\log\chi.
\]
Thus \(q=0\) is the physical basepoint; varying \(q\) below does not change
the carrier type-support.
\end{definition}

\begin{definition}[Reduced Born trace coordinates]
\label{p3:def:pre-realization-sector-weights}
For the later additive trace calculation, extend the positive ratio
\(\chi\) to a response coordinate on a fixed reduced Born product leaf and
write
\[
        X(\chi):=\chi,
        \qquad
        Y(\chi):=\chi^{-1}.
\]
They obey
\[
        X(\chi)Y(\chi)=1,
        \qquad
        \mathcal B:X\leftrightarrow Y,
\]
where \(\mathcal B\) denotes Born exchange.  The product constraint is the fixed-leaf condition;
the quotient coordinate is \(q=\frac12\log(X/Y)=\log\chi\).
These are central trace coefficients inside the declared reduced action
domain.  They are not physical cutoff capacities and are not used to select
the finite carrier.
\end{definition}

\begin{definition}[Additive Born central trace class]
\label{p3:def:pre-realization-spectral-functional}
For the off-shell reciprocal-response calculation around the already constituted physical cell, let
\[
        Z_{\rm sec}:=\RR e_x\oplus\RR e_p,
        \qquad
        e_xe_p=0,
        \qquad
        e_x+e_p=1,
\]
be the centre of the unlabelled two-sector direct-sum spectral datum.  The
positive reciprocal weights define the central element
\[
        W(X,Y):=Xe_x+Ye_p,
        \qquad X,Y\ge0.
\]
A reduced cell functional \(F_{\rm cell}(X,Y)\) belongs to the
\emph{additive reciprocal-response trace class} when it is obtained by restricting one
positive finite spectral trace with one common profile to this centre:
\[
        F_{\rm cell}(X,Y)
        =\mathcal T\!\left(W(X,Y)\right),
\]
where \(\mathcal T:Z_{\rm sec}\to\RR\) obeys:
\begin{enumerate}[label=(\roman*),leftmargin=2.1em]
\item \emph{positive trace linearity}: \(\mathcal T\) is a positive real-linear
      functional on \(Z_{\rm sec}\);
\item \emph{common profile}: the same finite spectral profile is used on both
      central summands;
\item \emph{Born symmetry certification}: for the exchange automorphism
      \(\sigma(e_x)=e_p\), \(\sigma(e_p)=e_x\),
\[
        \mathcal T\circ\sigma=\mathcal T.
\]
      On the native same-owner branch where this trace is the scalar output of
      a licensed constructor whose complete input is Born-unmarked, this is
      inherited from Paper~II-A,
      Theorem~\ref{p2a:thm:generated-born-inheritance}, rather than an
      independent physical premise.  It remains an explicit defining
      condition when the class is used abstractly without that origin
      certificate;
\item \emph{nondegeneracy}: \(\mathcal T(e_x)>0\).
\end{enumerate}
The real numbers \(X,Y\) are therefore central scalar coefficients in a finite
trace, not a postulated continuum of copy numbers.  No sector projector beyond
the two already constituted central idempotents, no channel-dependent cutoff,
and no mixed endomorphism is introduced.
\end{definition}

\begin{lemma}[Central two-sector trace representation]
\label{p3:lem:finite-trace-real-weights}
For the finite two-sector datum of
Definition~\ref{p3:def:pre-realization-spectral-functional}, restriction of a
finite block-diagonal spectral trace to \(Z_{\rm sec}\) is automatically a
positive real-linear functional.  Conversely, every positive real-linear
functional on \(Z_{\rm sec}\) has the form
\[
        \mathcal T(Xe_x+Ye_p)=aX+bY,
        \qquad a,b\ge0.
\]
Thus the use of real positive sector weights is exact finite-dimensional trace
linearity and requires no refinement or limiting multiplicity argument.
\end{lemma}

\begin{proof}
Let the common profile operator on the finite direct sum be
\[
        K=K_x\oplus K_p,
        \qquad K_x,K_p\ge0.
\]
For \(W(X,Y)=Xe_x+Ye_p\), ordinary matrix-trace linearity gives
\[
\Tr(WK)
 =X\Tr(K_x)+Y\Tr(K_p).
\]
Hence the restricted spectral trace is positive and real-linear.  Conversely,
\(Z_{\rm sec}\) is the two-dimensional real vector space with basis
\((e_x,e_p)\), so a real-linear functional is determined by
\[
        a:=\mathcal T(e_x),
        \qquad b:=\mathcal T(e_p),
\]
and positivity on the positive cone forces \(a,b\ge0\).  Therefore
\(\mathcal T(Xe_x+Ye_p)=aX+bY\) for all \(X,Y\ge0\).
\end{proof}


\begin{lemma}[Trace linearity and Born exchange force the reciprocal valuation]
\label{p3:lem:pre-realization-linear}
Let \(F_{\rm cell}\) be a additive reciprocal-response functional.  Then
there is a constant \(C_{\rm tr}>0\) such that
\[
        F_{\rm cell}(X,Y)=C_{\rm tr}(X+Y).
\]
Consequently, after normalizing the self-dual cell to unit relative action,
\[
        F_{\rm rel}(X,Y)
        :=\frac{F_{\rm cell}(X,Y)}{F_{\rm cell}(1,1)}
        =\frac{X+Y}{2}.
\]
\end{lemma}

\begin{proof}
Lemma~\ref{p3:lem:finite-trace-real-weights} gives
\[
        F_{\rm cell}(X,Y)=aX+bY
\]
with \(a,b\ge0\).  Born invariance under \((X,Y)\mapsto(Y,X)\) implies
\[
        aX+bY=aY+bX
\]
for all \(X,Y\ge0\), hence \(a=b\).  Nondegeneracy gives the common value
\(C_{\rm tr}:=a=b>0\).  Division by
\(F_{\rm cell}(1,1)=2C_{\rm tr}\) gives the normalized
functional.
\end{proof}


\begin{definition}[Normalized additive Born functional]
\label{p3:def:minimal-additive-born-functional}
A reciprocal-cell functional belongs to \(\mathcal S_{\rm Born}^{\rm add}\) when it is the restriction of one positive finite spectral trace to the positive cone of the central two-sector algebra, is invariant under Born exchange \(X\leftrightarrow Y\), uses a single common profile, and is normalized on the self-dual product leaf.  Positive real linearity is exact matrix-trace linearity on the central coefficients, as proved in Lemma~\ref{p3:lem:finite-trace-real-weights}.  This is the constitution class used in the reciprocal slack argument; it is not the class of all conceivable Born-even actions.
\end{definition}
\begin{remark}[Additive trace class versus arbitrary Born-even functions]
The uniqueness statement below is internal to the adopted additive
finite-trace subclass.  Nonlinear Born-even expressions such as
\((X+Y)^2\), \(XY\), or higher functions of \(X/Y\) are not ruled out by
GP1--GP3; they define different action data and lie outside the class used
here.
\end{remark}


\begin{theorem}[Reciprocal slack in the additive central trace class]
\label{p3:thm:slack-from-spectral-action}
On a fixed reduced Born product leaf, the normalized \(\chi\)-dependent part of the no-insertion
off-shell reciprocal spectral action is unique inside the additive finite-trace
Born-symmetric class of the unlabelled reciprocal spectral datum:
\[
        F_{\rm rel}(\chi)
        =\frac12\left(\chi+\chi^{-1}\right).
\]
Its excess over the self-dual geometric mean is exactly
\[
\boxed{
\mathfrak S_{\rm cell}(\chi)
:=F_{\rm rel}(\chi)-1
=\frac12\left(\chi+\chi^{-1}\right)-1
=\cosh(\log\chi)-1.
}
\]
Thus the reciprocal slack is the normalized surplus of the adopted additive
Born-symmetric trace in
Definition~\ref{p3:def:pre-realization-spectral-functional}, up to the common
normalization already quotiented in the relative functional.  Other
Born-even actions are not thereby excluded.
\end{theorem}

\begin{proof}
By Definition~\ref{p3:def:pre-realization-sector-weights}, the fixed-leaf
sector weights are
\[
  (X,Y)=(\chi,\chi^{-1}).
\]
Lemma~\ref{p3:lem:pre-realization-linear} gives
\[
        F_{\rm rel}(X,Y)=\frac{X+Y}{2},
\]
hence
\[
        F_{\rm rel}(\chi)=\frac12(\chi+\chi^{-1}).
\]
The geometric mean of the two reciprocal sector weights is fixed on the leaf:
\[
        \sqrt{XY}=\sqrt{\chi\chi^{-1}}=1.
\]
Therefore the relative arithmetic-volume surplus above the geometric mean is
\[
        \frac{X+Y}{2\sqrt{XY}}-1
        =\frac12(\chi+\chi^{-1})-1.
\]
Since \(q=\log\chi\), this is also \(\cosh q-1\).  The uniqueness follows from
the linearity lemma: any other continuous Born-even additive uninserted
functional differs only by an overall positive constant, which disappears in
the relative normalization and cannot change the minimizer.
\end{proof}


\begin{lemma}[Born-even positivity and strict convexity of reciprocal slack]
\label{p3:lem:slack-positive-convex}
For every \(\chi>0\),
\[
\mathfrak S_{\rm cell}(\chi)\ge 0,
\]
with equality if and only if \(\chi=1\).  In the logarithmic coordinate \(q=\log\chi\),
\(\mathfrak S_{\rm cell}(q)=\cosh q-1\) is strictly convex and has a unique critical point,
namely \(q=0\).
\end{lemma}

\begin{proof}
The arithmetic--geometric mean inequality gives
\[
\frac12\left(\chi+\chi^{-1}\right)
\ge
\sqrt{\chi\chi^{-1}}=1,
\]
with equality if and only if \(\chi=\chi^{-1}\), hence \(\chi=1\) since \(\chi>0\).
Equivalently, with \(q=\log\chi\),
\[
\mathfrak S_{\rm cell}(q)=\cosh q-1.
\]
Then
\[
\frac{d}{dq}\mathfrak S_{\rm cell}=\sinh q,
\qquad
\frac{d^2}{dq^2}\mathfrak S_{\rm cell}=\cosh q>0.
\]
Thus the only critical point is \(q=0\), and it is the unique global minimum.
\end{proof}

\begin{theorem}[Imported physical momentum seed and derived level]
\label{p3:thm:factor-local-momentum-shape}
\label{p3:thm:factor-local-momentum-m1}
\label{p3:prop:constitutive-minimization-Dcl}
\label{p3:thm:truncation_m1}
\label{p3:thm:minimal-momentum-carrier}
\label{p3:thm:paper-i-series-selected-cell-closure}
Paper~II-B proves, relative to the frozen momentum-seed signature,
\[
 \boxed{\mathcal H_p\cong\mathcal H_1,\qquad m_p=1.}
\]
The carrier type-support is fixed by GP2 exact/host-faithful constructor
closure, and GP3 removes only unmarked realization multiplicity.  The level
\(m_p=1\) is read off from the already fixed carrier; it is not a GP3
minimization variable.
\end{theorem}

\begin{corollary}[Physical reciprocal basepoint]
\label{p3:thm:unique-cell-selfdual}
\label{p3:cor:selfdual-implies-sharp-saturation}
On the imported physical momentum seed the actual component maximum is
\(\Pmax=\Punit\).  Hence
\[
 \boxed{\chi_{\rm phys}=1,\qquad q_{\rm phys}=0.}
\]
This is the basepoint about which the additive Born response is evaluated; the
strict convexity of Lemma~\ref{p3:lem:slack-positive-convex} is a stability
statement, not a seed selector.
\end{corollary}

\begin{remark}[Role of the flat hyper-K\"ahler cone]
\label{p3:rem:flat-cone-m1}
The flat hyper-K\"ahler cone of Paper~I-A is best read here as a rigidity and vacuum-consistency
statement for the constituted round parent and the derived momentum seed. It strengthens the interpretation of the
physical compact momentum realization, but it is not itself an independent functional of the
truncation level.
\end{remark}


\begin{remark}[No hidden cone-curvature source in the constituted-cell action]
\label{p3:rem:flat-cone-no-dynamical-source}
Since the Paper-I-A Einstein-normalized cone is flat, the reduced action below does
not acquire a scalaron potential, vacuum term, mass defect, or hierarchy
correction from cone curvature.  All dynamical scale data in Paper~III must come
from the no-insertion spectral functional, Hopf/contact residue, local heat-germ
completion, or explicitly declared finite carrier data.
\end{remark}


\begin{remark}[Fixed momentum seed versus position unfolding]
\label{p3:rem:mp-mx-distinction}
The momentum entry is already fixed upstream:
\[
 m_p=1,
 \qquad
 \mathcal H_p\cong\mathcal H_1.
\]
A position-side finite label is derived only after a frontier occurs.
Definition~\ref{p3:def:position-unfolding-level} gives its stagewise value
\(m_x(\widetilde t)\); cumulative frontier inclusion proves monotonicity
without a separate MON premise.  None of these statements reopens the
momentum-seed constitution.
\end{remark}

\begin{theorem}[Fixed-carrier invariance firewall]
\label{p3:thm:no-momentum-m-flow}
Inside the loop-derived class \(\Dcl\), the physical momentum datum is the fixed
structured realization with underlying carrier \(\mathcal H_p\cong\mathcal H_1\).
Automorphisms, and any unitary endomorphism dynamics separately specified on
that fixed carrier, preserve its dimension and irreducible type-support.  The
class \(\Dcl\) is a scalar action-form class, not by itself an equation of
motion, Hamiltonian, or evolution law.  A transition to an auxiliary module
\(\mathcal H_m\) with \(m\ne1\) requires a separately constituted
inter-carrier embedding, quotient, projection, or transition operation and
therefore lies outside \(\Dcl\) unless that datum is explicitly added.
\end{theorem}

\begin{proof}
The automorphism group and every unitary one-parameter endomorphism family on
the physical momentum realization act within its fixed finite presentation
algebra.  Such maps preserve carrier dimension and \(\operatorname{Sp}in(4)\)-type content.
Generic \(\mathcal H_m\) are auxiliary
representation-theoretic modules; passing from \(\mathcal H_1\) to one of
them is not an inner evolution of the fixed spectral datum.  Such a passage
therefore requires new typed structure.  Since Definition~\ref{p3:def:D1}
only classifies scalar spectral-action readouts, any physical evolution also
requires a separately specified equation of motion or Hamiltonian.  Neither
that dynamics nor an inter-carrier transition is generated by the action
class itself.
\end{proof}

\begin{proposition}[Paper-I-A dependency-split handoff on the derived seed]
\label{p3:prop:postselection-finite-cell-imports}
On the derived physical momentum carrier \(\mathcal H_p\cong\mathcal H_1\),
Paper-I-A representation-theoretic statements apply directly.  If the coordinate completion is specialized to \(\mathsf C_{\rm PI}\),
its absolute projected-coordinate and completion-dependent scale formulas
apply on the same carrier.  On the weaker \((\mathrm{ND}_{\rm sp})\)
stratum, the normalized protected even Clifford algebra already applies by
Paper~II-C.  If the core-support role \(\mathsf{Supp}_{2T}\) also occurs, the corresponding
Hurwitz/trimer support theorems apply.  Neither specialization nor role
occurrence is selected by the momentum-seed theorem.
\end{proposition}

\begin{proof}
The carrier statement is Theorem~\ref{p2:thm:momentum-host-H1} together with
Theorem~\ref{p2b:thm:momentum-carrier-rigidity}.  The additional coordinate
normalization and finite-support conclusions have exactly the independent
antecedents recorded in Papers~I-A--II; neither follows from \(m_p=1\) alone.
\end{proof}

\begin{remark}[No circular use of the finite seed]
\label{p3:rem:no-circular-use-finite-cell}
III-B uses \(\mathcal H_1\) only after it has been constituted upstream.  The
action theorems below therefore classify dynamics on a fixed seed and do not
feed back into the Paper-II carrier proof.
\end{remark}

\begin{remark}[Auxiliary transverse counting comparisons]
The transverse harmonic-cell, saturated-area, and diagonal \(m^4\)
calculations are presented as comparison tests in Paper~III-G,
Section~\ref{p3b:app:thc}.  They do not select the physical momentum seed or define
\(m_x(\widetilde t)\).  For the reciprocal-position frontier the active result
is now Paper~II-C,
Theorem~\ref{p2c:thm:born-transported-transverse-shell-reconstruction}, which
derives
\[
 E_{m_x}^x\simeq\mathcal T_{x,m_x+1}^{\perp},
 \qquad N_{\partial x}(m_x)=(m_x+1)^2.
\]
Thus the position-side tag \((\mathrm{ID}_{\rm THC})\) is not an
additional hypothesis: the result is a theorem, with the exact finite-level shift \(\mu_x=m_x+1\).  No analogous
identification of the full physical momentum carrier \(\mathcal H_1\) with an
auxiliary scalar transverse window is generated.
\end{remark}

\begin{definition}[Position-unfolding bookkeeping level]
\label{p3:def:position-unfolding-level}
Assume the Paper-II-C frontier-history occurrence
\((\mathrm{HIST}_x^{\rm fin})\), with actual finite frontiers
\(F_x(\widetilde t)\).  The symbol
\[
 m_x(\widetilde t)
 :=m_x^{\rm res}(F_x(\widetilde t))
\]
is the derived least scalar resolution, not an independently supplied integer
trajectory.  This paper does not identify it with a floor of a geometric
radius divided by a new Born length.  If the occurring history is cumulatively
retained, frontier inclusion proves monotonicity without a separate
\((\mathrm{MON}_x)\) premise.  The primitive Paper-II-C frontier already has
resolution one; an epoch-identification premise is needed only if a chosen
physical time is asserted to coincide with that primitive stage.
\end{definition}

\begin{proposition}[Position unfolding does not introduce momentum-sector insertions]
\label{p3:prop:position-unfolding-no-insertion}
A supplied position frontier history, and hence its derived statistic
\(m_x(\widetilde t)\), may change while the
physical momentum seed remains \(m_p=1\).  This is not momentum-level flow and
does not contradict Theorem~\ref{p3:thm:no-momentum-m-flow}.
\end{proposition}

\begin{proof}
The two labels have different typed roles.  The momentum value describes the
already fixed Paper-II seed, whereas \(m_x(\widetilde t)\) summarizes the
separate occurring position frontier.  No inter-carrier momentum map is
created by changing the latter.
\end{proof}

\begin{theorem}[Calibrated local cotangent phase-volume count]
\label{p3:thm:born-liouville-mp-mx-seesaw}
Let \(C\subset T^*U\) be a fixed local cotangent cell with canonical
symplectic form
\[
 \omega_{\rm can}=\sum_{\mu=1}^{4}dp_\mu\wedge dq^\mu
\]
and impose the physical phase-unit calibration \((\mathrm{CAL}_{\hbar})\).
Then the native local semiclassical count is the joint phase-volume count
\[
 \boxed{
 N_{\rm phase}(C)
 =\frac{1}{(2\pi\hbar)^4}
  \int_C\frac{\omega_{\rm can}^{4}}{4!}.}
\]
The coordinate position and momentum polarizations are Lagrangian:
\[
 \omega_{\rm can}|_{p={\rm const}}=0,
 \qquad
 \omega_{\rm can}|_{q={\rm const}}=0.
\]
Thus \((\mathrm{CAL}_{\hbar})\) calibrates the joint phase unit; it does not
generate two separately symplectic four-dimensional position/momentum factors
or separate dimensionless counts \(N_x^{\rm loc}\) and \(N_p^{\rm loc}\).
Under the eight-dimensional Born reversal the canonical two-form is
anti-symplectic while the Liouville volume is invariant,
\[
 R_B^*\omega_{\rm can}=-\omega_{\rm can},
 \qquad
 R_B^*(\omega_{\rm can}^{4})=\omega_{\rm can}^{4}.
\]
No identification with \(\dim\mathcal H_1\), \(N_p(1)\), the word label
\(m_p\), or a macroscopic position trajectory \(m_x(\widetilde t)\) is
asserted.
\end{theorem}

\begin{proof}
The canonical cotangent form is nondegenerate on \(T^*U\), and its fourth
exterior power is the Liouville volume.  Pullback to either coordinate
polarization kills every term because one member of each
\(dp_\mu\wedge dq^\mu\) pair vanishes.  Born reversal changes the sign of the
two-form, hence preserves its fourth power.  Division by the calibrated cell
volume \((2\pi\hbar)^4\) gives the displayed joint count.  The optional
product-symplectic factorization, when a genuinely different product model is
supplied, is recorded separately in Paper~III-G, Section~\ref{p3:app:seesaw}.
\end{proof}

Table~\ref{tab:p3-completed-carrier-ledger} records the current reciprocal geometry and finite interfaces.

\begin{table}[H]
\centering
\small
\begingroup
\setlength{\tabcolsep}{3pt}
\renewcommand{\arraystretch}{1.15}
\begin{tabular}{@{}L{3.0cm}L{9.5cm}L{\dimexpr\linewidth-12.5cm-4\tabcolsep\relax}@{}}
\toprule
\textbf{Ledger row}&\textbf{Claim}&\textbf{Discharged by}\\
\midrule
Reciprocal smooth roles
& Distinct round momentum and Born-derived reciprocal-position roles form the Born-reciprocal product; dimensionful \(\Rcs\) uses \((\mathrm{CAL}_{xp})\).
& Paper~II-C/GP2 and Paper~III-A, with \(\mathrm{CAL}_{xp}\) only for the dimensionful radius representative; \(\mathsf P_{\rm mc}\) is optional full-profile realization\\
Momentum finite seed
& The physical momentum carrier is \(\mathcal H_1\) with \(m_p=1\).
& Paper~II-B\\
Position finite interface
& The mathematical family is derived; an actual frontier history occurs and
its least current resolution is derived from that frontier, with primitive
value one.
& Paper~II-C derived tower and \((\mathrm{HIST}_x^{\rm fin})\)\\
Scalar phase pair
& The reciprocal scalar torus carries the primitive Poincar\'e class; its dimensionful representative uses \((\mathrm{CAL}_{\hbar})\).
& Paper~III-A\\
Global symplectic status
& The full compact smooth product is not a global symplectic manifold.
& Theorem~\ref{p3:thm:completed-carrier-not-global-symplectic}\\
Born ratio decomposition
& Diagonal fibres and the relative quotient form the Born-even symmetric-pair geometry.
& Theorem~\ref{p3:thm:born-parity-decomposition}\\
\bottomrule
\end{tabular}
\caption[Completed reciprocal-geometry ledger]{Completed reciprocal-geometry ledger.}
\label{tab:p3-completed-carrier-ledger}
\TableNotes{Forward gauge/gravity, dilation, hierarchy, and index-locking statements are not native reciprocal-geometry/action results of III-B and are therefore not included in this ledger.}
\endgroup
\end{table}

\subsection{The Completed Compact Reciprocal Geometry}

\begin{theorem}[Completed compact Born/contact geometry]
\label{p3:thm:phasespace}
By Paper~II-C intrinsic evaluation and the Paper-II-A GP2 reciprocal-role
closure, the reciprocal-role objectification, relation-free composition, and
Born exchange involution are already derived; no \(\mathsf P_{\rm mc}\) or separate
\(\mathrm{BR}_{\rm ind}\) or \(\mathrm{IND}_{xp}\) premise is required.
The abstract reciprocal smooth role is
\[
 (S^3_B\times S^3_p(\Rmom))\times(S^1_t\times S^1_E).
\]
Under \((\mathrm{CAL}_{xp})\) it is written
\[
 \boxed{\mathcal M_8\cong
 (S^3_x(\Rcs)\times S^3_p(\Rmom))\times(S^1_t\times S^1_E),
 \qquad \Rcs=\alpha_{xp}\Rmom.}
\]
Under \((\mathrm{CAL}_{\hbar})\) the already derived scalar Poincar\'e class
has the dimensionful representative with cell \(2\pi\hbar\).  These
calibrations do not introduce a scalar-circle metric.  The physical momentum
carrier attached to this geometry is \(\mathcal H_p\cong\mathcal H_1\); if a
position tower is constituted, the generic finite reciprocal family is
\(\mathsf K(m_x,1)\), with \(\mathsf K(1,1)\) only under the primitive-epoch
identification \((\mathrm{EPOCH}_{x,0}^{\rm prim})\).
\end{theorem}

\begin{proof}
The reciprocal smooth product and Born exchange are premise-minimal
Paper-II-C/GP2 consequences, imported and completed in Paper~III-A.  The
optional \(\mathsf P_{\rm mc}\) clause realizes the complete transported
profile family as G0 tests but is not needed for the product.  The two
dimensionful calibrations remain separate Paper~III-A boundaries and are not
consequences of that abstract Born closure.  The physical momentum seed is imported
from Paper~II-B.  The position finite statement is the separately typed Paper-II
position-resolution interface.  No lower-ideal or common finite-level selector is used.
\end{proof}

\subsubsection{The completed reciprocal host}

\begin{remark}[Host presentation and construction chronology]
\label{p3b:rem:m8-host-construction-chronology}
The compact reciprocal host has smooth support
\[
 M_8=(S_x^3\times S_p^3)\times(S_t^1\times S_E^1).
\]
The notation distinguishes the support \(S_p^3\) from the derived,
enriched momentum package \(S_1^3\) constituted on it.  That package
is not an extra product factor and is not a sphere contained in
another sphere.  The host roles are:
\begin{center}
\begin{tabular}{c|p{0.65\textwidth}}
 \(S_p^3\) & Smooth momentum support carrying the derived \(S_1^3\)
 package, Hopf/Spin geometry and finite geometric descendants.\\
 \(S_x^3\) & Evaluation-dual position role, with its own geometric
 configuration and induced reciprocal action.\\
 \(S_t^1\) & Transported scalar clock primitive and its character action.\\
 \(S_E^1\) & The dual scalar primitive, paired with the clock on the
 actual scalar torus.
\end{tabular}
\end{center}
The position and momentum configurations form the full Cartesian
product.  Both scalar circles are retained; a Hopf circle in a sphere
does not substitute for either one.  The scalar phase pairing does
not supply an independent native metric action on each circle.
The hosting relation of
Definition~\ref{p2a:def:host-constructor-incidence-contract} allows
arbitrary licensed depth inside these roles, as well as constructors
with inputs from several factors.  It does not make every object a
direct constructor child of the whole product.
\end{remark}

The shared mathematics of this labelled reciprocal product is collected in
Paper~I-C \cite{Kugelmann2026Paper1B}.  The references below use its exact
results after the host and common-model identifications constructed here;
in particular the relative quotient and its tangent split are
Theorems~\ref{p1b:thm:relative} and~\ref{p1b:thm:symmetric}.

\begin{theorem}[Derived subgroup, connected center, and abelianization of the completed parent]
\label{p3b:thm:v193-m8-derived-center-abelianization}
For the labelled completed compact parent
\[
 M_8=SU(2)_x\times SU(2)_p\times U(1)_t\times U(1)_E,
\]
one has
\[
 [M_8,M_8]=SU(2)_x\times SU(2)_p,
\]
\[
 Z(M_8)=\{\pm1\}_x\times\{\pm1\}_p
       \times U(1)_t\times U(1)_E,
 \qquad
 Z(M_8)^0=U(1)_t\times U(1)_E,
\]
and
\[
 M_8^{\rm ab}:=M_8/[M_8,M_8]
 \cong U(1)_t\times U(1)_E.
\]
The labelled scalar projection realizes the displayed abelianization.
Consequently every continuous one-dimensional character of \(M_8\) factors
uniquely through the scalar torus, and
\[
 X^*(M_8):=\operatorname{Hom}_{\rm cont}(M_8,U(1))\cong\mathbb Z^2.
\]
Moreover
\[
 \pi_1(M_8)\cong\mathbb Z^2,
 \qquad
 H^1(M_8;\mathbb Z_2)\cong\mathbb Z_2^2,
\]
so all nontrivial loop classes and all spin-structure difference classes are
scalar-supported.  The character lattice is the integral dual of the loop
lattice, while spin structures form a torsor over its mod-two dual; these are
related but are not identified as one object.
\end{theorem}

\begin{proof}
Apply Paper~I-C, Theorem~\ref{p1b:thm:group}, to the displayed labelled
group.  The loop and spin-torsor statements are respectively
Theorems~\ref{p1b:thm:cohomology} and~\ref{p1b:thm:spin}; their scalar
labels are the same \((t,E)\) labels used here.
\end{proof}

This is the enriched microscopic Born/contact/spectral host of
Theorem~\ref{p3:thm:phasespace}.  The compact completed carrier is not
thereby globally symplectic
(Theorem~\ref{p3:thm:completed-carrier-not-global-symplectic}), and
no physical eight-dimensional spacetime or new scalar metric is
assumed.  Both diagonal-fibre spin extensions remain constituents
of the joint geometry.  Positive determinant densities and response
owners retain positive real actions; they are not the complex
relative determinant line of the Hopf--Hurwitz comparison.


\begin{theorem}[The completed compact geometry is not a global symplectic manifold]
\label{p3:thm:completed-carrier-not-global-symplectic}
Let
\[
\mathcal C_8:=(S^3_x\times S^3_p)\times(\Sone_t\times\Sone_E)
\]
denote the smooth compact manifold underlying Theorem~\ref{p3:thm:phasespace}.  Then
\(\mathcal C_8\) admits no global symplectic form.  Consequently, throughout Paper~III and in
all later-layer uses, Theorem~\ref{p3:thm:phasespace} is a statement about the compact
Born/contact/spectral geometry together with local Darboux/Born charts and scalar universal-cover
Hamiltonian evolution; it is not a claim that the compact product itself is a global Hamiltonian phase space.
\end{theorem}

\begin{proof}
The underlying manifold is the \(M_8\) of Paper~I-D,
Proposition~\ref{p1b:prop:lcs}, which proves the obstruction.  The local
Darboux/Born forms used here are cotangent-model forms and hence do not
assert a global symplectic form on this compact product.
\end{proof}

\begin{corollary}[Global/local phase-geometry hierarchy]
\label{p3b:cor:global-local-phase-geometry-hierarchy}
The reciprocal construction has three distinct geometric levels:
\[
\boxed{\begin{gathered}
M_8\ \text{is compact and not globally symplectic};\\[0.25em]
(M_8,\Omega,H=-\mathrm d\Omega)\ \text{carries}\\
\text{the exact twisted-Dirac graph }L_\Omega;\\[0.25em]
T^*U\ \text{is locally canonical symplectic}.
\end{gathered}}
\]
Here the normalized global reciprocal two-form \(\Omega\) is nondegenerate
but nonclosed, while the Paper-III-A cotangent germ carries the closed
canonical form \(\omega_{\rm can}\).  Consequently the local Hamiltonian
phase geometry is not a declaration that the compact eight-manifold itself is
an ordinary symplectic phase space, and \(\omega_{\rm can}\) is not the
global nonclosed form \(\Omega\).
\end{corollary}
\begin{proof}
The first statement is the preceding theorem.  The reciprocal transgression
calculation and exact twisted-Dirac theorem below give
\(\mathrm d\Omega\neq0\), \(H=-\mathrm d\Omega\), and integrability of the graph
\(L_\Omega\) for the \(H\)-twisted Courant bracket.  Paper~III-A independently
generates \(T^*U\) with its closed canonical symplectic form.  The objects and
forms therefore have different domains and no contradiction arises.
\end{proof}


\begin{mdframed}[nobreak=true]
\noindent\textbf{Terminology convention.}
After Theorem~\ref{p3:thm:completed-carrier-not-global-symplectic}, the phrase ``compact
microscopic geometry'' means the smooth Born/contact/spectral product
\[
(S^3_x\times S^3_p)\times(S^1_t\times S^1_E).
\]
Hamiltonian evolution is always local in Darboux/Born charts, or on the scalar universal cover when
time and energy accumulation are required.  This smooth role space is not the finite ontic Hilbert
carrier.
\end{mdframed}


\begin{remark}[Export-layer hierarchy interface]
\label{p3:rem:downstream-hierarchy-export}
The hierarchy invariant exported from Paper~III is therefore not a global symplectic volume or
capacity of \(\mathcal C_8\).  The exported topological logarithmic step is the primitive Hopf
contact spectral residue of Theorem~\ref{p3:thm:Vlog},
\[
L_{\mathrm{Hopf}}^{(1)}
=\int_{S^3}\eta_0\wedge d\eta_0
=(2\pi)^2,
\]
with local smooth completion fixed only after the regularization prescription
\(({\rm REG})\) stated in Paper~III-C.  Later papers may cite that conditional theorem to
exclude any hidden global-symplectic foundation in the hierarchy construction.
\end{remark}

\begin{theorem}[G0-admissible non-destructive relative readout]
\label{p3b:thm:relative-readout-admission}
Let \(\mathcal M_8=K_x\times K_p\),
\(K=SU(2)\times U(1)\), and let
\(\iota_x:K_x\xrightarrow{\sim}K\) and
\(\iota_p:K_p\xrightarrow{\sim}K\) be the generated common-model
isomorphisms of Paper~III-A.  Before evaluation, declare a relative target
sort \(\mathsf{Eval}_{\rm rel}(K)\) and a kernel-pair relation sort on
\(\mathcal M_8\times\mathcal M_8\).  Put
\[
 D_{\rm rel}:=(\mathcal M_8,\iota_x,\iota_p,={}_K)
\]
and freeze the finite relative-readout word
\[
 w_r:\quad
 (a,b)\longmapsto
 r(a,b):=\iota_x(a)\iota_p(b)^{-1},
 \qquad
 \widehat r(a,b):=((a,b),r(a,b)),
\]
and let \(\Gamma_r:=\operatorname{im}\widehat r\) be the operation graph.
Define the typed kernel-pair relation
\[
 \mathcal R_r
 :=\bigl\{(c,c')\in\mathcal M_8\times\mathcal M_8:
 r(c)=r(c')\bigr\}
 \subseteq\mathcal M_8\times\mathcal M_8,
\]
using equality on \(K\) as separated by its already generated tests.  This is
a relation-valued output, not a partially specified numerical diagnostic.
Then
the complete constructor
\[
 (r,\Gamma_r,\operatorname{im}r,\mathcal R_r)
\]
is the required output of \(w_r\).  Product projection, the retained
\(\iota\)-maps, multiplication, inversion, graph/image formation, and the
generated equality comparison are all already licensed in the frozen
\(\Sigma_{\rm rel}\) grammar at this stage.  Graph, image, and kernel-pair
relation formation are declared constructor operations in that same grammar,
with their intermediates retained and their output sorts predeclared.  Hence
\[
 \operatorname{Orig}_{\Sigma_{\rm rel}}
 (w_r;D_{\rm rel},\mathsf{Eval}_{\rm rel}(K))
\]
holds.  Paper~II-A,
Theorem~\ref{p2a:thm:rooted-generated-image-corealization}, then implies that
on the nonobstructed branch its
relative readout doctrine is co-realized with the same complete reciprocal
root tuple and requires no second G0 occurrence.  The full relation-free
product and all factor-separating tests remain present.

This is the quotient/readout case of
Theorem~\ref{p2a:thm:universal-corealization-bridge-normal-form}.  The
co-realized source and readout are connected by
\((r,\Gamma_r,\mathcal R_r)\), but \(r\) is not injective: its fibres are the
diagonal \(K\)-orbits displayed below.  Therefore this bridge is not a Born
exchange between \(\mathcal M_8\) and \(K\), and no inverse process is
generated.

Moreover, the kernel-pair classes of \(r\) are exactly the orbits of the free
diagonal right action
\[
 (a,b)\cdot k
 :=\bigl(\iota_x^{-1}(\iota_x(a)k),
         \iota_p^{-1}(\iota_p(b)k)\bigr).
\]
Thus \(\mathcal M_8/{\sim_r}\cong K\) is a generated homogeneous
readout-role presentation.  For the nonabelian \(SU(2)\) factor,
\(\Delta K\) is not normal in \(K\times K\); therefore \(\sim_r\) is not a
congruence for componentwise source-group multiplication.  Any group law on
the readout is transported through the induced bijection with \(K\), not
inherited as a quotient-group law on the source.
\end{theorem}

\begin{proof}
The frozen word \(w_r\) uses only the product projections, the frozen maps
\(\iota_x,\iota_p\), multiplication, inversion, and the already generated
equality comparison; its graph, image, and kernel-pair outputs are declared
before evaluation and retained for reuse.  It is candidate-independent,
host-natural, type
preserving, and contains no representative choice, scale, profile, or new
parameter.  The graph lift \(\widehat r\) is injective because its first
coordinate is the identity on \(\mathcal M_8\).  Hence it reflects every
source equality, inequality, factor distinction, and frozen intrinsic defect.
In particular, no factor test is identified or deleted.

With \(e\in K\) the identity, the canonical section
\[
 s(k):=(\iota_x^{-1}(k),\iota_p^{-1}(e))
\]
satisfies \(r\circ s=\operatorname{id}_K\).  Thus the proposed output is
exactly the operation graph, its complete essential image \(K\), and the
typed kernel-pair relation \(\mathcal R_r\); it inserts no support, connection, metric,
calibration, state, or history.  Compactness of the source and continuity of
the finite constructor word give GP1 admissibility.

For the enlarged presentation, the insertion defect is zero because every
operation is generated; the omission defect is zero because the source,
factor tests, graph, and full image are retained; the relation defect is zero
because equality of readout values is not imposed as equality of source
points; the typing defect is zero because source, graph, kernel pair, and
readout sorts remain distinct; and the closure defect is zero because the
complete graph, image, and kernel-pair relation are present.  Hence the joint
five-part GP2 realization defect vanishes.  This exhibited realization rules
out \(\operatorname{Obs}_{\rm FRC}(w_r,D_{\rm rel})\), so the rooted
generated-image theorem gives co-realization.  The naturality and zero-defect
checks verify the realization rather than licensing \(w_r\).  The conclusion
is subject to the role/token firewall: it does not
actualize every \(k\in K\) as an independent token.

Diagonal invariance follows from
\[
 (\iota_x(a)k)(\iota_p(b)k)^{-1}
 =\iota_x(a)\iota_p(b)^{-1}.
\]
Conversely, if \(r(a,b)=r(a',b')\), set
\(k:=\iota_p(b)^{-1}\iota_p(b')\).  Then the second coordinates agree after
the action, and the ratio equality gives agreement of the first coordinates.
Thus each kernel-pair class is exactly one diagonal orbit; freeness follows
from the second factor.  The induced map from the compact orbit space to the
Hausdorff group \(K\) is a continuous bijection, hence a homeomorphism, and
the surjective-submersion/free-proper-action theorem makes it a diffeomorphism.
The generated comparison makes \(\sim_r\) a congruence in the declared
readout doctrine, so Paper~II-A's generated-orbit inheritance theorem applies
in that role.  Finally, \(\Delta K\) is normal in \(K\times K\) only when the
relevant factor is abelian; conjugating \((k,k)\) by \((x,y)\) gives
\((xkx^{-1},yky^{-1})\), which is generally not diagonal for \(SU(2)\).
This proves the source-group-law firewall.
\end{proof}

\begin{remark}[Right fibre action versus left conjugation action]
\label{p3b:rem:v193-diagonal-action-typing}
For \(r(a,b)=ab^{-1}\) the free right diagonal action is the principal fibre
action,
\[
 r(ak,bk)=r(a,b),
\]
whereas the left diagonal action descends to conjugation,
\[
 r(ka,kb)=k\,r(a,b)\,k^{-1}.
\]
The first is fibre redundancy; the second is a genuine action on the relative
readout.  They must not be identified.
\end{remark}

These two action laws are the distinct identities in Paper~I-C,
Theorem~\ref{p1b:thm:relative}.

\begin{remark}[Supplementary pair-groupoid analysis]
The relative readout admits a typed pair-groupoid and cocycle presentation,
proved in Paper~III-G, Section~\ref{p3b:app:pair-groupoid-analysis}.  Its convolution
completion is \(\mathcal K(L^2(K))\), and the tangent deformation has zero
fibre \(C_0(T^*K)\).  These are analytic consequences of the admitted
readout, not a Hamiltonian, state, trajectory, physical occurrence, or
determinant ancestry.
\end{remark}

\begin{theorem}[Born-equivariant ratio fibration and symmetric-pair decomposition]
\label{p3:thm:born-parity-decomposition}
Write the completed compact role space in dimensionless group coordinates as
\[
        \mathcal M_8\cong K_x\times K_p,
        \qquad
        K:=SU(2)\times U(1),
\]
where \(K_x=S^3_x\times S^1_t\) and
\(K_p=S^3_p\times S^1_E\) are labelled reciprocal copies.  Let
\[
 \iota_x:K_x\xrightarrow{\sim}K,
 \qquad
 \iota_p:K_p\xrightarrow{\sim}K
\]
be the canonical common-model identifications inherited from the reciprocal
construction of Paper~III-A.  Define the typed ratio map
\[
        \Pi_{\rm rel}:K_x\times K_p\longrightarrow K,
        \qquad
        \Pi_{\rm rel}(a,b):=\iota_x(a)\,\iota_p(b)^{-1}.
\]
This is the non-destructive readout admitted in
Theorem~\ref{p3b:thm:relative-readout-admission}.  Every orbit-space statement
below presents that output role while the source product remains retained.
Then:
\begin{enumerate}[label=(\roman*),leftmargin=2.2em]
\item \(\Pi_{\rm rel}\) is a surjective submersion and is the quotient by the
      free diagonal right \(K\)-action
\[
 (a,b)\cdot k
 :=\Bigl(\iota_x^{-1}(\iota_x(a)k),
          \iota_p^{-1}(\iota_p(b)k)\Bigr).
\]
Its fibres are copies of the diagonal fibre
\[
        F_4\cong K_{\rm diag}
        \cong S^3_{\rm diag}\times S^1_{\rm diag},
\]
while the quotient is the relative role space
\[
        B_4\cong K_{\rm rel}
        \cong S^3_{\rm rel}\times S^1_{\rm rel}.
\]
\item After transporting the two identity tangent spaces through
\(d\iota_x\) and \(d\iota_p\), the typed Born-swap differential on
\(\mathfrak k\oplus\mathfrak k\) has the eigenspace decomposition
\[
\mathfrak k\oplus\mathfrak k
 =\mathfrak k_+\oplus\mathfrak k_-,
\qquad
\mathfrak k_+=\{(A,A)\},
\qquad
\mathfrak k_-=\{(A,-A)\}.
\]
It is a symmetric-pair decomposition:
\[
[\mathfrak k_+,\mathfrak k_+]\subseteq\mathfrak k_+,
\quad
[\mathfrak k_+,\mathfrak k_-]\subseteq\mathfrak k_-,
\quad
[\mathfrak k_-,\mathfrak k_-]\subseteq\mathfrak k_+.
\]
Here \(\mathfrak k_+\) is the vertical Lie algebra of the diagonal fibre and
\(\mathfrak k_-\) is the tangent module of the relative quotient.  In
particular, \(\mathfrak k_-\) is not a gauge Lie algebra.
\item The connected isometry group of the round diagonal vector fibre is
\[
        \operatorname{Isom}_0(S^3_{\rm diag})\cong SO(4),
\]
with spin lift
\[
        \operatorname{Sp}in(4)_{\rm fib}
        \cong SU(2)^a_L\times SU(2)^a_R,
\]
and hence
\[
        \mathfrak{isom}(S^3_{\rm diag})
        \cong
        \mathfrak{su}(2)^a_L\oplus\mathfrak{su}(2)^a_R.
\]
The diagonal scalar fibre contributes \(\mathfrak u(1)_{\rm fib}\).
\end{enumerate}
The typed Born exchange
\[
 \sigma(a,b)
 :=\Bigl(\iota_x^{-1}(\iota_p(b)),
          \iota_p^{-1}(\iota_x(a))\Bigr)
\]
descends to inversion on the quotient:
\[
 \Pi_{\rm rel}(\sigma(a,b))
 =\iota_p(b)\,\iota_x(a)^{-1}
 =\Pi_{\rm rel}(a,b)^{-1}.
\]
Algebraically, a quotient/readout scalar is Born-even precisely when
its representing function \(F\) also satisfies \(F(k^{-1})=F(k)\).
For the \(SU(2)\) factor this follows for class functions because \(k\) and
\(k^{-1}\) are conjugate; for the abelian \(U(1)\) factor it is a genuine
inversion-evenness condition and does not follow from conjugation invariance.
If a bi-invariant auxiliary metric on \(K\) is separately supplied, inversion
is an isometry of that metric.  No physical quotient metric or scalar-circle
metric is derived here.  The single diagonal orbit has the displayed compact
isometry-role presentation; it is not thereby a gauge connection or a
physical internal metric.
\end{theorem}

\begin{proof}
Transport through the already constructed \(\iota_x,\iota_p\) identifies
\(\Pi_{\rm rel}\) with \(r(a,b)=ab^{-1}\).  Paper~I-C,
Theorems~\ref{p1b:thm:relative} and~\ref{p1b:thm:symmetric}, give its free
diagonal fibres, quotient, exchange inversion and symmetric-pair brackets.
In particular \(d\Pi_{\rm rel}(A,B)=A-B\): the odd space represents the
quotient tangent and is not a gauge Lie algebra.  The admission of this
readout remains Theorem~\ref{p3b:thm:relative-readout-admission}.

Each vector fibre is one round three-sphere, whose connected isometry group
is \(SO(4)\) by Paper~I-A, Proposition~\ref{p1:prop:isom-s3-so4}, with spin
cover \(SU(2)_L\times SU(2)_R\).  The two simple summands therefore come
from the isometries of that fibre, while its scalar factor supplies the
\(U(1)\) translation role.  On \(SU(2)\), inversion preserves every
conjugacy class; on \(U(1)\), conjugation is trivial, so Born-evenness
requires the stated inversion identity itself.  The auxiliary metric clause
uses only the separately supplied bi-invariant metric.
\end{proof}

\begin{theorem}[Canonical scalar-diagonal quotient of the compact reciprocal carrier]
\label{p3b:thm:canonical-compact-seven-quotient}
In the notation of Theorem~\ref{p3:thm:born-parity-decomposition}, let
\[
 Z_{\rm sc}:=U(1)_{\rm diag}\triangleleft K_{\rm diag}
 \cong SU(2)_{\rm diag}\times U(1)_{\rm diag}
\]
be the diagonally embedded scalar factor already typed by the reciprocal
scalar role.  It is central and normal in \(K_{\rm diag}\), and its free
right action on \(\mathcal M_8\) defines the compact quotient
\[
 \boxed{X_7^{\rm cmp}:=\mathcal M_8/Z_{\rm sc}.}
\]
Equivalently, in coordinates
\((g_x,z_t;g_p,z_E)\in
SU(2)_x\times U(1)_t\times SU(2)_p\times U(1)_E\), declare the target before
evaluation and use the total readout
\[
 Q_7(g_x,z_t;g_p,z_E):=(g_x,g_p,z_tz_E^{-1})
 \in SU(2)_x\times SU(2)_p\times U(1)_{\rm rel}.
\]
The graph lift \(c\mapsto(c,Q_7(c))\) satisfies the witness-transport
criterion, its kernel pair is exactly the \(Z_{\rm sc}\)-orbit relation, and
its exact image is canonically diffeomorphic to \(X_7^{\rm cmp}\).  Hence
\(X_7^{\rm cmp}\) is a generated seven-dimensional readout-role descendant
requiring no second G0 occurrence.  This notation does not replace
\(\mathcal M_8\), which remains present with all of its factor tests.
Then \(\dim X_7^{\rm cmp}=7\), the residual group
\[
 K_{\rm diag}/Z_{\rm sc}\cong SU(2)_{\rm diag}\cong S^3
\]
acts freely on \(X_7^{\rm cmp}\), and there is a canonical principal bundle
\[
 \boxed{
 SU(2)_{\rm diag}\longrightarrow X_7^{\rm cmp}
 \longrightarrow K_{\rm rel}\cong S^3_{\rm rel}\times S^1_{\rm rel}.}
\]
The Born exchange descends to this quotient and induces the same inversion on
its base \(K_{\rm rel}\) as in
Theorem~\ref{p3:thm:born-parity-decomposition}.
\end{theorem}

\begin{proof}
The word for \(Q_7\) uses only projections and the already generated scalar
multiplication and inversion.  It is natural, typed, and selector-free.  Its
graph lift is injective and therefore conservative on the entire source
presentation.  The section
\[
 (g_x,g_p,z)\longmapsto(g_x,z;g_p,1)
\]
proves that its essential image is exactly the displayed target.  As in the
proof of Theorem~\ref{p3b:thm:relative-readout-admission}, compactness gives
GP1 admissibility and the five joint GP2 defect components vanish: no
operation is inserted, no source test or output is omitted, equal output
values are confined to the declared comparison sort, all ports retain their
types, and the graph/image/kernel closure is complete.  Paper~II-A witness
transport therefore co-realizes the readout without a new occurrence root.

If two source points have equal \(Q_7\)-values, their two \(SU(2)\)
coordinates agree.  With \(u:=z_E^{-1}z_E'\), equality of the relative scalar
coordinates gives both \(z_Eu=z_E'\) and \(z_tu=z_t'\).  Thus the two points
lie in the same diagonal scalar orbit; the converse is immediate.  This
proves that the readout kernel pair is exactly the quotient relation and that
its image is the claimed orbit-space presentation.

The factor \(Z_{\rm sc}\) is the scalar circle of the already typed
diagonal product group.  The identified orbit space, its residual free
\(SU(2)_{\rm diag}\)-action, the principal bundle and the descended Born
inversion are then Paper~I-C, Theorem~\ref{p1b:thm:relative}, applied to
this same readout.  Thus the quotient presentation retains the constructed
source and introduces no further occurrence root.
\end{proof}


\begin{corollary}[Relative scalar character lattice and normalized generator]
\label{p3b:cor:v181-relative-scalar-generator}
Use the normalized scalar angles of Paper~III-A,
\[
 \theta=2\pi t/\tper,\qquad \phi=2\pi E/\Eper,
 \qquad
 N_t=-i\partial_\theta,\quad N_E=-i\partial_\phi,
\]
and put
\[
 \chi:=\frac{\theta+\phi}{2},\qquad \psi:=\theta-\phi.
\]
The scalar component of the quotient map in
Theorem~\ref{p3b:thm:canonical-compact-seven-quotient} is
\[
 \boxed{z_{\rm rel}=z_tz_E^{-1}=e^{i\psi}.}
\]
A character $e^{i(n\theta+w\phi)}$ descends through the diagonal scalar
circle if and only if $n+w=0$.  Hence the descended character lattice is
\[
 \boxed{\Lambda_{\rm rel}=\{(n,-n):n\in\mathbb Z\}\cong\mathbb Z,}
\]
and its primitive dimensionless generator is
\[
 \boxed{N_{\rm rel}:=-i\partial_\psi=\frac12(N_t-N_E).}
\]
On descended states $N_t+N_E=0$, so
$N_{\rm rel}=N_t=-N_E$.  Born exchange sends $\psi\mapsto-\psi$ and hence
$N_{\rm rel}\mapsto-N_{\rm rel}$.
\end{corollary}
\begin{proof}
This is Paper~I-C, Corollary~\ref{p1b:cor:relative-characters}, with
\((\theta_t,\theta_E)=(\theta,\phi)\).  The difference angle is
\(\psi=\theta-\phi\), differentiated at fixed mean angle \(\chi\), so
the generator convention and its factor \(1/2\) agree.
\end{proof}

\begin{corollary}[Circle-valued relative primitive]
\label{p3b:cor:v181-relative-circle-moment}
For the calibrated scalar Poincar\'e form of Paper~III-A,
\[
 \omega_{\rm sc}=\frac{\hbar}{2\pi}\,d\phi\wedge d\theta,
\]
the diagonal generator $V_+=\partial_\theta+\partial_\phi$ obeys
\[
 \boxed{\iota_{V_+}\omega_{\rm sc}
 =\frac{\hbar}{2\pi}\,d(\theta-\phi)
 =\frac{\hbar}{2\pi}\,d\psi.}
\]
Therefore the quotient coordinate $z_{\rm rel}=e^{i\psi}$ is the canonical
circle-valued primitive of the diagonal scalar action.  Since $\psi$ is only
circle-valued globally, this does not manufacture a single-valued real moment
map or a scalar-circle metric modulus.
\end{corollary}
\begin{proof}
In the ordered coordinates \((\theta,\phi)\), the dimensionless form
of Paper~I-C satisfies
\(\omega_{\rm sc}=-2\pi\hbar\,\omega_{1B}\).  Multiplying
Proposition~\ref{p1b:prop:scalar-moments} by this factor gives the stated
contraction and the relative primitive \(e^{i\psi}\), with the
Paper~III-A phase calibration unchanged.
\end{proof}

\begin{remark}[Supplementary ratio-map consequences]
The canonical one-factor restrictions of \(\Pi_{\rm rel}\) and the
odd-jet vanishing theorem for Born-even relative scalars are recorded in
Paper~III-G, Section~\ref{p3b:app:pair-groupoid-analysis}.  They do not discharge
\((\mathrm{GRED}_7)\) or identify the compact readout with a local physical
universal-cover base.
\end{remark}


\begin{remark}[Compact quotient versus the later physical GRED carrier]
\label{p3b:rem:compact-seven-not-gred}
The scalar-diagonal quotient derives a compact seven-dimensional readout-role
descendant and its principal \(SU(2)\) orbit fibration from already
constituted reciprocal geometry.  The derivation is non-destructive:
\(\mathcal M_8\) and all factor tests remain co-realized with the readout.
The quotient \emph{by itself} does not select a connection, supply a fibre
metric or normalized reduction measure, identify \(K_{\rm rel}\) with a local
four-dimensional physical base, or fix a generator convention.  However,
Theorem~\ref{p3b:thm:parent-cartan-scalar-diagonal-descent} uses the additional
already generated parent Cartan connection to produce one unique compact
residual \(SU(2)_{\rm diag}\)-connection by basic descent.  Thus the correct
implications are
\[
 \boxed{\begin{gathered}
 X_7^{\rm cmp}\ \text{alone}\not\Rightarrow\\
 \mathrm{GRED}_7,\ \text{connection},\ \text{fibre metric},\\
 \text{local four-dimensional base or reduction measure},
 \end{gathered}}
\]
but
\[
 \boxed{
 (X_7^{\rm cmp},q_{\rm sc},\mathcal A_{M8})
 \Longrightarrow \omega_{\rm cmp}
 }
\]
for the compact principal bundle.  Even this strengthened parent-descended
connection does not by itself discharge Paper~III-F's localization/reduction
requirements: the local physical base, fibre metric, reduction measure,
normalization, fluctuation content, and action remain later typed data.
\end{remark}


\paragraph{Technical proof ownership and constituted scope.}
Every extended auxiliary derivation, explicit matrix verification,
and optional comparison moved during this reorganization is
preserved in Paper~III-G, Section~\ref{p3g:sec:spectral}.
The physical momentum carrier is inherited from Paper~II-B:
$\mathcal H_p\cong\mathcal H_1$ and $m_p=1$, without later
scalar, action, regulator, or source data reselecting it.

\section{Dynamical Cartan geometry and connection restrictions}
\label{paper3b:sec:organized-cartan}
\subsection{Cartan deformation as a connection-family constructor}
\label{p3b:sec:v199-cartan-deformation}
The nonabelian connection and solder form of the completed parent are
\(\mathcal A=(\theta_x+\theta_p)/2\) and
\(\mathcal E=(\theta_x-\theta_p)/2\).  Here the diagonal right action is
the principal action: \(\mathcal E\) is horizontal and adjoint-equivariant,
whereas \(\mathcal A\) transforms as a connection.  Splitting the two
Maurer--Cartan equations gives
\(D_{\mathcal A}\mathcal E=0\) and
\(F_{\mathcal A}=-\mathcal E\wedge\mathcal E\).
These are the local identities of
Theorem~\ref{p3b:thm:v187-m8-cartan-split}; the global anti-diagonal
limitations remain as stated there.

\begin{theorem}[Connection family and its exact curvature]
\label{p3b:thm:v199-cartan-deformation}
For every smooth real scalar \(u\) basic under the principal action,
\begin{equation}
 \mathscr A_u=\mathcal A+u\mathcal E,\qquad
 F_u=du\wedge\mathcal E+(u^2-1)\mathcal E\wedge\mathcal E.
 \label{p3b:eq:v199-cartan-curvature}
\end{equation}
The first expression is a connection on the same bundle.  On the
labelled-factor section, \(\mathcal A=\mathcal E=\Theta/2\), and hence
\begin{equation}
 A_f=f\Theta,\quad
 F_f=df\wedge\Theta+f(f-1)\Theta\wedge\Theta,
 \qquad f=(1+u)/2.
 \label{p3b:eq:v199-reduced-connection}
\end{equation}
On the momentum restriction of
Theorem~\ref{p3b:thm:v188-momentum-cartan-restriction}, the solder sign is
reversed, so \(f=(1-u)/2\).  Born exchange sends \(u\) to \(-u\), together
with reciprocal base transport, and interchanges \(f\) and \(1-f\).
\end{theorem}
\begin{proof}
The connection and curvature are Paper~I-C,
Theorem~\ref{p1b:thm:deformation}, on the same smooth real basic domain.
For the momentum inclusion the solder has the opposite sign by
Theorem~\ref{p3b:thm:v188-momentum-cartan-restriction}; hence the local
substitution is \(f=(1-u)/2\), and exchange interchanges it with
\((1+u)/2\).
\end{proof}

\begin{remark}[Evaluated nonconstant Cartan provider]
\label{p3b:rem:v2217-nonconstant-cartan-import}
For the round spatial realization, Paper~I-C,
Theorem~\ref{p1b:thm:v2217-nonconstant-cartan-dirac}, evaluates a nonconstant
curve of this smooth basic family, including its gauge projection and signed
Dirac coefficient. The coordinate $u$ remains part of the connection;
it is not an additional independently constituted scalar. A represented
time history still consumes the actual metric, action and Gauss reduction.
\end{remark}


The mixed term \(du\wedge\mathcal E\) supplies a succession-direction leg
when the configuration varies along that direction.  Paper~VII-A composes the
momentum restriction with the existing \(\iota_1\) transport to reciprocal
position before forming its native gauge links
(Theorem~\ref{p10:thm:v199-exact-cartan-links}).  The relative group alone
is not being identified with physical position.  This constructs an
admitted family and its curvature; \(u\) is a configuration coordinate,
not a newly constituted fundamental scalar or a selected G0 history.

On the retained round spatial restriction, the represented canonical
doublet connection is the Levi--Civita spin connection. Paper~III-F,
Theorem~\ref{p3f:thm:m8td-hodge-dirac}, identifies its twisted Dirac
operator with the scalar/one-form Hodge operator and gives the pointwise
Cartan perturbation for arbitrary smooth basic $u$. In the subfamily
$u(q^{-1})=-u(q)$, the resulting Born grading protects a zero mode;
Theorem~\ref{p3f:thm:m8td-protected-zero} proves uniqueness and a gap for
$\|u\|_\infty\le1$, while
Corollary~\ref{p3f:cor:m8td-constructive-protected-zero} gives a convergent
reduced-resolvent construction of that represented zero line. Its
perturbative expansion separates longitudinal from coexact one-form
dressing, and Corollary~\ref{p3f:cor:m8td-transverse-zero-response} shows
that dipole mixing with the reciprocal-odd part of scalar shell $E_k$
spans the complete reciprocal-odd coexact shell at the same vector level.
These are represented consequences on that spatial domain, with the same
reciprocal transport and configuration typing as above. They do not select
a Cartan history, promote the zero-state one-form to a gauge connection, or
constitute a new fundamental scalar occurrence.

\subsection{Complete $M_8$ Cartan split and finite defect}
\label{p3b:sec:v187-m8-cartan-defect}
Let $K\cong SU(2)\times U(1)$ and write the completed reciprocal parent as
$M_8=K_x\times K_p$.  Denote the two left Maurer--Cartan forms by
$\theta_x,\theta_p$ and define
\[
 \boxed{\mathcal A_{M8}:=\tfrac12(\theta_x+\theta_p),\qquad
 \mathcal E_{M8}:=\tfrac12(\theta_x-\theta_p).}
\]

\begin{theorem}[Parent Cartan connection--solder decomposition]
\label{p3b:thm:v187-m8-cartan-split}
Paper~I-C, Theorem~\ref{p1b:thm:cartan}, gives, under the diagonal right
action $(a,b)\mapsto(ah,bh)$,
\[
 \mathcal A_{M8}\mapsto
 \operatorname{Ad}_{h^{-1}}\mathcal A_{M8}+h^{-1}dh,
 \qquad
 \mathcal E_{M8}\mapsto
 \operatorname{Ad}_{h^{-1}}\mathcal E_{M8}.
\]
The second form vanishes on vertical vectors $(X,X)$ and evaluates to $X$ on
anti-diagonal vectors $(X,-X)$; hence it is the canonical relative solder
form.  Splitting the two Maurer--Cartan equations gives the exact parent
identities
\[
 \boxed{D_{\mathcal A}\mathcal E=0,\qquad
 F_{\mathcal A}=-\mathcal E\wedge\mathcal E.}
\]
No field equation or history is used.
\end{theorem}

Under parent exchange, \(\mathcal A_{M8}\) is even and
\(\mathcal E_{M8}\) is odd; hence \(F_{\mathcal A}\) is even.
The rank-four solder determinant has sign \((-1)^4=+1\).
Likewise exchange of the two four-dimensional \(K\) factors has
orientation sign \((-1)^{4\cdot4}=+1\): the odd \(S^3\)-exchange
and odd scalar-torus exchange cancel. The scalar phase two-form's odd
sign is therefore not the sign of a full metric volume density.


The global labelled-factor trivialization
\[
 \Phi:K_{\rm rel}\times K_{\rm diag}\to M_8,
 \qquad (k,h)\mapsto(kh,h)
\]
has inverse $(a,b)\mapsto(ab^{-1},b)$ and yields
\[
 \mathcal A_{M8}=\theta_h+\tfrac12\operatorname{Ad}_{h^{-1}}\theta_k,
 \qquad
 \mathcal E_{M8}=\tfrac12\operatorname{Ad}_{h^{-1}}\theta_k.
\]
On the canonical labelled-factor section $h=e$ both reduce to
$\tfrac12\theta_k$.  The residual vector connection below is therefore not a
new selector: it is the unique scalar-diagonal descent of the already
generated parent Cartan connection.

\begin{theorem}[Scalar-diagonal descent of the parent Cartan connection]
\label{p3b:thm:parent-cartan-scalar-diagonal-descent}
Let
\[
 q_{\rm sc}:M_8\longrightarrow X_7^{\rm cmp}=M_8/Z_{\rm sc}
\]
be the generated scalar-diagonal quotient of
Theorem~\ref{p3b:thm:canonical-compact-seven-quotient}.  Let
\[
 \pi_{\rm su}:\mathfrak k=\mathfrak{su}(2)\oplus\mathfrak u(1)
 \longrightarrow \mathfrak k/\mathfrak u(1)\cong\mathfrak{su}(2)
\]
be the canonical quotient map and set
\[
 \widehat\omega:=\pi_{\rm su}\circ\mathcal A_{M8}.
\]
Then \(\widehat\omega\) is \(Z_{\rm sc}\)-basic and therefore descends
uniquely to a principal \(SU(2)_{\rm diag}\)-connection
\[
 \boxed{\omega_{\rm cmp}\in
 \Omega^1(X_7^{\rm cmp};\mathfrak{su}(2))}
\]
satisfying
\[
 \boxed{q_{\rm sc}^{*}\omega_{\rm cmp}
 =\pi_{\rm su}(\mathcal A_{M8}).}
\]
In the canonical presentation
\(X_7^{\rm cmp}\cong SU(2)_x\times SU(2)_p\times U(1)_{\rm rel}\),
\[
 \boxed{\omega_{\rm cmp}
 =\frac12\bigl(\theta_x^{\mathfrak{su}(2)}
                 +\theta_p^{\mathfrak{su}(2)}\bigr).}
\]
The descended Born exchange fixes this connection.  Its curvature is the
projected parent curvature,
\[
 q_{\rm sc}^{*}\Omega_{\rm cmp}
 =\pi_{\rm su}(F_{\mathcal A}),
\]
and on anti-diagonal horizontal vectors
\[
 \boxed{\Omega_{\rm cmp}((X,-X),(Y,-Y))=-[X,Y].}
\]
For the local section \(s(r,z)=(r,e,z)\),
\[
 \boxed{A_B=s^*\omega_{\rm cmp}=\frac12\theta_r,\qquad
 F_B=-\frac14\theta_r\wedge\theta_r.}
\]
This theorem generates the compact residual principal connection only.  It
does not generate a local four-dimensional base, fibre metric, reduction
measure, generator normalization, fluctuation field, or physical action.
\end{theorem}

\begin{proof}
The quotient of Theorem~\ref{p3b:thm:canonical-compact-seven-quotient}
and the parent Cartan datum have already been generated above.
Apply the independent mathematical constructor in Paper~I-C,
Theorem~\ref{p1b:thm:minimal-native-provider}, to this actual diagram.
Semisimple projection kills scalar vertical values and is scalar-invariant,
so it descends uniquely through the surjective submersion.  Residual
reproduction and equivariance make the descent a principal connection.
The same theorem gives the average formula and Born invariance.

Curvature commutes with the central Lie-algebra quotient, so
\(q_{\rm sc}^*\Omega_{\rm cmp}=\pi_{\rm su}(F_{\mathcal A})\).
Using \(F_{\mathcal A}=-\mathcal E\wedge\mathcal E\) and
\(\mathcal E(X,-X)=X\) proves the horizontal formula.  The section-frame
potential and curvature are also Paper~I-C,
Theorem~\ref{p1b:thm:conditional-residual-descent}; equivalently, pull
back the average and use the Maurer--Cartan equation.

This is a proof for the actual generated source, using mathematical results
that do not assume its physical occurrence.  Together with the generated
scalar quotient it supplies the constructor stage of Paper~I-C's
provisional native datum.  Paper~III-F,
Corollary~\ref{p3f:cor:v259-native-su2-discharge}, records the exact
weak-role discharge.  No local-base, measure or physical-action conclusion
is added by the connection descent alone.
\end{proof}

\begin{theorem}[Cartan embedding and finite order-swap defect]
\label{p3b:thm:v187-cartan-embedding-commutator}
By Paper~I-C, Theorem~\ref{p1b:thm:symmetric}, the Born involution
$\sigma(a,b)=(b,a)$ generates the symmetric pair
$(K\times K,\Delta K,\sigma)$ and its canonical Cartan embedding
\[
 \boxed{c(k)=(k,k^{-1}),\qquad r(c(k))=k^2,\qquad
 (dc)_eX=(X,-X).}
\]
The finite order-swap defect
\[
 \boxed{\mu_{\rm com}(a,b)=ab(ba)^{-1}=aba^{-1}b^{-1}}
\]
is diagonal-conjugation equivariant and Born-odd by Paper~I-C,
Proposition~\ref{p1b:prop:commutator},
$\mu_{\rm com}(b,a)=\mu_{\rm com}(a,b)^{-1}$.  Near the identity,
\[
 \log\mu_{\rm com}(e^{\epsilon X},e^{\epsilon Y})
 =\epsilon^2[X,Y]+O(\epsilon^3).
\]
Paper~I-C, Theorem~\ref{p1b:thm:relational-cartan-commutator}
\cite{Kugelmann2026Paper1B}, sharpens the infinitesimal identification at the
Cartan basepoint.  At \(c(e)=(e,e)\), put
\(U_X=(X,-X)\) and \(U_Y=(Y,-Y)\).  For any local horizontal extensions
\(\widetilde U_X,\widetilde U_Y\),
\[
 F_{\mathcal A}(U_X,U_Y)=-[X,Y],\qquad
 \mathcal A([\widetilde U_X,\widetilde U_Y])=[X,Y]
 =-F_{\mathcal A}(U_X,U_Y).
\]
Away from the identity, a left-trivialized tangent to the Cartan image is
\[
 (dc)_k(kX)=\bigl(X,-\operatorname{Ad}_kX\bigr);
\]
the anti-diagonal \((X,-X)\) is the horizontal symmetric-pair complement but
is not generally tangent to the nonabelian Cartan image.  The finite
commutator satisfies only the corresponding quadratic BCH jet
\[
 \log\mu_{\rm com}(e^{\epsilon X},e^{\epsilon Y})
 =-\epsilon^2F_{\mathcal A}(U_X,U_Y)+O(\epsilon^3).
\]
Thus the Lie bracket, horizontal Frobenius defect and curvature are exactly
the same infinitesimal defect; \(\mu_{\rm com}\) is its finite group
commutator with that defect as its second-order logarithmic jet.  No exact
finite equality between curvature and \(\mu_{\rm com}\) is asserted.
Occupation of a particular commuting or conjugacy stratum remains an
occurrence claim.
\end{theorem}

\begin{remark}[Infinitesimal anti-diagonal versus global subgroup]
\label{p3b:rem:v193-antidiagonal-global-firewall}
The anti-diagonal tangent space \(\{(X,-X)\}\) is the Born-odd symmetric-pair
complement, but its Cartan image
\(c(K)=\{(k,k^{-1})\}\) is a subgroup exactly when \(K\) is abelian.  Thus
the global labelled trivialization of \(M_8\to K_{\rm rel}\) does not
integrate the nonabelian Born-odd vector direction to a second subgroup
factor.  On the scalar \(U(1)\) restriction the obstruction vanishes.
\end{remark}


\subsection{$M_8$ restriction to the scalar phase parent}
\label{p3b:sec:v190-m8-scalar-restriction}
On the $U(1)_t\times U(1)_E$ component write the normalized Maurer--Cartan
forms as $d\theta_t,d\theta_E$.  The scalar restriction of the parent Cartan
pair is
\[
 \boxed{\mathcal A_{\rm sc}=\tfrac12(d\theta_t+d\theta_E),\qquad
 \mathcal E_{\rm sc}=\tfrac12(d\theta_t-d\theta_E).}
\]
Because $\mathfrak u(1)$ is abelian,
\[
 \boxed{d\mathcal A_{\rm sc}=0,\qquad
 d\mathcal E_{\rm sc}=0,\qquad
 F_{\mathcal A,{\rm sc}}=0,\qquad
 \mu_{\rm com}^{\rm sc}=1.}
\]
Thus the scalar branch is the flat abelian restriction of the same parent
Cartan geometry whose $SU(2)$ branch carries the nonzero commutator/curvature
defect.

\begin{theorem}[Scalar Born eigenphase cover and primitive-phase host]
\label{p3b:thm:v190-scalar-eigenphase-host}
The diagonal and anti-diagonal scalar subgroups
\[
 D=\{(z,z)\},\qquad A=\{(z,z^{-1})\}
\]
obey
\[
 D\cap A=\{(1,1),(-1,-1)\}\cong\mathbb Z_2,
\]
and multiplication gives
\[
 \boxed{T^2_{tE}\cong(D\times A)/\mathbb Z_2.}
\]
Equivalently the Born eigenphase map $Q=(q_+,q_-)$ is the two-fold cover
of Paper~I-C, Theorem~\ref{p1b:thm:eigenphase}.  Its integral pullback shows
that the primitive scalar Poincar\'e class of Paper~III-A is not the pullback of a
primitive integral class on the Born-eigenphase torus, although its double
class descends.  Therefore the primitive scalar torus is the least generated
global host that simultaneously retains the two scalar roles and their
primitive phase cell.
\end{theorem}

\begin{corollary}[Born action on global scalar spin twists]
\label{p3b:cor:v190-scalar-spin-twists}
By Paper~I-C, Theorem~\ref{p1b:thm:spin}, the two $S^3$ factors contribute
no $H^1(-;\mathbb Z_2)$, and
\[
 H^1(M_8;\mathbb Z_2)\cong\mathbb Z_2^2
\]
is carried entirely by the two scalar circles.  Writing a spin-twist class as
$(\varepsilon_t,\varepsilon_E)$, Born exchange sends
\[
 (\varepsilon_t,\varepsilon_E)\longmapsto
 (\varepsilon_E,\varepsilon_t).
\]
Hence the strictly Born-fixed spin-twist classes are exactly
\[
 \boxed{(0,0),\qquad(1,1).}
\]
Reducing the eigenphase pullback matrix of Paper~I-C,
Theorem~\ref{p1b:thm:eigenphase}, modulo two gives image precisely
these two classes.  Explicitly, the target twists \((0,0),(1,1)\) pull back to
\((0,0)\), while \((0,1),(1,0)\) pull back to \((1,1)\).  Thus the two parent
image classes are distinct and the ambiguity is two-to-one on target
preimages.  This is a classification of possible global lifts, not a
physical selector of one spin structure.
\end{corollary}

\begin{remark}[Born quotient and fermionic lift firewall]
\label{p3b:rem:v190-born-quotient-spin-firewall}
The direct quotient $T^2_{tE}/\langle\mathcal B\rangle$ is the symmetric
square of a circle and hence a M\"obius band.  It is nonorientable and does not
carry an ordinary spin structure.  Thus quotienting by Born exchange cannot
silently replace the parent scalar torus in a fermionic owner; a Pin or
separately generated equivariant lift would be extra data.
\end{remark}

\subsection{$M_8$ restriction to the momentum polarization and finite supports}
\label{p3b:sec:v188-m8-momentum-restriction}

Let
\[
 \iota_x:K\to M_8,\quad a\mapsto(a,e),
 \qquad
 \iota_p:K\to M_8,\quad b\mapsto(e,b)
\]
be the labelled position and momentum factor inclusions.

\begin{theorem}[Momentum-polarized restriction of the parent Cartan pair]
\label{p3b:thm:v188-momentum-cartan-restriction}
The complete parent Cartan pair restricts as
\[
 \boxed{\iota_x^*\mathcal A_{M8}=\frac12\theta_x,
 \quad \iota_x^*\mathcal E_{M8}=\frac12\theta_x,}
\]
\[
 \boxed{\iota_p^*\mathcal A_{M8}=\frac12\theta_p,
 \quad \iota_p^*\mathcal E_{M8}=-\frac12\theta_p.}
\]
Equivalently, for $Y\in\mathfrak k$ a pure momentum tangent vector decomposes
canonically as
\[
 (0,Y)=\frac12(Y,Y)+\frac12(-Y,Y),
\]
so that the same momentum generator has a diagonal-connection component
$Y/2$ and a relative-solder component $-Y/2$.  Born exchange reverses the
relative component and preserves the diagonal one.

On the $SU(2)_p\simeq S_p^3$ factor use the Paper-I-A Maurer--Cartan coframe
$(\sigma_1,\sigma_2,\sigma_3)$ with chosen Hopf form
$\eta_0=\sigma_3/2$ and $d\sigma_3=-\sigma_1\wedge\sigma_2$.  Writing the
restricted relative coframe as $\varepsilon_i=-\sigma_i/2$ gives the exact
identity
\[
 \boxed{d\eta_0=-\frac12\sigma_1\wedge\sigma_2
       =-2\,\varepsilon_1\wedge\varepsilon_2.}
\]
Thus the Hopf/contact bracket-generating defect on the momentum sphere is the
Hopf-polarized restriction of the same quadratic solder defect that globally
appears as $F_{\mathcal A}=-\mathcal E\wedge\mathcal E$.  No finite carrier
or occurrence statement is used.
\end{theorem}

\begin{proof}
The two labelled factor restrictions are Paper~I-C,
Proposition~\ref{p1b:prop:finite-defect}.  On the retained momentum
sphere, the Paper-I-A convention \(\eta_0=\sigma_3/2\) and
\(d\sigma_3=-\sigma_1\wedge\sigma_2\), with
\(\varepsilon_i=-\sigma_i/2\), gives the stated contact identity.
The parent structure equation in Theorem~\ref{p1b:thm:cartan} identifies
its quadratic solder origin.
\end{proof}

\begin{theorem}[Finite support restriction of the Cartan square and order-swap defect]
\label{p3b:thm:v188-finite-defect-restriction}
Let $\Gamma\le SU(2)_p$ be any retained subgroup support.  The intrinsic
parent group words restrict without a new selector:
\[
 c_\Gamma(g)=(g,g^{-1}),\qquad
 r(c_\Gamma(g))=g^2,
\]
\[
 \mu_\Gamma(a,b)=aba^{-1}b^{-1}\in\Gamma.
\]
The subgroup generated by the values of $\mu_\Gamma$ is exactly the derived
subgroup $[\Gamma,\Gamma]$.

For $\Gamma\in\{2T,2O,2I\}$, whose center is $Z=\{\pm1\}$, both operations
are insensitive to changing a lift by the central sign and therefore factor
through the projective quotient in the set-theoretic sense
\[
 s_\Gamma:\Gamma/Z\to\Gamma,\qquad [g]\mapsto g^2,
\]
\[
 \bar\mu_\Gamma:(\Gamma/Z)^2\to\Gamma,\qquad
 ([a],[b])\mapsto aba^{-1}b^{-1}.
\]
They satisfy
\[
 \pi(s_\Gamma([g]))=[g]^2,
 \qquad
 \pi(\bar\mu_\Gamma([a],[b]))=[[a],[b]].
\]
The lift can nevertheless retain a nontrivial central defect invisible in the
projective quotient.  In the common quaternion subgroup,
\[
 \boxed{\bar\mu([i],[j])=[i,j]=-1,}
\]
although the projective half-turns $[i]$ and $[j]$ commute in the Klein four
subgroup of $SO(3)$.  Hence the parent order-swap defect detects the spinorial
central holonomy erased by projectivization.
\end{theorem}

\begin{proof}
Apply Paper~I-C, Theorem~\ref{p1b:thm:symmetric} and
Proposition~\ref{p1b:prop:finite-defect}, to the retained subgroup
\(\Gamma\); its closure under the group words gives their restrictions.
For central signs \(z,w\), \((zg)^2=g^2\) and
\([za,wb]=[a,b]\), proving representative independence.  Projection gives
the corresponding words in \(\Gamma/Z\).  Finally \([i,j]=-1\) for
quaternion units, while their projective half-turns commute in the Klein
four group, which proves the asserted central defect interpretation.
\end{proof}

\begin{corollary}[Exceptional finite defect filtration]
\label{p3b:cor:v188-exceptional-defect-filtration}
Define $\mathsf D^0(\Gamma)=\Gamma$ and
$\mathsf D^{n+1}(\Gamma)=[\mathsf D^n(\Gamma),\mathsf D^n(\Gamma)]$.
Using the exceptional derived-subgroup identities of Paper~I-B,
\[
 \boxed{2O\to2T\to Q_8\to\{\pm1\}\to1,}
\]
\[
 \boxed{2T\to Q_8\to\{\pm1\}\to1,}
\qquad
 \boxed{2I\to2I.}
\]
Thus $2O/2T\cong C_2$ is precisely the outer abelian residue of the
normalizer, $2T/Q_8\cong C_3$ is precisely the outer abelian residue of the
Hurwitz core, and $2I$ is defect-closed/perfect.  These are restrictions of
one $M_8$ parent operation; they do not turn the exceptional fork into a
subgroup chain and do not assert occurrence of any particular finite support.
\end{corollary}

\begin{proof}
The restrictions in Theorem~\ref{p3b:thm:v188-finite-defect-restriction}
use the derived-subgroup identities of Paper~I-C,
Proposition~\ref{p1b:prop:finite-defect}, together with the exceptional subgroup identities of Paper~I-B.  Iteration gives the three displayed series for those supports;
it does not select an occurring support.
\end{proof}


\section{Born-even relative connection: Yang--Mills and spectral consequences}
\label{p3b:sec:v346-born-meron}
The compact relative $SU(2)$-connection already established by
Theorem~\ref{p3b:thm:parent-cartan-scalar-diagonal-descent} has a further
\emph{mathematical} consequence on the round relative three-sphere.  The
Yang--Mills action used in the following proposition is a comparison
functional, not an assertion that the four-dimensional physical action
contains an independently propagating gauge copy on this quotient.

\begin{theorem}[Born meron is a nonflat stationary Yang--Mills connection]
\label{p3b:thm:v346-born-meron-ym}
Write $\theta=r^{-1}dr$ on round $SU(2)_{\rm rel}$ and
$A_\lambda=\lambda\theta$ for real constant $\lambda$.  Then
\[
 F_\lambda=\lambda(\lambda-1)\theta\wedge\theta,
 \qquad \operatorname{CS}(A_\lambda)
 =3\lambda^2-2\lambda^3\pmod{\mathbb Z},
\]
where the Chern--Simons convention assigns unit winding to $A_1=\theta$.
For the bi-invariant round metric the source-free Yang--Mills equation is
\[
 \boxed{D_{A_\lambda}\star F_\lambda=0
 \quad\Longleftrightarrow\quad
 \lambda\in\{0,\tfrac12,1\}.}
\]
The Born-even value $A_B=A_{1/2}$ is nonflat and has
$\operatorname{CS}(A_B)=1/2\pmod{\mathbb Z}$.  Along this homogeneous path
$\mathcal E_{\rm YM}(\lambda)=N\lambda^2(1-\lambda)^2$, with $N>0$,
and $\mathcal E_{\rm YM}''(1/2)=-N<0$; it is therefore a
\emph{saddle}, not a stable Yang--Mills vacuum.
\end{theorem}
\begin{proof}
Maurer--Cartan gives $d\theta=-\theta\wedge\theta$.
Bi-invariance and three-dimensional Hodge duality give
$\star(\theta\wedge\theta)=k\theta$ for a nonzero convention-dependent
constant $k$, whence
$D_{A_\lambda}\star F_\lambda
 =k\lambda(\lambda-1)(2\lambda-1)\theta\wedge\theta$.
The Yang--Mills norm is proportional to
$\lambda^2(1-\lambda)^2$.  Finally substitute $A_\lambda$ into
$\operatorname{tr}(A\wedge dA+\frac23 A\wedge A\wedge A)$, divide by its
value at $A_1$, and obtain $3\lambda^2-2\lambda^3$.
Detailed normalization and the associated index construction are in
Paper~III-G, Section~\ref{p3g:sec:v346-meron-proof}.
\end{proof}

\begin{corollary}[Central Wilson loop and representation-parity spectral shift]
\label{p3b:cor:v346-born-meron-wilson}
On $r(\varphi)=e^{i\varphi\sigma_3}$,
$0\le\varphi\le2\pi$, the Born parallel transport in the fundamental
representation is
\[
 \boxed{\operatorname{Hol}_{1/2}(A_B)
 =e^{-i\pi\sigma_3}=-I_2.}
\]
Its action in spin $j$ is $(-1)^{2j}I_{2j+1}$.  In the gauge induced by
this closed circle, the restricted covariant Laplacian has mode values
$\lambda_{n,m}=(n+m)^2$, $n\in\mathbb Z$ and
$m=-j,-j+1,\ldots,j$.  This is a loop restriction, not the full
three-dimensional twisted Dirac spectrum and not a particle mass.
\end{corollary}
\begin{proof}
The pullback is $A_B=i\sigma_3 d\varphi/2$ and exponentiation gives
$-I_2$.  The central element acts by $(-1)^{2j}$ and the periodic
Fourier modes of $\partial_\varphi+iJ_3$ give the eigenvalue formula.
A physical observable requires a licensed representation and actual
coupling of a field to this compact connection.
\end{proof}

\clearpage
\section{Generated spectral and action principles}
\label{paper3b:sec:organized-principles}
\subsection{Dynamical Principles Beyond the Kinematic Completion}
\label{p3:sec:dyn-principles}

Table~\ref{tab:p3-dcl-ledger} states the finite intrinsic action problem discharged in this section.

\begin{table}[H]
\centering
\small
\begingroup
\setlength{\tabcolsep}{3pt}
\renewcommand{\arraystretch}{1.15}
\begin{tabular}{@{}
  L{3.0cm}
  L{9cm}
  L{\dimexpr\linewidth-12cm-4\tabcolsep\relax}
@{}}
\toprule
\textbf{Ledger row}
& \textbf{Claim}
& \textbf{Discharged by} \\
\midrule
Scalar realization operation
& A completed reciprocal geometry without a scalar valuation is not yet a physical dynamics; the valuation problem is the Paper-III analogue of the Paper-II operation-closure tests.
& Theorem~\ref{p3:thm:scalar-valuation-necessity} \\

Kinematic covariance
& The completed geometry has no preferred geometric coordinate or
\(\operatorname{Sp}in(4)\)-related frame; full-unitary spectral naturality follows later
from the unmarked closed occurrence.
& Theorem~\ref{p3:thm:spectral_covariance} \\

Intrinsic class \(\Dcl\)
& \(\Dcl\) is the loop-derived positive Dirac-even single-process spectral form class.
& Definition~\ref{p3:def:D1}, Theorem~\ref{p3:thm:gp2-dcl-derivation} \\

No hidden insertion
& Every additional channel profile or mark requires generated provenance;
the occurring common process itself remains free.
& Theorem~\ref{p3:thm:no_residual_regularity_D1} \\

No mixed vector contraction
& Before diagonal identification there is no invariant mixed vector contraction.
& Lemma~\ref{p3:lem:no-invariant-mixed-vector-contraction} \\

Sectorwise additivity
& Sectorwise additivity follows from the licensed direct sum and loop trace.
& Theorem~\ref{p3:thm:sectorwise_additivity} \\

Shellwise coarse graining
& Coarse graining is imported only through the Paper-II shellwise theorem.
& Theorem~\ref{p3:thm:shellwise_coarse_graining} \\

Optional full-cell scalar profile
& On an actualized \(\mathsf{Supp}_{2T}\) core-support role, degree six is the first available nonconstant fully cell-symmetric scalar profile; physical activation is separate.
& Proposition~\ref{p3:prop:first-symmetric-finite-cell-correction} \\
\bottomrule
\end{tabular}
\caption[Intrinsic spectral-action ledger]{Intrinsic spectral-action ledger.}
\label{tab:p3-dcl-ledger}
\TableNotes{This table marks the transition from kinematic constitution to
dynamics.  On an enlarged-profile branch a closed process token is
G0 data and its outer trace form is derived.  On the native heat branch,
Theorem~\ref{p3b:thm:v173-joint-heat-process-corealization} generates the
complete process role from its self-adjoint operator owner; only an actually
performed token remains instantiational.}
\endgroup
\end{table}


The completed kinematic package \(\mathcal A_{\mathrm{book}}\) combines:
\[
\begin{gathered}
\text{the derived scalar Hamiltonian/Born constitution,}\\
\text{the imported scalar classification theorem }
\ref{p3:thm:scalar_classification_import},\\
\text{and the compact-cell constitution results of }
\text{Sec.~\ref{p3:sec:m_selection}}.
\end{gathered}
\]
Together these determine the compact reciprocal geometry \(\mathcal M_8\)
and its fixed physical momentum carrier \(\mathcal H_p\cong\mathcal H_1\).
When a reciprocal-position frontier actually occurs, its finite datum is the
Born-derived spectral window \(\mathcal P_x(m_x)\); \(m_x\) labels that
window and does not constitute a second finite position module or microscopic
carrier.
To derive dynamics one must close the finite intrinsic action problem on this
completed geometry.  Here an ``action'' means a scalar realization valuation on
the already constituted finite reciprocal datum, not an independently postulated
Lagrangian or a claim about arbitrary dynamics.  This does not alter the compact
microscopic arena; it determines which invariant functional is physically
admissible on the completed geometric spectral datum, how a smooth macroscopic
position-space description emerges, and how the shellwise coarse-grained flow is
reduced.

\begin{theorem}[Scalar realization operation]
\label{p3:thm:scalar-valuation-necessity}
The physical momentum seed and reciprocal basepoint are fixed upstream by Paper~II-B.
The normalized restriction of the separately declared additive Born trace
class to the two-dimensional reciprocal centre provides an independent
stability check.
Full physical dynamics then requires a scalar valuation of the completed
intrinsic spectral datum.  The
remaining action problem is therefore to classify the common functional
calculus profile and its finite-window extension, not to reselect the already
derived value \(\chi=1\).
\end{theorem}

\begin{proof}
The upstream kinematic constitution determines the microscopic geometry and finite momentum support.  Inside the adopted
additive Born trace class, the normalized central restriction gives its
reciprocal stability valuation, but it does not determine a common cutoff
profile on the remaining spectral clusters, define field equations, or
produce local curvature coefficients.  The geometry therefore leaves an
action-profile problem open at this chronological stage.  On the separate enlarged-profile branch the
occurrence \((\mathrm{CLOOP}_F)\) supplies one physical closed process and
Theorem~\ref{p3:thm:gp2-dcl-derivation} derived its \(\Dcl\) form.  The
later-stage theorem
\ref{p3b:thm:v173-joint-heat-process-corealization} removes that separate
\emph{role-level} occurrence on the native heat branch once the generated
joint operator parent is available.  A genuinely enlarged non-heat profile or
an actually performed process/readout token remains separately typed.
\end{proof}

\begin{theorem}[Kinematic equivariance inherited from Paper II]
\label{p3:thm:spectral_covariance}
\leavevmode\par\noindent
For the completed compact Born/contact geometry, every \(\operatorname{Sp}in(4)\)-equivariant
scalar observable is independent of geometric coordinates and frames related
by the constituted \(\operatorname{Sp}in(4)\) action.  Equivalently, the completed compact
geometry has no preferred geometric coordinate or \(\operatorname{Sp}in(4)\)-related frame.
This statement is not invariance under every unitary on a later Hilbert-space
realization.
\end{theorem}

\begin{proof}
This is the direct product-space lift of Paper~II,
Theorem~\ref{p2:thm:no-preferred-basis}, restricted to equivariant scalar
observables.  Each reciprocal vector sector inherits that finite geometric
equivariance statement, and the scalar pair contributes homogeneous circle
factors with no preferred origin.  It follows that geometric coordinate and
frame choices related by the constituted symmetry do not change such an
observable.  Full-unitary naturality of an action is not a kinematic
conclusion here; it follows only after the unmarked closed-readout occurrence
of Theorem~\ref{p3:thm:gp2-dcl-derivation}.
\end{proof}


\begin{remark}[Kinematic naturality versus a physical readout occurrence]
\label{p3:rem:why-action-spectral}
The completed Born/contact geometry supplies intrinsic kinetic and Dirac
operators together with finite spectral windows.
Theorem~\ref{p3:thm:spectral_covariance} shows that no geometric coordinate
or \(\operatorname{Sp}in(4)\)-related frame is kinematically preferred.  This makes
spectral processes available, but does not by itself say which enlarged
process is physically read.  Paper~II-A exhausts the primitive outer readout
roles and Paper~III-A imports the loop/boundary handoff.  On the later native
joint-heat branch, however, functional calculus and the generated heat law
supply the complete process role without another occurrence; GP1--GP3 still do
not actualize a performed token or select a law outside that minimal branch.
\end{remark}

\begin{definition}[Dynamical readout atlas]
\label{p3:def:dynamical-branch-atlas}
For generated finite spectral data, the outer provenance and its generated
decorations are
\[
\begin{array}{c|l}
\Rloop&\text{unmarked closed process readout},\\
\Rbdry&\text{prepared, tested, sourced, or external-field readout},\\
\mathcal D_{\rm corr}&\text{generated multiloop or cross-block composite},\\
\mathcal D_{\rm gr}&\text{loop with a generated grading or orientation mark},\\
\mathcal D_{\rm ch}&\text{loop with generated physical channel projectors}.
\end{array}
\]
The last three rows are marked subclasses or generated composites, not new
primitive outer roles.  On an enlarged-profile branch G0 supplies the actual
process/readout occurrence.  On the native heat branch a generated
self-adjoint operator owner supplies the complete process role by
Theorem~\ref{p3b:thm:v173-joint-heat-process-corealization}; G0 is then needed
only for an actually performed process/readout token.  GP2 excludes hidden
insertions and omissions relative to the frozen constructor data.
\end{definition}

\begin{definition}[Derived closed spectral form class
\texorpdfstring{$\Dcl$}{D1}]
\label{p3:def:D1}
For an unlabelled finite spectral datum \((\mathcal H,D,\Lcut)\), define
\(\Dcl\) to be the form class
\[
 \boxed{\begin{gathered}
 \Dcl:=\biggl\{S_F(D)=\operatorname{Tr}_{\mathcal H}
 F(D^2/\Lcut^2):\\
 F\text{ is one generated positive admissible}\\
 \text{spectral process}\biggr\}.
 \end{gathered}}
\]
Here ``one process'' means that the same generated functional-calculus law is
transported to every unlabelled reciprocal block.  In the generic enlarged
atlas the function \(F\), its process-token occurrence, its regularity domain,
and the absolute scale calibration are not part of the outer-loop theorem.
On the native heat branch, Theorem~\ref{p3b:thm:v173-joint-heat-process-corealization}
derives the role-level process from a generated self-adjoint operator owner and
fixes \(F(u)=e^{-u}\) up to the existing rate--scale orbit.  A fixed non-scalar endomorphism, grading,
projector, source, channel label, or sector-dependent process law is an
additional generated datum and is typed accordingly.
\end{definition}

\begin{theorem}[Closed-loop derivation of \texorpdfstring{$\Dcl$}{D1}]
\label{p3:thm:gp2-dcl-derivation}
Let \(D=D^*\) be the intrinsic finite-window Dirac operator and assume the G0
occurrence
\[
 (\mathrm{CLOOP}_{F}):\quad
 W_F:=F(D^2/\Lcut^2):\mathcal H\to\mathcal H,
 \qquad
 r_F:\mathbf1\to\mathbf1,
\]
where \(W_F\) is one generated positive process and \(r_F\) is the scalar
output obtained by applying precisely the licensed canonical duality
closing word
\[
 r_F=\operatorname{ev}_{\mathcal H}\circ
 (1_{\mathcal H^*}\otimes W_F)\circ\operatorname{coev}_{\mathcal H},
\]
with the tensor symmetry understood. Assume \(\operatorname{FSR}_{\Sigma}(r_F)\), empty external
boundary, unmarked closure provenance, and a licensed direct-sum constructor
on disjoint finite spectral blocks.
Then its scalar readout is forced to be
\[
 \boxed{S_F(D)=r_F
 =\operatorname{Tr}_{\mathcal H}F(D^2/\Lcut^2)\in\Dcl.}
\]
It is unitary-natural, positive, Dirac-even, and additive on the licensed
blocks.  The theorem derives those properties and the single-profile trace
form from the occurrence; it does not select \(F\) or assert that every
admissible \(F\) occurs.
\end{theorem}

\begin{proof}
By hypothesis the scalar output \(r_F\) of the licensed closure constructor,
rather than the open endomorphism \(W_F\), satisfies the finite
structural-readout predicate and has empty external boundary.  Paper~II-A
Theorem~\ref{p2:thm:loop-boundary-readout-exhaustion} therefore puts that
closure in \(\Rloop\). The specified canonical duality word then gives
the ordinary trace by Lemma~\ref{p2:lem:canonical-finite-duality-trace}
and Theorem~\ref{p2:thm:unmarked-natural-closure}.  Hence
\(r_F=\operatorname{Tr}_{\mathcal H}W_F\); the trace formula is a conclusion,
not an antecedent.  Functional calculus makes
\(W_F\ge0\), so its trace is nonnegative.  Because \(D\) enters only through
\(D^2\), replacing \(D\) by \(-D\) changes nothing.  On licensed disjoint
blocks,
\[
 F((D_1\oplus D_2)^2/\Lcut^2)
 =F(D_1^2/\Lcut^2)\oplus F(D_2^2/\Lcut^2),
\]
and finite trace linearity gives additivity.  One transported process
\(F\) gives one common profile.  Any extra endomorphism or independent profile
would be a retained mark, so it lies outside the unmarked occurrence.  These
are exactly the defining form conditions of \(\Dcl\).
\end{proof}


\begin{theorem}[Ordinary closed-loop trace detects replicated finite carriers]
\label{p3b:thm:v295-trace-owner-multiplicity-firewall}
Let $\mathcal A_0=\bigoplus_{c=1}^r\End(W_c)$ be a finite
semisimple test algebra with $d_c=\dim W_c>0$, and consider a
faithful represented implementation
$H=\bigoplus_c(W_c\otimes\mathbb C^{k_c})$, $k_c\geq1$.
Let $j(a)=\bigoplus_c a_c\otimes I_{k_c}$ and
$H_0=\bigoplus_c W_c$.  For the \emph{ordinary carrier traces},
\[
 \boxed{\Tr_H j(a)=\sum_c k_c\Tr_{W_c}(a_c),\qquad
 \Tr_{H_0}(a)=\sum_c\Tr_{W_c}(a_c).}
\]
Equality of these two unnormalized scalar functionals on all of
$\mathcal A_0$ holds exactly when every $k_c=1$.  Equality of the
\emph{normalized} carrier traces on all of $\mathcal A_0$ holds exactly
when all $k_c$ have one common value $k$ (not necessarily one).
Sector-conditioned normalized traces are always equal.

Suppose in addition that the \emph{same operator owner and profile}
are compared on both implementations, with
$D=\bigoplus_c(D_c\otimes I_{k_c})$, $D_c=D_c^*$, and
$D_0=\bigoplus_cD_c$. Then for every common finite functional-calculus
profile $F$,
\[
 \Tr_H F(D^2/L^2)=\sum_c k_c\Tr_{W_c}F(D_c^2/L^2).
\]
For $F(u)=e^{-tu}$ and $t>0$, this scalar is strictly greater than
$\Tr_{H_0}e^{-tD_0^2/L^2}$ whenever any $k_c>1$; hence a single
\emph{same-owner, unnormalized finite heat trace} separates every
nontrivial replicated implementation in this restricted class.

The state/effect quotient of
Theorem~\ref{p2b:thm:v293-operational-multiplicity-quotient}
therefore does not automatically preserve a \emph{closed-loop scalar
trace port}.  If such a port occurs as part of the frozen physical
signature under Theorem~\ref{p3:thm:gp2-dcl-derivation}, its
normalization and scalar output must be preserved by a putative GP3
whole-candidate comparison.  The theorem does not assert that every
finite presentation has that physically occurring closed-loop port or
that the compared operator owners have been identified.
\end{theorem}
\begin{proof}
Taking tensor-product traces gives the first formula. Test the
central unit $e_c$ of each simple block.  Its two unnormalized values
are $k_cd_c$ and $d_c$, so equality for every $c$ forces $k_c=1$.
Writing $N=\sum_c k_cd_c$ and $N_0=\sum_cd_c$, normalized equality
on $e_c$ is $k_cd_c/N=d_c/N_0$ and thus
$k_c=N/N_0$ for all $c$. The converse follows by substitution.
Conditioning a block trace divides by $k_cd_c$, canceling $k_c$.
Functional calculus commutes with tensoring by an identity, which
proves the spectral formula. Each finite heat trace
$\Tr_{W_c}e^{-tD_c^2/L^2}$ is strictly positive, so a positive
replication surplus gives a strictly positive trace difference.
The last qualification follows from the distinction between a
state/effect evaluation and the separately licensed scalar duality
closing word of Theorem~\ref{p3:thm:gp2-dcl-derivation}.
\end{proof}


\begin{theorem}[Two common-owner heat profiles reconstruct coordinate-block multiplicities]
\label{p3b:thm:v296-two-heat-multiplicity-reconstruction}
Suppose the \emph{fully represented} intrinsic-coordinate algebra of
Theorem~\ref{p2c:thm:v294-native-multiplicity-algebra} has the two
components $\End(W_g)\oplus\End(W_s)$ of dimensions $11$ and $4$,
with finite multiplicities $k_g,k_s\ge1$. Suppose one specified
same-owner self-adjoint spectral operator is
$D=(D_g\otimes I_{k_g})\oplus(D_s\otimes I_{k_s})$ and that the
scalar readouts at two supplied parameters $t_1,t_2>0$ are ordinary
unnormalized traces of this operator with the \emph{same} scale $L$.
Set
\[
 h_c(t)=\operatorname{Tr}_{W_c}\!\left(e^{-tD_c^2/L^2}\right),
 \qquad y(t)=\operatorname{Tr}_{H}e^{-tD^2/L^2},
 \quad c\in\{g,s\}.
\]
Then
\[
 \binom{y(t_1)}{y(t_2)}=
 \begin{pmatrix}h_g(t_1)&h_s(t_1)\\h_g(t_2)&h_s(t_2)\end{pmatrix}
 \binom{k_g}{k_s}.
\]
For arbitrary real multiplicity weights this pair determines both
weights uniquely exactly when
\[
 \Delta(t_1,t_2)=h_g(t_1)h_s(t_2)-h_g(t_2)h_s(t_1)\ne0,
\]
in which case
\[
 k_g=\frac{y(t_1)h_s(t_2)-y(t_2)h_s(t_1)}{\Delta},\qquad
 k_s=\frac{h_g(t_1)y(t_2)-h_g(t_2)y(t_1)}{\Delta}.
\]
A pair with $\Delta\ne0$ exists if and only if the finite positive
spectral counting measures of $D_g^2/L^2$ and $D_s^2/L^2$ are not
proportional. In that case $\Delta\ne0$ for generic pairs of positive
parameters. The claim of real-linear unique reconstruction is stronger
than needed for nonnegative integer weights; $\Delta=0$ does not by
itself rule out uniqueness after imposing integer constraints.
This theorem does \emph{not} generate two occurring G0 closed-loop
profiles, identify a native owner, or force $k_g=k_s=1$; the trace
occurrence and GP3 comparisons of
Theorem~\ref{p3b:thm:v295-trace-owner-multiplicity-firewall}
remain separate.
\end{theorem}
\begin{proof}
Functional calculus and the tensor-product trace prove the displayed
$2\times2$ system; Cramer's rule gives the reconstruction. The
functions $h_c(t)$ are finite Laplace transforms of positive finite
atomic counting measures. If they are proportional for all $t>0$,
subtract the corresponding finite measures; ordering their distinct
atoms and repeatedly multiplying the vanishing Laplace transform by
the exponential of the smallest remaining eigenvalue, then taking
$t\to\infty$, proves each signed atomic coefficient zero. Thus their
spectral counting measures are proportional. Conversely proportional
measures give proportional heat functions and $\Delta\equiv0$.
If the functions are not proportional, there are $t_1,t_2>0$ with
nonzero determinant. Analyticity and the identity theorem show its
zero set has empty interior (indeed measure zero) in the positive
parameter quadrant. The occurring-owner qualifications do not follow
from this algebraic calculation.
\end{proof}
\begin{definition}[Pre-response scalar-action coefficient compiler]
\label{p3b:def:scalar-action-coefficient-compiler}
Let $S_\star$ be one occurring scalar action in $\Dcl$, with its complete
local density, operator, trace/localizer and all retained marks. Write
$\mathcal D_{\rm act}$ for its already generated positive action-density
line. Its coefficient pairing is
\[
 \operatorname{ev}_{\rm act}:
 \mathcal D_{\rm act}^{-1}\otimes\mathcal D_{\rm act}\longrightarrow\mathbf1.
\]
This is the density line of this action; it is not identified by notation
with the Hopf/contact determinant line of Paper~III-C. If a scheme/depth
family of the actual action is present, its density and coefficient
connections are transported as duals. Let $U_0=\Lambda_{\rm EH}^2>0$ be the
normalized upper EH coefficient of this same action at the identity stage.
The response constructor retains that section and its actual coefficient
family. It does not choose a connection or a character weight on a separately
rooted determinant owner.
\end{definition}

\begin{theorem}[Generated action-coefficient dual and identity-stage normalization]
\label{p3b:thm:generated-action-response-owner}
On the frozen action-density pairing of
Definition~\ref{p3b:def:scalar-action-coefficient-compiler},
\begin{equation}
 \label{p3b:eq:generated-action-response-owner}
 \boxed{\mathcal L_U=\mathcal D_{\rm act}^{-1},\qquad
 a_{\rm coeff}=\operatorname{id}_{\mathcal D_{\rm act}^{-1}}.}
\end{equation}
The identity-stage coefficient continuation retains
\begin{equation}
 \label{p3b:eq:generated-action-response-initial-value}
 \boxed{\mathcal R_\star(0)=U_0=\Lambda_{\rm EH}^2.}
\end{equation}
These statements fix the coefficient slot and its initial section. They do
not identify $\mathcal D_{\rm act}$ with the contact determinant line,
select a scheme/depth transport, or prove that the actual coefficient follows
$e^{-\Phi}$. That further claim is exactly the executable certificate below.
\end{theorem}
\begin{proof}
The scalar evaluation of a density in $\mathcal D_{\rm act}$ has its
coefficient in the dual line. The frozen pairing therefore gives the
identity map to this dual; no independently normalized duplicate line is
needed. Exact continuation preserves its specified identity section $U_0$.
Neither assertion determines the connection on the density family: any
connection on a line paired with its dual preserves evaluation. Thus a
contact-character identification requires its own normalized intertwining
square, rather than an inference from duality.
\end{proof}

\begin{theorem}[Actual-action derivative and exact contact-character certificate]
\label{p3b:thm:v2211-action-contact-certificate}
Let $B_{\rm sd}$ be a connected simply connected scheme/depth domain, with
$\Phi(\sigma,q)=qL_1^{(\sigma)}$ and basepoint $b_0$. On a fixed finite
pairing presentation of the actual action let
\[
 S_b=\operatorname{Tr}\bigl(M_b e^{-B_b}\bigr),\qquad B_b=B_b^*,
\]
where $M_b$ retains the actual trace density, localizer, support and every
insertion. Assume $M_b,B_b$ are $C^1$. For a parameter direction $A$, the
complete derivative is
\begin{equation}
\label{p3b:eq:v2211-complete-action-duhamel}
 \boxed{\partial_A S_b=
 \operatorname{Tr}(M_Ae^{-B})
 -\int_0^1\operatorname{Tr}\bigl[
 M e^{-(1-s)B}B_Ae^{-sB}\bigr]ds.}
\end{equation}
The same formula applies on a trace-class parent only with its stated
trace-norm differentiation domain. If a licensed EH coefficient extraction
$\mathscr P_{\rm EH}$ is present and commutes with this derivative on a
uniform asymptotic/effective-action domain, put
\[
 c_{\rm EH}(b)=\mathscr P_{\rm EH}S_b>0,\qquad
 \mathcal U(b)=8c_{\rm EH}(b),\qquad \mathcal U(b_0)=U_0.
\]
Then the actual normalized coefficient has the contact character if and
only if the following complete residual vanishes:
\begin{equation}
\label{p3b:eq:v2211-action-contact-residual}
 \boxed{\eta_U:=d\log\mathcal U+d\Phi=0
 \quad\Longleftrightarrow\quad
 \mathcal U(b)=U_0e^{-[\Phi(b)-\Phi(b_0)]}.}
\end{equation}
In a different generated line presentation, a candidate positive owner map
$a_U$ must instead pass
$\nabla_H(a_Us)=a_U\nabla_Us$, with its retained root normalization,
host/mark naturality and actual operator origin. In local components the
connection defect is $d\log a_U+d\Phi-\omega_U$. Vanishing of that defect
and the root comparison is equivalent to the transport square; a freely
chosen depth-dependent gauge is not a generated ancestry certificate.
We denote an evaluated, origin-certified successful comparison by
$(\mathrm{ACD}_H)$. This is an exact test on an existing action, not a new
postulate that its residual is zero.
\end{theorem}
\begin{proof}
For finite matrices the power series or differentiation of the product
exponential gives
\[
 d e^{-B}[B_A]=-\int_0^1e^{-(1-s)B}B_Ae^{-sB}ds.
\]
The product rule retains $M_A$ and gives the first formula, with no
commutation assumption on $M,B,B_A$. Applying the admitted coefficient
extractor gives the actual coefficient derivative; no fixed-window
small-time coefficient is inferred. Positivity permits its logarithm.
Equation $\eta_U=0$ states that
$d\log(\mathcal U e^{\Phi})=0$; connectedness and the base value give the
unique displayed coefficient, and differentiation gives the converse.
For the line version, direct substitution into $\nabla_H(a_Us)-a_U\nabla_Us$
gives $a_U(d\log a_U+d\Phi-\omega_U)s$. Vanishing is necessary and,
with the base comparison on this simply connected domain, sufficient for
transport compatibility. Naturality and constructor origin are separate
necessary parts of a physical use; this differential identity does not
manufacture them.
\end{proof}

\begin{corollary}[Minimal-heat evaluation and the surviving counterfamily]
\label{p3b:cor:v2211-heat-contact-counterfamily}
On the scalar heat/EH comparison domain with fixed process and fixed
ambient trace convention,
\[
 \mathcal U(b)=\frac{f_{2,g}}6\Lambda_g(b)^2,\qquad
 f_{2,g}=(4\pi)^{-2}.
\]
Hence $(\mathrm{ACD}_H)$ holds exactly when
\[
 \boxed{\Lambda_g(b)^2=\Lambda_g(b_0)^2
 e^{-[\Phi(b)-\Phi(b_0)]}.}
\]
For every $\lambda>0$, the actual smooth heat family with
$\Lambda_{g,\lambda}^2=\Lambda_g(b_0)^2e^{-\lambda[\Phi-\Phi(b_0)]}$
has the same upper coefficient and positive common profile, but
\[
 \eta_{U,\lambda}=(1-\lambda)d\Phi.
\]
A fixed cutoff instead has $\eta_U=d\Phi$. Thus common heat, duality,
initial normalization and GP1 finiteness on compact parameter windows do not
make the residual vanish. If the complete actual coefficient contains other
terms, its full extractor and derivative, not this scalar subchannel, must
be evaluated. A bare geometric contact residue cannot replace that work.
\end{corollary}
\begin{proof}
The scalar $a_2$ coefficient gives the first formula in the declared
four-dimensional smooth-parent regime. Its fixed positive moment cancels in
the normalized ratio. Substitute each cutoff trajectory into
\eqref{p3b:eq:v2211-action-contact-residual}. The finite-window heat matrices
remain positive and continuous with bounded image on compact parameter
windows. Their different coefficient trajectories therefore provide genuine
counterfamilies to any deduction of $(\mathrm{ACD}_H)$ from pairing or the
single heat profile alone. They are not counterexamples after that evaluated
certificate has been supplied.
\end{proof}

\begin{remark}[Independent response and contact ancestry]
\label{p3b:rem:independent-response-atlas}
The initial section $U_0$ is fixed on the frozen coefficient branch, whereas
$U_\kappa=\kappa U_0$ changes that initial assignment. In contrast the
weight-$\lambda$ counterfamily above keeps $U_0$ unchanged and changes the
actual transport. These are different tests. The duality theorem excludes
neither transport family merely by naming its coefficient line. Physical
contact ancestry requires the actual-action certificate, and a later
Lorentzian/source-action claim retains its separate complete-action port.
\end{remark}

\begin{definition}[Role-valued sector trace and normalized reciprocal role scalarizer]
\label{p3b:def:sector-trace-born-state}
Let
\[
 Z_{\rm role}:=\mathbb R e_x\oplus\mathbb R e_p,
 \qquad
 e_xe_p=0,\qquad e_x+e_p=1,
\]
be the two-sector label algebra.  For positive trace-class sector operators,
define
\[
\begin{aligned}
 \mathfrak T_{\rm sec}:&
 \mathcal L^1(\mathcal H_x)_+
 \oplus\mathcal L^1(\mathcal H_p)_+
 \longrightarrow (Z_{\rm role})_+,\\
 \mathfrak T_{\rm sec}(A_x\oplus A_p)
 &:=
 \operatorname{Tr}_{\mathcal H_x}(A_x)e_x
 +\operatorname{Tr}_{\mathcal H_p}(A_p)e_p .
\end{aligned}
\]
The physical scalar-action constructor then uses a positive real-linear
role scalarizer
\[
 \lambda_{\rm B}:Z_{\rm role}\to\mathbb R,
 \qquad
 \lambda_{\rm B}(1)=1.
\]
On the native reciprocal owner the complete consumed presentation of this
scalarization is Born-unmarked and the target is an ordinary scalar.  Hence
Paper~II-A, Theorem~\ref{p2a:thm:generated-born-inheritance}, gives
\[
 \lambda_{\rm B}\circ\sigma=\lambda_{\rm B},
 \qquad
 \sigma(e_x)=e_p,\quad \sigma(e_p)=e_x.
\]
Accordingly the normalized two-sector scalar functional is
\[
 \mathcal S_{\rm B}
 :=\lambda_{\rm B}\circ\mathfrak T_{\rm sec}.
\]
The map \(\mathfrak T_{\rm sec}\) performs ordinary carrier traces first;
\(\lambda_{\rm B}\) acts only on the resulting role labels.  It is a
normalization/scalarization map in the action owner, not a prepared physical
state.  This definition does not identify unequal carrier identities or
normalize their matrix dimensions.  When the same two-dimensional algebra is
used abstractly without the displayed origin certificate, exchange invariance
remains an explicit defining condition rather than an inherited physical
conclusion.
\end{definition}

\begin{lemma}[Inherited Born covariance and unit normalization fix the reciprocal role weights]
\label{p3b:lem:equal-born-role-weights}
Every positive real-linear scalarizer
\(\lambda:Z_{\rm role}\to\mathbb R\) has the form
\[
 \lambda(ae_x+be_p)=w_xa+w_pb,
 \qquad w_x,w_p\ge0.
\]
On the native Born-unmarked scalar-action owner,
Theorem~\ref{p2a:thm:generated-born-inheritance} implies
\(\lambda\circ\sigma=\lambda\), and therefore
\[
 w_x=w_p.
\]
The independent unit normalization \(\lambda(1)=1\) gives
\(w_x+w_p=1\).  Hence
\[
 \boxed{w_x=w_p=\frac12,}
 \qquad
 \boxed{\lambda_{\rm B}(ae_x+be_p)=\frac{a+b}{2}.}
\]
Thus Born inheritance fixes the ratio \(1:1\), while the unit normalization
fixes the common magnitude.  Reciprocity by itself would determine only
\((w_x,w_p)=(c,c)\).
\end{lemma}

\begin{proof}
The two minimal idempotents form a basis, so linearity gives the displayed
form and positivity gives \(w_x,w_p\ge0\).  Because the role scalarizer is a
licensed scalar output of the Born-unmarked native action constructor,
Paper~II-A gives
\[
 \lambda(e_x)=\lambda(\sigma e_x)=\lambda(e_p),
\]
hence \(w_x=w_p\).  Separately,
\[
 1=\lambda(1)=\lambda(e_x+e_p)=w_x+w_p.
\]
The unique solution is \((w_x,w_p)=(1/2,1/2)\).  No state occurrence, matrix
block-dimension normalization, or response conversion coefficient is used.
\end{proof}

\begin{corollary}[All-orders Born-involution selection on a smooth unmarked stratum]
\label{p3b:cor:v1100-born-involution-derivative-selection}
Let $B$ be the inherited Born-exchange involution, acting linearly in a
licensed local configuration chart, and let the \emph{unmarked scalar}
functional $\mathcal A$ obey $\mathcal A\circ B=\mathcal A$. At a
Born-fixed configuration $q_0=Bq_0$, suppose $\mathcal A$ is $C^n$ and
$Bv_i=\epsilon_i v_i$ with $\epsilon_i\in\{+1,-1\}$. Then
\[
 \boxed{\prod_{i=1}^{n}\epsilon_i=-1
 \quad\Longrightarrow\quad
 D^n\mathcal A(q_0)[v_1,\ldots,v_n]=0.}
\]
This statement applies on the constituted smooth Born-unmarked
scalar-action owner, not to signed, marked, oriented, charged,
state-conditioned, or other physical sectors without the stated
invariance. At K16 or another nonsmooth depth threshold, the derivative
order is restricted to the actual smoothness class.
\end{corollary}
\begin{proof}
Differentiating $\mathcal A(Bq)=\mathcal A(q)$ at $Bq_0=q_0$
$n$ times gives
\[
\begin{aligned}
 D^n\mathcal A(q_0)[Bv_1,\ldots,Bv_n]
 &=D^n\mathcal A(q_0)[v_1,\ldots,v_n].
\end{aligned}
\]
Multilinearity multiplies the left side by $\prod_i\epsilon_i$;
if that product is $-1$, the derivative equals its negative and
vanishes.
\end{proof}


\begin{remark}[Persistence of the role scalarizer under unequal finite levels]
\label{p3b:rem:role-scalarizer-unequal-level-persistence}
The coefficients \(1/2,1/2\) belong to the constituted map
\(\lambda_{\rm B}:Z_{\rm role}\to\mathbb R\), not to the dimensions of the
later finite carrier blocks.  If the position role subsequently unfolds to
level \(m_x\) while the compact momentum role remains at \(m_p=1\), the same
action owner evaluates
\[
 \lambda_{\rm B}\!\left(T_x(m_x)e_x+T_p(1)e_p\right)
 =\frac12T_x(m_x)+\frac12T_p(1).
\]
The unequal level marks mean that Born exchange is then reciprocal branch
transport rather than an internal symmetry of the selected finite
presentation.  They do not alter the already constituted role scalarizer.
Any replacement \((1/2,1/2)\mapsto(w_x(m_x,m_p),w_p(m_x,m_p))\) would require
a new weighting constructor or calibration not present in the frozen action
grammar.
\end{remark}


\begin{theorem}[Invariant spectral-composite exhaustion]
\label{p3:thm:invariant-spectral-composite-exhaustion}
For \(D=D^*\in M_n(\mathbb C)\), every unitary-invariant scalar function
factors through the unordered eigenvalue multiset in \(\mathbb R^n/S_n\).
For polynomial functions,
\[
 \boxed{
 \mathbb C[\operatorname{Herm}(n)]^{U(n)}
 =\mathbb C[p_1,\ldots,p_n],
 \qquad p_k(D)=\operatorname{Tr}(D^k).}
\]
Thus nonlinear functions of the power traces exhaust the invariant polynomial
composites; only the additive subclass is a primitive single loop.
\end{theorem}

\begin{proof}
The spectral theorem writes
\(D=U\operatorname{diag}(\lambda_1,\ldots,\lambda_n)U^{-1}\), so an
invariant is a symmetric function of the eigenvalues.  The fundamental
theorem of symmetric polynomials expresses every polynomial invariant in the
elementary symmetric functions \(e_k\).  Newton's identities
\[
 k e_k=\sum_{j=1}^k(-1)^{j-1}e_{k-j}p_j,
 \qquad e_0=1,
\]
recursively express \(e_k\) in \(p_1,\ldots,p_k\); the converse is immediate
because every \(p_k\) is symmetric.  The first \(n\) power traces determine
the characteristic polynomial and hence the orbit.  A nonlinear polynomial
in several \(p_k\)'s correlates loop outputs and is therefore a generated
composite rather than a new outer closure.
\end{proof}

\begin{definition}[Heat-process class and atomic heat-scalar subclass]
\label{p3:def:Dheat}
The occurring heat-process class consists of typed semigroup families
\[
 \mathcal P_{\rm heat}
 :=\left\{\begin{gathered}
 (V_t)_{t\ge0}:\quad V_0=\mathbf1,\\
 V_{s+t}=V_sV_t,\quad V_t=e^{-tD^2},\\
 \text{the family has a G0 process occurrence}
 \end{gathered}\right\}.
\]
Its scalar readout subclass is
\[
 \Dheat:=\left\{
 \operatorname{Tr}_{\mathcal H}V_t:
 (V_s)_{s\ge0}\in\mathcal P_{\rm heat},\ t>0
 \right\}\subset\Dcl.
\]
By Paper~II-D Theorem~\ref{p2d:thm:heat-semigroup-rigidity}, the semigroup
law fixes the exponential.  At the spectral-action scale one takes
\(t=\Lambda^{-2}\).  Thus \(V_t\) is an endomorphism/process and
\(\operatorname{Tr}V_t\) is the associated scalar; neither is typed as the
other.  The existence of mathematical heat flow does not by itself supply
the physical process occurrence.
\end{definition}

\begin{remark}[Relation to the Paper~III-F reduction]
\label{p3:rem:dcl-reduction-crossref}
Paper~III-F may consume Theorem~\ref{p3:thm:gp2-dcl-derivation} after an
actual closed-process occurrence is supplied.  Its common-action or gauge
reduction cannot retroactively identify a boundary matter line with the
gravitational loop, and sharing a process profile requires a separate
same-process bridge.
\end{remark}

\subparagraph{Framework export: \(\Dcl\).}
Definition~\ref{p3:def:D1} exports the derived spectral form class.  A later
statement ``inside \(\Dcl\)'' means that a closed occurrence has supplied one
generated positive Dirac-even process of this form together with the local
analytic hypotheses displayed there.  It never means that GP1--GP3 selected
the profile.  The particular \(F\) belongs to the occurring process datum;
identifying profiles attached to independently rooted processes requires a
separate same-process bridge and is not an inheritance theorem.  This is the
class-allocation statement.  At the subsequent dynamical profile allocation,
Section~\ref{p3b:sec:dynamical-profile-invariance} reapplies GP1--GP3 to the
complete action presentations; Paper~III-E then proves that the
premise-minimal gravitational heat fibre is a singleton.

\begin{remark}[Action profile versus physical resolution frontier]
\label{p3:rem:action-profile-resolution-frontier}
The common \(\Dcl\) cutoff profile weights spectral data in the action.  It
does not by itself objectify the nonzero spectral range as the complete
physical reciprocal-position test frontier.  A compactly supported or sharp
profile is also not selected by \(\Dcl\), which leaves the common profile
free.  The later sharp-frontier occurrence \(\mathsf{SF}_\Lambda\) is
therefore a separate typed operational bridge, not a hidden consequence of
the one-profile action class.
\end{remark}

\subsubsection{Recursive dynamical closure of profile freedom}
\label{p3b:sec:dynamical-profile-invariance}

The profile allocation is a new physical role and therefore receives the same
GP1--GP3 closure pass as every earlier role.  This is not a fourth principle
and does not reopen any kinematic selection.  The complete action presentation,
rather than the symbol \(F\) in isolation, is the object compared.

\begin{definition}[Exact dynamical profile moduli and coefficient fibres]
\label{p3b:def:exact-dynamical-profile-moduli}
Fix the profile-independent prefix already frozen before the dynamical
allocation, and let a profile candidate be the complete typed presentation
\[
 p=(F,L;S_F,\delta S_F,\delta^2S_F,\ldots),
 \qquad
 S_F(D)=\operatorname{Tr}F(D^2/L^2),
\]
containing every action readout and variation generated by the frozen
constructor word.  Only profiles admitted by that word are candidates; the
definition does not promote all positive functions to physical processes.
Let
\[
 \mathfrak P_{\rm dyn}
 :=\operatorname{Sol}^{\rm ex}(\Delta_{\rm dyn})/
   (\cong,\sim_{\rm cal})
\]
be the complete GP1-admissible, jointly GP2-exact solution moduli, quotiented
only by structured isomorphisms and calibrations generated before evaluation.
When the Einstein--Hilbert coefficient compiler is present, define
\[
 c_{\rm EH}([p])
 :=\frac{f_2(F)}6L^2,
 \qquad
 \mathfrak P_\lambda:=c_{\rm EH}^{-1}(\lambda).
\]
The notation \(f_2(F)\) denotes the declared rescaled second profile moment in
the relevant branch.  A licensed readout
\(Y:\mathfrak P_{\rm dyn}\to\mathsf Y\) is
\emph{profile invariant at \(\lambda\)} when its restriction to every
GP3-surviving class in \(\mathfrak P_\lambda\) is constant.
\end{definition}

\begin{theorem}[GP1--GP3 dynamical profile-invariance and factorization]
\label{p3b:thm:gp-profile-invariance-factorization}
Assume the exact-solution constructor, its output sort, the coefficient map
\(c_{\rm EH}\), and the calibration congruence in
Definition~\ref{p3b:def:exact-dynamical-profile-moduli} have frozen
constructor-origin certificates.  Then the dynamical profile role closes as
follows.
\begin{enumerate}[label=(\roman*),leftmargin=2.1em]
\item GP1 applies to the complete realized finite-window action/readout
      family: every constituted scalar-readout image and complete timed-rate
      image is totally bounded in its generated uniform structure.  A smooth
      profile may remain a symbolic or asymptotic parent presentation; GP1
      does not select a function by counting its parameters.
\item GP2 retains the complete action presentation of every candidate.  It
      forbids choosing a preferred \(F\) or \(L\), and it also forbids
      quotienting two profiles merely because one moment agrees when another
      generated action readout or variation distinguishes them.
\item GP3 retains a unique least structured profile class when one exists.
      Otherwise it retains the complete unresolved irreducible family and
      freezes only readouts common to every surviving class.
\item Consequently, among licensed readouts whose owner, target sort, and
      remaining ancestry are already complete, the exact condition for a
      member-independent generated prediction is profile invariance on the
      surviving coefficient fibre.  Equivalently, every such prediction has
      a unique factorization
      \begin{equation}
       \label{p3b:eq:profile-invariant-factorization}
       \boxed{Y(F,L)=\overline Y\!\left(\Lambda_{\rm EH}^{2}\right),
       \qquad
       \Lambda_{\rm EH}^{2}=\frac{f_2(F)}6L^2,}
      \end{equation}
      on the full nonempty generated domain.  If two surviving members of one
      fibre give different values, no value is generated by that profile
      stage: an advertised output remains a mandatory unresolved prediction
      until a unique class or an additional common-readout theorem is derived.
\end{enumerate}
Thus profile invariance is a theorem about all generated predictions, not a
claim that every writable spectral functional is independent of its profile.
\end{theorem}

\begin{proof}
The certified exact-solution word and its quotient put
\(\mathfrak P_{\rm dyn}\) directly in the domain of Paper~II-A,
Theorem~\ref{p2a:thm:generated-exact-solution-moduli}.  GP1 gives the first
statement at the realized readout level; it supplies no profile selector.
The five zero-defect clauses of GP2 give the second statement: no insertion
forbids an after-the-fact representative, while no omission and typed host
faithfulness retain every generated distinction in the action and its
variations.  Part~(ii) of that theorem gives the GP3 alternative in~(iii).

For a licensed readout on a nonempty fibre, Paper~II-A
Theorem~\ref{p2a:thm:generated-exact-solution-moduli}(iii) and
Theorem~\ref{p2a:thm:rooted-generated-parent-cascade} say precisely that a
constant value is common generated content and may feed the next stage,
whereas a member-dependent value may not.  This proves the prediction
criterion.  If \(Y\) is constant on every fibre, define
\(\overline Y(\lambda):=Y(p)\) for any
\(p\in\mathfrak P_\lambda\).  Fibre constancy makes this well defined and
gives \(Y=\overline Y\circ c_{\rm EH}\); surjectivity onto the attained
coefficient image makes the factorization unique.  Conversely, every such
factorization is constant on each fibre.  This proves~(iv).
\end{proof}

\begin{proposition}[Equal Einstein moment does not identify raw profiles]
\label{p3b:prop:profile-invariance-sharpness}
On the optional enlarged positive-profile class, equality of
\((f_2/6)L^2\) alone does not make the raw finite spectral actions equal.
Indeed, the smooth positive normalized profiles
\[
 F_0(u)=e^{-u},\qquad
 F_1(u)=\frac13e^{-u/2}+\frac23e^{-2u}
\]
satisfy
\[
 F_0(0)=F_1(0)=1,\qquad
 \int_0^\infty F_0(u)\,du
 =\int_0^\infty F_1(u)\,du=1,
\]
and hence have the same rescaled second moment in the convention of
Paper~III-E.  Nevertheless \(F_0(1)\ne F_1(1)\).  On a one-dimensional
finite datum with \(D^2/L^2=1\), one therefore has
\[
 \operatorname{Tr}F_0(D^2/L^2)
 \ne
 \operatorname{Tr}F_1(D^2/L^2).
\]
Accordingly GP2 host faithfulness keeps these profiles distinct unless the
frozen dynamical signature has already removed the distinction.  Theorem
\ref{p3b:thm:gp-profile-invariance-factorization} is sharp: it derives the
factorization of predictions and permits a later singleton profile theorem;
it does not impose a false quotient on raw actions.
\end{proposition}

\begin{corollary}[No downstream profile insertion]
\label{p3b:cor:no-downstream-profile-insertion}
After the dynamical profile closure pass, a descendant may consume either the
unique GP3-retained profile class or a readout proved constant on the complete
surviving profile fibre.  Supplying a representative profile, a new moment,
or a cutoff convention at the descendant in order to obtain a number is a
GP2 insertion and cannot be a framework prediction.
\end{corollary}

\begin{proof}
This is the recursive-stage clause of
Theorem~\ref{p2a:thm:rooted-generated-parent-cascade} applied to
Theorem~\ref{p3b:thm:gp-profile-invariance-factorization}.
\end{proof}

\subparagraph{Speed and principal-ancestry boundary.}
Conditional on
\((\mathrm{CAL}_{t})\) and the spatial-radius representative,
\(c_{\rm op}=2\pi R_x/T_t\) is the derived clocked-Hopf operational ceiling,
and Paper~III-E writes the Newton/Einstein--Hilbert calibration in that host
metric rate.  Paper~II-A proves that a later undeformed same-owner neutral
principal descendant inherits the same front speed.  Paper~III-B neither
constructs such a field channel nor supplies an empirical light-channel
identification.

\begin{definition}[Finite-window reading of \texorpdfstring{$\Dcl$}{D1}]
\label{p3:def:D1-finite-window}
Whenever the sector Hilbert space contains an infinite scalar-circle factor, the classification theorem is applied to the finite spectral window
\[
        \mathcal K_{m,N}:=\mathcal K_m\otimes P_{|k|\le N}L^2(\Sone)\otimes\CC^2,
        \qquad
        D_{m,N}:=P_{m,N}DP_{m,N},
\]
with the operator-specific unlabelled datum
\((\mathcal K_{m,N},D_{m,N},\Lcut)\).  The action at finite window is the finite trace classified by Theorem~\ref{p3:thm:spectral-characterization}.  The continuum heat-kernel formulae used later are the standard semiclassical asymptotics of this monotone spectral-window exhaustion, with one common admissible profile.  Thus \(\Dcl\) never uses the finite-dimensional classification theorem as a proof about arbitrary untruncated functionals on the full infinite Hilbert space; it uses the theorem at each finite spectral window and then passes to the usual spectral-action limit.
Fixed-time trace convergence of these windows is not, by itself, convergence
of their strict small-time coefficients.  Continuum heat coefficients belong
to the compact parent and are approximated by the windows only in the joint
large-window/intermediate-time regime stated in
Theorem~\ref{p3:thm:directed-spectral-window-exhaustion}.
\end{definition}

\begin{theorem}[Regularity bookkeeping inside the common-profile class]
\label{p3:thm:no_residual_regularity_D1}
Inside the dynamical class \(\Dcl\), the regularity conditions used below are either consequences of the finite spectral construction or admissibility
conditions of the single spectral-action profile:
\begin{enumerate}[label=(\roman*),leftmargin=2.1em]
\item unitary naturality follows from the unmarked loop closure on the finite
spectral datum \((\mathcal H,D)\);
\item block additivity follows from closure on licensed direct sums and is realized by
the trace on orthogonal direct-sum blocks;
\item evenness follows from the generated positive input, represented through
\(D^2/\Lcut^2\);
\item continuum regularity is part of the common-profile admissibility
required of the occurring process in \(\Dcl\);
\item any grade-resolved, channel-resolved, or sector-dependent profile is a datum-enlarged inserted
extension and requires extra endomorphism data not supplied by \(\mathcal A_{\mathrm{book}}\).
\end{enumerate}
Consequently the later one-parameter reduction uses no regularity hypothesis
beyond the occurrence and analytic content explicitly displayed for \(\Dcl\).
\end{theorem}

\begin{proof}
Unitary naturality, block additivity, and evenness are the conclusions of
Theorem~\ref{p3:thm:gp2-dcl-derivation}; common-profile regularity is an
analytic condition on its occurring process.  On licensed direct sums,
block additivity is realized by the direct-sum identity
\[
        \Tr f((D_1\oplus D_2)^2/\Lcut^2)
        =\Tr f(D_1^2/\Lcut^2)+\Tr f(D_2^2/\Lcut^2),
\]
which is Lemma~\ref{p3:lem:additivity} in spectral form.  The generated
positive process takes \(D^2/\Lcut^2\) as its input.
The reduced valuation inherits continuity from the common admissible profile used in the finite
functional calculus and in the continuum heat-kernel asymptotics.  Finally,
Corollary~\ref{p3:cor:no-channel-profiles} proves that every non-scalar insertion, grade
projector, or independent profile leaves the uninserted class.  Thus the listed properties are not
additional hidden regularity assumptions; they are trace identities, input
typing, or displayed analytic conditions on the occurring process.
\end{proof}


\begin{lemma}[No invariant mixed vector contraction before diagonal identification]
\label{p3:lem:no-invariant-mixed-vector-contraction}
Let \(V_x\) and \(V_p\) denote the vector tangent/first-harmonic representations carried by the reciprocal position and momentum factors before any diagonal identification of their covariance groups.  The unlabelled compact product carries independent covariance
\[
        G=\operatorname{Sp}in(4)_x\times\operatorname{Sp}in(4)_p
\]
up to the Born exchange involution.  Then there is no nonzero canonical bilinear scalar
\[
        B:V_x\otimes V_p\longrightarrow \RR
\]
which is invariant under \(G\).  Equivalently,
\[
        (V_x^*\otimes V_p^*)^{\operatorname{Sp}in(4)_x\times\operatorname{Sp}in(4)_p}=0 .
\]
A mixed contraction exists only after adjoining extra data, such as a diagonal identification of the two \(\operatorname{Sp}in(4)\) factors, a metric-valued intertwiner, a chosen local reciprocal chart, or a sector-resolving projector.
\end{lemma}

\begin{proof}
For a product group acting separately on the two factors,
\[
        (V_x^*\otimes V_p^*)^{\operatorname{Sp}in(4)_x\times\operatorname{Sp}in(4)_p}
        \cong
        (V_x^*)^{\operatorname{Sp}in(4)_x}\otimes (V_p^*)^{\operatorname{Sp}in(4)_p} .
\]
The vector representation of \(\operatorname{Sp}in(4)\) has no nonzero invariant linear functional, so both invariant factors on the right vanish.  Hence no invariant bilinear mixed scalar is carried by the unlabelled product.  The Born exchange involution identifies the two sectors only as reciprocal copies; it does not supply a pointwise diagonal contraction between their vector representations.  Local cotangent expressions such as \(dp_\mu\wedge dx^\mu\) arise only after a reciprocal chart has been chosen and therefore do not define an unlabelled global sector-action term.
\end{proof}

\begin{theorem}[Sectorwise additivity in the loop-derived \(\Dcl\) realization]
\label{p3:thm:sectorwise_additivity}
For the completed compact Born/contact geometry of
Theorem~\ref{p3:thm:phasespace}, no canonical mixed \(x\)--\(p\) contraction
or intertwiner is carried by the unlabelled geometry beyond Born exchange.
Thus no genuinely mixed \(x\)--\(p\) action term is kinematically forced.
The closed-process occurrence on the licensed direct sum gives the
\(\Dcl\) realization, whose action is sectorwise additive.  A larger generated
class may include a mixed term after specifying an intertwiner.

More precisely, the licensed direct-sum constructor supplies the block action
datum.  Writing
\((\mathcal A_x,\mathcal H_x,D_x)\) and
\((\mathcal A_p,\mathcal H_p,D_p)\) for their sector representatives (with
the finite-window representatives displayed below), it is
\begin{equation}
 \mathcal A_{\rm act}=\mathcal A_x\oplus\mathcal A_p,
 \qquad
 \mathcal H_{\rm act}=\mathcal H_x\oplus\mathcal H_p,
 \qquad
 D_{\rm act}=D_x\oplus D_p.
 \label{p3:eq:primitive-reciprocal-action-coproduct}
\end{equation}
The normalized reciprocal role scalarizer on its two-dimensional centre is
\begin{equation}
       \lambda_{\rm B}(ae_x+be_p)=\frac{a+b}{2}.
       \label{p3:eq:born-central-state}
\end{equation}
Its equality of reciprocal weights is inherited from Born covariance of the licensed scalar-action constructor, while the common magnitude is fixed by the independent unit normalization; therefore the two reciprocal sector traces carry coefficients \(1/2\) each.
The tensor-product triple remains the kinematical conjunction; it does not
replace the loop-derived sectorwise action direct sum.
\end{theorem}

\begin{proof}
By Theorem~\ref{p3:thm:position}, the position and momentum sectors are related
only by the involutive Born exchange; no preferred pairing of individual
\(x\)- and \(p\)-degrees of freedom is part of the completed kinematics.
Lemma~\ref{p3:lem:no-invariant-mixed-vector-contraction} shows that, before an
additional diagonal identification or local chart is supplied, the
product reciprocal covariance group carries no mixed invariant scalar
contraction.  Therefore a genuinely mixed \(x\)--\(p\) action term would
require extra contraction or intertwiner data not supplied by the completed
geometry.

Theorem~\ref{p3:thm:gp2-dcl-derivation} represents its derived direct-sum additivity by
the block direct sum
\eqref{p3:eq:primitive-reciprocal-action-coproduct}; functional calculus and
ordinary trace additivity give
\[
 \Tr f(D_{\rm act}^2/\Lcut^2)
 =\Tr f(D_x^2/\Lcut^2)+\Tr f(D_p^2/\Lcut^2).
\]
On the centre \(\mathbb Re_x\oplus\mathbb Re_p\),
Lemma~\ref{p3b:lem:equal-born-role-weights} applies directly to the licensed
role scalarizer: hereditary Born covariance gives equality of the two weights,
and the independent unit normalization gives their sum.  Thus
\eqref{p3:eq:born-central-state} follows with weights \(1/2,1/2\), without a
state occurrence or a separate Einstein--Hilbert matching convention.
\end{proof}

\begin{corollary}[No canonical boundary-to-position transfer channel]
\label{p3:cor:no-boundary-to-position-transfer}
Let \(\Delta_x\) and \(\mathcal L_r^\sigma\) denote the Paper~I-A compact
coordinate-closure defect and its declared finite symmetric-degree boundary
face.  The Born exchange involution does not generate a canonical linear map
\[
  J_r:\mathcal L_r^\sigma\longrightarrow\mathcal H_x
\]
from that excluded compact face to the reciprocal-position role.  Any such
map would pair individual momentum-side boundary directions with individual
position-side degrees of freedom and would therefore be an additional mixed
intertwiner excluded by Theorem~\ref{p3:thm:sectorwise_additivity}.
Consequently, compact coordinate nonclosure may constrain which finite
supports are retained, but it is not a physical flux law and cannot by itself
generate, suppress, or threshold a reciprocal-position source action.
\end{corollary}

\begin{proof}
A nonzero canonical \(J_r\) would be constitutive mixed \(x\)--\(p\) data.
Theorem~\ref{p3:thm:sectorwise_additivity} proves that the completed unlabelled
Born geometry carries no such intertwiner beyond exchange of the two sectors
as wholes.  The Paper~I-A quotient therefore remains a compact-side
representability statement.  Reciprocal-position dynamics is supplied
independently by the realized position polarization and its sectorwise local
spectral action.
\end{proof}

\begin{remark}[Forward source criterion]
The native III-B theorem chain supplies no metric source law.  Conditional on
Paper~III-F's later common-action theorem, Hilbert variation gives a block-
additive reciprocal-position source criterion; the full conditional statement
and proof are recorded in
Paper~III-G, Section~\ref{p3b:app:forward-synthesis}.  Without that later antecedent no
source criterion is asserted here.
\end{remark}


\begin{remark}[Paper~VI-A chain-closed cosmological realization]
\label{p3:rem:p7-chain-closed-cosmology}
\leavevmode\par\noindent
No additional global Born-manifold principle is introduced here for the
existence of the reciprocal role.  Paper~III-A derives the exact local
operation/generator polar target independently of any diagonal identification.
Paper~III-F, Theorem~\ref{p3:thm:local_reciprocal_atlas}, then identifies the
already constituted global reciprocal role with those derived cotangent
generator germs only under the narrower global--local interface
\((\mathrm{GL}_{xp})\).  Thus neither the global Born involution nor the
present independent-covariance data silently supply a pointwise tangent--
cotangent identification.

The only nonlocal realization statement needed beyond Paper~III is the
physical cosmological realization of the semiclassical derivative-expansion
regime on the homogeneous branch.  Paper~VI-A supplies that closure.  Paper~III
then continues with the local consequences on reciprocal-position charts in
that regime.
\end{remark}


\begin{theorem}[Shell-uniform functional calculus and sharp shell removal]
\label{p3:thm:shellwise_coarse_graining}
Let
\[
 \widehat N=\sum_{n\in\mathcal I}nP_n
\]
be the canonical total-weight operator on a finite truncation, where \(P_n\)
projects onto the complete symmetry-equivalent shell \(\mathcal H_n\).  Then
every functional-calculus profile satisfies
\[
 g(\widehat N)=\sum_{n\in\mathcal I}g(n)P_n
\]
and therefore acts by one scalar on each complete shell.  It never splits a
shell.  It literally retains or removes complete shells precisely when
\(g(n)\in\{0,1\}\) on the retained spectrum, equivalently when
\(g(\widehat N)\) is a spectral projector there.  Thus shell uniformity is
imported kinematic structure; sharp shell removal is the projector-valued
specialization, not a property of every cutoff profile.
\end{theorem}

\begin{proof}
Paper~II supplies the displayed spectral resolution.  Finite functional
calculus gives the formula for \(g(\widehat N)\), so its restriction to
\(\mathcal H_n\) is \(g(n)I_{\mathcal H_n}\).  This proves shell uniformity.
Such an operator is an orthogonal projector exactly when
\(g(n)^2=g(n)=\overline{g(n)}\) for every spectral value, which is equivalent
to \(g(n)\in\{0,1\}\).  Only in that case is each shell kept or deleted.
\end{proof}

\begin{remark}[Optional full-cell scalar profile]
\label{p3:rem:cubature-f4-use}
On an actualized core-support role \(\mathsf{Supp}_{2T}\),
degree six is the first available nonconstant fully \(24\)-cell-symmetric
scalar profile after spherical normalization.  This is an availability
statement on the core-support role, not an activated term of the native
III-B action.  The calculation and its normalized Hurwitz cubature defect are
deferred to Paper~III-G, Section~\ref{p3b:app:finite-support-scalar}.
\end{remark}


\begin{remark}[Paper~IV closure of the quantum coarse-graining interpretation]
\label{p3:rem:p4-coarse-graining-closure}
\leavevmode\par\noindent
The theorem above is a Paper~III geometric/kinematic statement about complete
symmetry shells.  When the same shell removal is read as a quantum channel, RG
bookkeeping operation, or threshold-matching map, the required finite quantum
statement is deferred to Paper~IV.

There, the Paper~IV finite coarse-graining semigroup theorem proves that
compatible finite coarse-graining maps form an exact CPTP semigroup and supply
the shell-flow law.  Thus Paper~III derives the shellwise reduced dynamics,
while Paper~IV closes the finite quantum/RG interpretation of the same
bookkeeping.
\end{remark}


\begin{remark}[Status of the one-parameter reduction]
\label{p3:rem:D4-status}
\leavevmode\par\noindent
The one-parameter shell-decoupling reduction is derived later in
Sec.~\ref{p3:sec:reduced_reciprocity}.  Within the fully invariant
round-sector spectral-action framework, shellwise coarse-graining reduces the
effective action to a single running parameter.

The reduced coupling is then fixed by the visible reciprocal rigidity theorem,
Theorem~\ref{p3:thm:geff_from_reciprocity}.
\end{remark}


\begin{remark}[Supplementary free Casimir trajectories]
The exact relation between monotone Casimir Hamiltonians, unparameterized
round geodesics, and the independent GP1 timing calibration is proved in the
appendices.  It is not used in the native spectral-action classification.
\end{remark}


\subsection{GP1--GP3 generated profile law}
\label{p3b:sec:v164-generated-profile-law}

\begin{theorem}[Premise-minimal scalar heat-profile specialization]
\label{p3b:thm:v164-generated-heat-profile}
On the \emph{real scalar-function} signature, let the occurring metric root
generate an unmarked, positive, unit-preserving, strongly continuous, local
formally self-adjoint second-order continuation whose principal symbol is
\[
 \sigma_2(L)(x,\xi)=a\,g^{\mu\nu}(x)\xi_\mu\xi_\nu,
 \qquad a>0,
\]
with no auxiliary line/bundle connection or additional endomorphism datum.
Then reapplying GP1--GP3 gives the single positive scalar generator ray
\[
 L_a=a\Delta_+^{\rm scal},\qquad a>0,
\]
and hence
\[
 H_t^{(a)}=e^{-ta\Delta_+^{\rm scal}},
 \qquad F_a(u)=e^{-au}.
\]
Metric naturality and the stated scalar principal symbol permit only
\(a\Delta_+^{\rm scal}+E\); preservation of the unit forces \(E=0\), exact
composition exponentiates the generator, and GP3 retains the unique
projective scalar generator class.  This theorem does not classify arbitrary
spin, gauge, bundle or joint Dirac-square owners.  Nonlocal subordinations,
higher-order terms, potentials, auxiliary connections, channel-dependent
profiles, and mixture measures are enlarged parents rather than competing
members of this scalar branch.
\end{theorem}

\begin{corollary}[Exact rate-scale orbit]
\label{p3b:cor:v164-profile-scale-orbit}
At spectral scale \(t=\Lambda^{-2}\),
\[
 A_{a,\Lambda}(L)=e^{-aL/\Lambda^2}
 =e^{-L/\Lambda_{\rm eff}^2},
 \qquad
 \Lambda_{\rm eff}^2:=\frac{\Lambda^2}{a}.
\]
For every \(\rho>0\),
\[
 A_{\rho a,\sqrt\rho\Lambda}(L)=A_{a,\Lambda}(L).
\]
Thus the premise-minimal profile class is one heat profile together with one
effective scale.  The general profile moduli proved earlier in this paper
remain the correct enlarged-parent classification; the present theorem
specializes the generated scalar heat branch and does not identify arbitrary
raw cutoff functions.
\end{corollary}


\subsection{generated native heat-process role}
\label{p3b:sec:v173-joint-heat-process}

\begin{theorem}[Heat-process co-realization from a prescribed generated operator square]
\label{p3b:thm:v173-joint-heat-process-corealization}
Let \(D_\Sigma\) be a jointly occurring/generated root tuple whose certified
constructor word produces a complete operator family \(\mathfrak D_\Sigma\)
in a predeclared self-adjoint operator sort.  Restrict to the frozen subdomain
on which each \(D\in\mathfrak D_\Sigma\) has a self-adjoint realization and
the generated continuation generator is its nonnegative square
\[
 A_D:=D^2\ge0.
\]
Assume either that the carrier heat operator is trace class or that an
explicitly declared positive localizer gives a trace-class localized heat
operator.  Then spectral functional calculus uniquely gives the strongly
continuous semigroup
\[
 \boxed{W_t(D)=e^{-tD^2}},\qquad t\ge0,
\]
with
\[
 \begin{aligned}
 W_0(D)&=1,\qquad W_{s+t}(D)=W_s(D)W_t(D),\\
 -\left.\frac{d}{dt}W_t(D)x\right|_{t=0+}&=D^2x
 \qquad(x\in\operatorname{Dom}(D^2)),
 \end{aligned}
\]
Typed unitary intertwiners transport the complete process naturally.  At
\(t=\Lambda_{\rm eff}^{-2}\),
\[
 W_{\Lambda_{\rm eff}^{-2}}(D)
 =e^{-D^2/\Lambda_{\rm eff}^{2}}.
\]
Its unmarked scalar readout is, according to the declared trace domain,
\[
 S_{\rm heat}(D)=\operatorname{Tr}e^{-D^2/\Lambda_{\rm eff}^{2}}
\]
or, for a positive localizer,
\[
 S_{\rm heat}^{U}(D)=
 \operatorname{Tr}\!\left(
 \varphi_U^{1/2}e^{-D^2/\Lambda_{\rm eff}^{2}}\varphi_U^{1/2}
 \right),
\]
whenever this operator is trace class.  Under localization covariance the
pair \((D,\varphi_U)\), not merely \(D\), is transported.

The family-valued process graph and the declared readout are generated
functional-calculus descendants of the operator parent.  Hence the rooted
generated-image theorem co-realizes them without a second structural
closed-loop occurrence.  Denote this result by
\[
 \boxed{(\mathrm{CLOOP}_{\rm heat}^{\rm der}).}
\]
The scalar-selection Theorem~\ref{p3b:thm:v164-generated-heat-profile} is an
independent specialization that explains why a premise-minimal scalar metric
continuation has no competing generator ray; it is not used to classify the
general operator family here.  The separate \((\mathrm{CLOOP}_F)\) premise
remains only for an enlarged profile not generated by the prescribed square
or for a claim that one process/readout token is actually performed.  This
theorem does not identify the spectral cutoff with elapsed Euclidean time,
adopt the Euclidean trace as a Lorentzian action, choose an operator-family
fibre, or fix the absolute scale.
\end{theorem}

\begin{proof}
The spectral theorem for the nonnegative self-adjoint operator \(D^2\) gives
the unique semigroup with generator \(-D^2\) and the displayed composition
law.  Functional calculus gives
\(UDU^{-1}=D'\Rightarrow Ue^{-tD^2}U^{-1}=e^{-tD'^2}\).  Trace or localized
trace is applied only after its domain has been declared.  The rooted
co-realization theorem then transports the occurrence witness from the
generated operator parent to the complete process/readout graph.  No scalar
Markov or unit-preserving classification is needed for this general
prescribed-generator statement.
\end{proof}

\begin{proposition}[Exact topology of the heat-generator readout]
\label{p3b:prop:v284-heat-generator-domain}
Let $A=A^*\geq0$ be a (possibly unbounded) operator on a Hilbert space
and $W_t=e^{-tA}$ for $t\geq0$.  Then $W_t$ is a strongly continuous
contraction semigroup, and its right derivative at zero satisfies
\[
 \lim_{t\downarrow0}\frac{W_tx-x}{t}=-Ax
 \qquad(x\in\operatorname{Dom}A).
\]
Moreover, $t\mapsto W_t$ is right differentiable at zero in operator norm
if and only if $A$ is bounded.  For an unbounded positive generator,
\[
 \left\|\frac{I-W_t}{t}\right\|=\frac1t\qquad(t>0),
\]
so no operator-norm derivative at zero exists.  Consequently recovering
$D^2$ from the general heat process in
Theorem~\ref{p3b:thm:v173-joint-heat-process-corealization}
is a strong-generator operation on $\operatorname{Dom}(D^2)$; it does not
construct a globally bounded instantaneous derivative or an analytic
continuation beyond the declared operator domain.
\end{proposition}

\begin{proof}
The spectral theorem yields $W_t=\int e^{-t\lambda}\,dE_A(\lambda)$.
The bound $|e^{-t\lambda}|\leq1$ and dominated convergence give strong
continuity.  On $\operatorname{Dom}A$,
$|(e^{-t\lambda}-1)/t|\leq\lambda$ and
$(e^{-t\lambda}-1)/t\to-\lambda$; the square-integrable spectral
majorant $\lambda^2$ proves the strong generator identity.  For bounded
$A$, the exponential series converges in operator norm and yields the
norm derivative.  Conversely, if $A$ is unbounded, then its positive
spectrum is unbounded and functional calculus gives, for every $t>0$,
\[
 \|I-e^{-tA}\|=\sup_{\lambda\in\sigma(A)}(1-e^{-t\lambda})=1.
\]
Dividing by $t$ proves the displayed norm divergence and excludes
norm differentiability.  No additional physical occurrence is inferred.
\end{proof}

\begin{corollary}[Application to the generated joint operator family]
\label{p3b:cor:v173-joint-heat-process-application}
The complete family \(\mathfrak D_J(U;F)\) of
Theorem~\ref{p3f:thm:generated-fibre-spin-joint-operator} satisfies the
constructor-origin hypothesis of
Theorem~\ref{p3b:thm:v173-joint-heat-process-corealization}.  Hence, on its
self-adjoint trace/localizer domain, its native heat-process graph and closed
scalar functional are co-realized descendants and require no independent
\((\mathrm{CLOOP}_F)\) role occurrence.
\end{corollary}

\begin{remark}[Role closure versus a performed process]
\label{p3b:rem:v173-heat-role-token-firewall}
The theorem generates the heat-semigroup role and its closed functional over
the admissible operator family.  It does not assert that every semigroup fibre
is physically executed, measured, or interpreted as an actual history.  Such
a token-level claim remains G0-instantiational and adds no new structural
profile on the native branch.
\end{remark}


\section{Spectral triples, heat action, and complete coefficients}
\label{paper3b:sec:organized-spectral}
\subsection{The Spectral Triple Formulation}
\label{p3:sec:spectraltriple}

Table~\ref{tab:p3-spectral-octet-ledger} records the spectral-triple and active-mode data used by the later hierarchy and gravitational modules.

\begin{table}[H]
\centering
\small
\begingroup
\setlength{\tabcolsep}{3pt}
\renewcommand{\arraystretch}{1.15}
\begin{tabular}{@{}
  L{3cm}
  L{9cm}
  L{\dimexpr\linewidth-12cm-4\tabcolsep\relax}
@{}}
\toprule
\textbf{Ledger row}
& \textbf{Claim}
& \textbf{Discharged by} \\
\midrule
Finite spectral windows
& Each auxiliary finite module supplies its block and presentation algebras; the closed-process occurrence uses the latter as its spectral envelope.
& Definition~\ref{p3:def:block-algebra-envelope} \\

Position and momentum triples
& The position and momentum sectors admit dimensionless auxiliary phase-product representatives; no scalar-circle metric is inferred.
& Definitions~\ref{p3:def:pos_triple_4d} and~\ref{p3:def:mom_triple_4d} \\

8-dimensional product triple
& The doubled reciprocal role has an auxiliary normalized product representative.
& Definition~\ref{p3:def:kin_8d_triple} \\

Reduced sector action
& The reduced Born-symmetric sector action follows from the loop-derived \(\Dcl\) form and Born transport.
& Corollary~\ref{p3:thm:sectorwise-from-D} \\

Phase-space algebra
& The direct-sum phase-space algebra is the reduced finite comparison algebra.
& Definition~\ref{p3:def:phase-space-algebra} \\

No bare round asymmetry
& The round \(S^3\) Dirac spectrum carries no bare eta asymmetry.
& Corollary~\ref{p3:cor:round-dirac-eta-zero} \\

Protected octet
& On \((\mathrm{ND}_{\rm sp})\), Paper~II-C supplies the protected even Clifford algebra \(A_{\rm ev}\cong\Cl(0,3)\) independently of the nonzero completion scale and relative chiral phase; the lowest Dirac projector is an independent analytic window.
& Definition~\ref{p3:def:lowest_dirac_projector}, Theorem~\ref{p3:thm:8mode_operator_basis} \\

Scalar projection
& The regular trace extracts the scalar coefficient; the factor \(1/8\)
belongs only to the explicitly adopted rank-one scalar-corner readout.
& Lemma~\ref{p3:lem:clifford_scalar_projection},
Theorem~\ref{p3:thm:normalized-protected-octet-trace},
Definition~\ref{p3:def:scalar-corner-readout} \\

\bottomrule
\end{tabular}
\caption[Spectral octet ledger]{Spectral octet ledger.}
\label{tab:p3-spectral-octet-ledger}
\TableNotes{The following subsections discharge the finite comparison-triple, octet, and scalar-projection ingredients that the hierarchy section uses as protected data.}
\endgroup
\end{table}


\subsubsection{Spectral Triples}

\begin{definition}[Spectral Triple]
A spectral triple $(\mathcal{A}, \mathcal{H}, D)$ consists of:
\begin{itemize}
    \item a represented algebra $\mathcal{A}$,
    \item its representation Hilbert space $\mathcal{H}$,
    \item a Dirac operator $D$ encoding the spectral geometry,
    \item and, on an even-dimensional spectral representation, a grading \(\gamma\) with
      \(\gamma^2=1\), \(\gamma^\dagger=\gamma\) and \(\gamma D=-D\gamma\),
\end{itemize}
satisfying appropriate compatibility conditions \cite{Connes1996,GraciaBondia2001}.
This is presentation-level spectral-triple terminology.  It assigns neither
quantum-event semantics to the elements of \(\mathcal A\) nor quantum-state
semantics to vectors or positive functionals; those interpretations belong to
Paper~IV.
\end{definition}

\begin{remark}[Role of the grading]
\label{p3:rem:grading-role}
The constitutional smooth product
\(\mathcal M_8=S^3\times S^3\times\Sone\times\Sone\) has topological dimension
\(3+3+1+1=8\), but III-A does not constitute independent physical metrics on
the scalar circles.  Therefore this dimension count alone does not supply a
physical eight-dimensional geometric Dirac operator.  On the \emph{auxiliary
normalized metricized comparison representative} of
Definitions~\ref{p3:def:pos_triple_4d}, \ref{p3:def:mom_triple_4d}, and
\ref{p3:def:kin_8d_triple}, the geometric spin Dirac operator is even-dimensional,
carries a chirality grading \(\gamma\), and obeys \(\gamma D=-D\gamma\).
For that graded odd comparison representative,
\(-D=\gamma D\gamma^{-1}\), so unitary naturality implies Dirac-reversal
evenness.  This argument does not apply to an arbitrary self-adjoint
finite-window kinetic representative: for example, \(D=I\) is not unitarily
conjugate to \(-I\).  Accordingly evenness is supplied on the general window
class by the generated positive-process input \(D^2\), not by loop topology;
it becomes redundant only
on its explicitly graded odd comparison subdomain.
\end{remark}


\subsubsection{The Phase Space Spectral Triple}

For auxiliary finite-window comparisons of the constituted reciprocal geometry, the position factor may carry a level \(m_x\) while the physical momentum seed remains fixed at \(m_p=1\).  Thus the comparison notation below is based on \(\mathsf K(m_x,1)\), together with normalized scalar phase circles.  Paper~III-A already supplies the positive universal-cover lift \(\widetilde t\in\mathbb R\) and its generated clocked-Hopf doctrine; a particular causal evolution still requires its sector Hamiltonian and actual trajectory.

\begin{definition}[Finite block algebra and spectral envelope]
\label{p3:def:block-algebra-envelope}
At level \(m\), the Paper~I-A block algebra is
\[
\mathcal A_m^{\rm blk}
:=\bigoplus_{j_L+j_R\le m}\End\bigl(V_{j_L}\otimes V_{j_R}\bigr),
\]
acting on
\[
\mathcal H_m=\bigoplus_{j_L+j_R\le m}V_{j_L}\otimes V_{j_R}.
\]
For any auxiliary finite module \(\mathcal H_m\), its full
endomorphism algebra exists without any action-class or semantic hypothesis:
\[
\mathcal A_m^{\rm env}:=\End_{\CC}(\mathcal H_m)\cong M_{N_m}(\CC),
\qquad N_m:=\dim_{\CC}\mathcal H_m .
\]
The algebra \(\End_{\CC}(\mathcal H_m)\) is canonical as the algebra of all
complex-linear endomorphisms of that finite module; an isomorphism with
\(M_{N_m}(\CC)\) requires a choice of basis and is not itself extra physical
structure.  Thus the existence of \(\mathcal A_m^{\rm env}\) is an
unconditional algebraic/presentation consequence of \(\mathcal H_m\), not a
consequence of \(\Dcl\).

The block algebra is the canonical equivariant auxiliary truncation inherited from
Paper~I-A and retains its displayed block decomposition.  The envelope is the
full matrix presentation algebra on the same underlying module, including
off-block endomorphisms.  After an unmarked closed occurrence has placed the
action in \(\Dcl\), its derived unitary-conjugacy doctrine selects the
\emph{use} of this already existing algebra as the spectral-action envelope on
which \((\mathcal H_m,D,\Lcut)\) and its presentation changes are expressed.
It does not derive the algebra's existence and does not assert that every
element is a physical yes/no event, effect, or quantum observable.  Paper~IV
owns that quantum-event interpretation.  The sector representatives below use
\(\mathcal A_m^{\rm env}\) in the \(\Dcl\) spectral-action sense; replacing it
by the block-preserving subalgebra records the same Paper-I-A truncation but with
fewer admissible presentation unitaries.

At fixed finite \(m\), this is a finite-matrix approximation to the \(S^3\) sector rather than a literal continuum spectral triple.  Its strict finite-matrix spectral dimension is zero on that finite factor.  When a spinorial Dirac spectral window is used, we write \(\mathcal S_m\) conceptually for that finite spinor window; \(\mathcal H_m\) remains the Paper-I-A/II auxiliary \(\operatorname{Sp}in(4)\)-equivariant module and should not be identified dimension-for-dimension with the Dirac eigenspace window.
For statements that do not depend on which kinetic representative is used, we write \(\mathcal K_m\) for the finite sector Hilbert window.  Thus \(\mathcal K_m=\mathcal S_m\) for spinorial Dirac windows, \(\mathcal K_m=\mathcal H_m\) for full-module kinetic representatives, and \(\mathcal K_m=\mathcal H_m^{\rm diag}\) for scalar diagonal Laplacian windows.
The corresponding operator-specific presentation algebra is
\[
 \mathcal A_m^{\rm win}:=\End_{\CC}(\mathcal K_m).
\]
It equals \(\mathcal A_m^{\rm env}\) exactly on the full-module branch
\(\mathcal K_m=\mathcal H_m\); it is a different matrix algebra on a
spinorial or scalar-diagonal window.

\end{definition}
\begin{remark}[Neutral notation for finite operator windows]
When no particular operator choice is being emphasized, one may write \((\mathcal K_m,D_m^{\rm win})\) for the finite sector Hilbert window appropriate to the operator being used: \((\mathcal S_m,D_m^{\rm spin})\) for a spinorial Dirac window, \((\mathcal H_m,K_m^{\rm full})\) for the full auxiliary-module kinetic representative, or the scalar diagonal window for scalar heat traces.  This prevents dimension-counting the Paper-II auxiliary module \(\mathcal H_m\) as though it were always a Dirac eigenspace window.
\end{remark}


\begin{remark}[Finite-window spectral representatives, not continuum uniqueness claims]
All spectral triples and product triples in this section are auxiliary finite-window comparison representatives of the compact reciprocal geometry.  They are the comparison class needed for \(\Dcl\), heat-kernel matching, and sectorwise Born symmetry.  They are not being used as standalone analytic continuum spectral triples with a claimed unique completion outside the finite-window/semiclassical regime.  Whenever continuum language is used below, it refers either to the homogeneous parent geometry or to the directed large-window comparison needed for heat-kernel extraction.
\end{remark}

\begin{definition}[Position normalized phase-product comparison representative]
\label{p3:def:pos_triple_4d}
Let \(D_{\Sone_t}:=-i\partial_{\theta_t}\) be the normalized dimensionless
\(U(1)\) generator on \(L^2(S^1_t)\), and let
\(\widetilde D^{\rm win}_{x,m_x}\) denote the chosen dimensionless normalized
finite position window.  The auxiliary even product representative is
\begin{align}
\mathcal A_x^{\rm cmp}&=\mathcal A_{x,m_x}^{\rm win}\otimes C^\infty(S^1_t),\\
\mathcal H_x^{\rm cmp}&=\mathcal K_{x,m_x}\otimes L^2(S^1_t)\otimes\mathbb C^2,\\
\widetilde D_x^{\rm cmp}&=\widetilde D^{\rm win}_{x,m_x}\otimes1\otimes\sigma_1
 +1\otimes D_{\Sone_t}\otimes\sigma_2,\\
\Gamma_x&=1\otimes1\otimes\sigma_3.
\end{align}
No physical scalar-circle circumference is assigned here.  A dimensionful
\(S^3_x\times S^1_t\) metric Dirac operator would require additional scalar
spectral-metric data not constituted in III-A and is not part of the present
core.
\end{definition}

\begin{definition}[Momentum normalized phase-product comparison representative]
\label{p3:def:mom_triple_4d}
Let \(D_{\Sone_E}:=-i\partial_{\theta_E}\) be the normalized dimensionless
\(U(1)\) generator on \(L^2(S^1_E)\), and let
\(\widetilde D^{\rm win}_{p,1}\) denote the chosen dimensionless normalized
finite momentum window on the physical seed.  The auxiliary even product
representative is
\begin{align}
\mathcal A_p^{\rm cmp}&=\mathcal A_{p,1}^{\rm win}\otimes C^\infty(S^1_E),\\
\mathcal H_p^{\rm cmp}&=\mathcal K_{p,1}\otimes L^2(S^1_E)\otimes\mathbb C^2,\\
\widetilde D_p^{\rm cmp}&=\widetilde D^{\rm win}_{p,1}\otimes1\otimes\sigma_1
 +1\otimes D_{\Sone_E}\otimes\sigma_2,\\
\Gamma_p&=1\otimes1\otimes\sigma_3.
\end{align}
Again this is a normalized comparison representative, not a claim that the
scalar phase circle carries an independently derived physical metric.
\end{definition}

\begin{remark}[Finite-window operators on auxiliary \(S^3\) modules]
\leavevmode\par\noindent
The symbol \(D^{\rm win}_{m}\) denotes the finite-window operator chosen
for the sector currently under discussion.  Three windowed operators must be
kept distinct:
\[
\begin{aligned}
D^{\rm spin}_{m}
&\quad\text{on the spinorial Dirac window } \mathcal S_m,\\
K^{\rm full}_{m}
&\quad\text{on the full } \operatorname{Sp}in(4)\text{-module } \mathcal H_m,\\
\Delta^{\rm diag}_{m}
&\quad\text{on the scalar diagonal Laplacian window}.
\end{aligned}
\]
When a genuine Dirac spectral action is intended,
\(D^{\rm win}_{m}=D^{\rm spin}_m\), and the Hilbert space is the spinorial
window.  When the full auxiliary module \(\mathcal H_m\) is used,
\(D^{\rm win}_{m}\) means the module kinetic/Casimir representative
\(K^{\rm full}_m\), not an untruncated continuum Dirac operator.  Scalar
diagonal heat traces use \(\Delta^{\rm diag}_m\).
\end{remark}

\begin{definition}[Auxiliary normalized reciprocal product representative]
\label{p3:def:kin_8d_triple}
The normalized phase-product comparison representatives may be combined as
\[
 (\mathcal A^{\rm cmp},\mathcal H^{\rm cmp},\widetilde D_{\rm kin},\Gamma_{\rm kin})
 :=(\mathcal A_x^{\rm cmp},\mathcal H_x^{\rm cmp},\widetilde D_x^{\rm cmp},\Gamma_x)
 \times
 (\mathcal A_p^{\rm cmp},\mathcal H_p^{\rm cmp},\widetilde D_p^{\rm cmp},\Gamma_p).
\]
This is an auxiliary normalized comparison object.  It does not supply the
physical \(\Dcl\) action and does not insert a scalar-circle metric.  The
physical action below is instead the Born-symmetric sum of exact finite
sector traces on the explicitly declared sector operators.
\end{definition}

\begin{remark}[Finite-window representatives as the relevant comparison class]
At fixed truncation level the use of a finite-dimensional Hilbert space is not ad hoc: fuzzy spaces are naturally modeled by finite or matrix spectral triples. The framework-specific content of Paper~III is therefore not the bare use of a finite spectral triple, but the particular auxiliary $S^3$ finite-module family constructed in Papers~I-A--II and its reciprocal finite-window comparison representative \cite{Barrett2015,LaiTeh2012}. For a fixed finite matrix window the strict \(t\to0\) spectral dimension is zero on the finite factor; continuum dimensions are recovered only in the appropriate large-window or intermediate-scale regime before the finite cutoff is resolved.  Thus every later heat-kernel or Einstein--Hilbert comparison is a finite-window/semiclassical extraction under the stated normalization, not an assertion that the fixed matrix block itself has continuum spectral dimension.
\end{remark}


\begin{corollary}[Normalized Born action on the constituted reciprocal sectors]
\label{p3:thm:sectorwise-from-D}
Inside \(\Dcl\), let \(D_x\) denote the finite position-sector action
operator on the currently constituted position window and let \(D_p\) denote
the exact finite action operator on the fixed physical momentum seed
\(\mathcal H_p\cong\mathcal H_1\).  The normalized reciprocal role scalarizer gives
\[
 S_{\rm phys}[\Lcut]
 =\frac12\Tr f(D_x^2/\Lcut^2)
 +\frac12\Tr_{\mathcal H_1} f(D_p^2/\Lcut^2).
\]
The auxiliary normalized phase-product representatives above are comparison
objects and are not required to define this action.
\end{corollary}

\begin{definition}[Physical reduced spectral action]
\label{p3:def:physical_spectral_action}
The displayed Born-symmetric exact finite trace is the physical action inside
\(\Dcl\):
\[
 \boxed{S_{\rm phys}[\Lcut]
 =\frac12\Tr f(D_x^2/\Lcut^2)
 +\frac12\Tr_{\mathcal H_1} f(D_p^2/\Lcut^2).}
\]
The factor \(1/2\) is the normalized reciprocal role-scalarizer weight: Born
inheritance fixes the ratio \(1:1\), while unit normalization fixes its common
magnitude.  It is not an adjustable coupling or a prepared physical state.  No
full-product spectral action and no independently metricized
scalar circle is assumed.
\end{definition}

\begin{proposition}[Exact finite Born-heat Hessian before terminal scalarization]
\label{p3b:prop:v245-exact-finite-heat-two-jet}
Let $q\mapsto D_x(q)$ be a $C^2$ family of self-adjoint matrices
on a fixed finite position-action window, and keep $\Lambda_g>0$
and the momentum seed $D_p$ fixed under the physical configuration
variations $u,v$. Write $X(q)=D_x(q)^2/\Lambda_g^2$ and
$S_{\rm B}(q)=\frac12\Tr e^{-X(q)}+\frac12\Tr_{\mathcal H_1}
 e^{-D_p^2/\Lambda_g^2}$. Then the exact first and mixed second
configuration jets are
\begin{align}
 DS_{\rm B}[u]&=-\frac12\Tr(e^{-X}X_u),\\
 D^2S_{\rm B}[u,v]
 &=\frac12\int_0^1\Tr\!\left(
  e^{-(1-s)X}X_v e^{-sX}X_u\right)ds
  -\frac12\Tr(e^{-X}X_{uv}).
 \label{p3b:eq:v245-finite-heat-two-jet}
\end{align}
The momentum trace has zero first and second \emph{configuration} jets
on this frozen-seed, fixed-cutoff domain.  The formula keeps the
noncommuting first-jet factors and the complete contact $X_{uv}$;
it must not be replaced by $\Tr(e^{-X}X_uX_v)$ without a commuting
hypothesis.  A varying cutoff, moving carrier, source density, or
boundary map must be differentiated in the \emph{complete} occurring
word and is not covered by this fixed-window statement.
\end{proposition}
\begin{proof}
Use Duhamel's formula
$D(e^{-X})[v]=-\int_0^1 e^{-(1-s)X}X_v e^{-sX}ds$.
Cyclicity of the finite trace gives the first derivative; differentiating
it again gives the integral and $-\Tr(e^{-X}X_{uv})$.  Both derivatives
of the fixed momentum term vanish. No commutativity is invoked.
\end{proof}

\begin{theorem}[Positive finite-heat Gram and exact stationary incidence discriminant]
\label{p3b:thm:v286-finite-heat-incidence-gram}
Let $E$ be a nonzero finite-dimensional Hermitian carrier, with
$A=A^*>0$ and $B,C_1,\ldots,C_m$ self-adjoint on $E$.  The intended
protected-grade specialization is $E=\operatorname{ran}P_4$;
identifying this mathematical sector with a block of the \emph{actual}
physical action is a separate incidence hypothesis. Define on the real
space of self-adjoint matrices the bilinear form
\begin{equation}
\label{p3b:eq:v286-heat-gram}
 \mathcal G_A(V,W):=\frac12\operatorname{Re}\int_0^1
 \operatorname{Tr}\!\left(e^{-(1-s)A}V e^{-sA}W\right)\,ds.
\end{equation}
Then $\mathcal G_A$ is a positive definite real inner product, and
\begin{equation}
\label{p3b:eq:v286-gram-comparison}
 \frac{e^{-\lambda_{\max}(A)}}2\|V\|_{\rm HS}^2
 \le\mathcal G_A(V,V)\le
 \frac{e^{-\lambda_{\min}(A)}}2\|V\|_{\rm HS}^2.
\end{equation}
For the local affine operator-square chart
$X(u,z)=A+uB+\sum_{i=1}^m z_iC_i$ (positive for sufficiently small
$(u,z)$), consider the finite Born-half-weight heat action
\[
 S(u,z)=\tfrac12\operatorname{Tr}e^{-X(u,z)}+L(u,z),
\]
where $L$ is an already specified affine term, and suppose its actual
first jet satisfies $S_z(0,0)=0$.  Assume that the $C_i$ are linearly
independent on the physical interior quotient. Then the local stationary
branch $z_*(u)$ exists and its exact reduced Hessian is
\begin{align}
 \chi_{\rm red}&:=\left.\frac{d^2}{du^2}
 S(u,z_*(u))\right|_{u=0}
 =\mathcal G_A(B,B)-g^{\mathsf T}M^{-1}g
 \nonumber\\
 &=\min_{c\in\mathbb R^m}\mathcal G_A
 \Bigl(B+\sum_i c_iC_i,B+\sum_i c_iC_i\Bigr),
 \label{p3b:eq:v286-gram-schur}\\
 M_{ij}&=\mathcal G_A(C_i,C_j),\qquad
 g_i=\mathcal G_A(B,C_i).\nonumber
\end{align}
With the empty-interior convention $\chi_{\rm red}=\mathcal G_A(B,B)$,
for $m\ge1$ the equivalent basis-independent scalar is
\begin{equation}
 \label{p3b:eq:v286-gram-determinant}
 \boxed{\chi_{\rm red}=\frac{\det\operatorname{Gram}_{\mathcal G_A}
 (B,C_1,\ldots,C_m)}{\det\operatorname{Gram}_{\mathcal G_A}
 (C_1,\ldots,C_m)}\ge0.}
\end{equation}
It is strictly positive exactly when
$B\notin\operatorname{span}_{\mathbb R}\{C_1,\ldots,C_m\}$, and obeys
\[
 \chi_{\rm red}\ge\tfrac12e^{-\lambda_{\max}(A)}
 \operatorname{dist}_{\rm HS}
 \bigl(B,\operatorname{span}_{\mathbb R}\{C_i\}\bigr)^2.
\]
The criterion is invariant under a change of real basis in the
interior span and simultaneous unitary transport of all finite data.
\end{theorem}
\begin{proof}
Diagonalize $A$ with eigenvalues $\lambda_i$. For Hermitian $V$,
\[
 \mathcal G_A(V,V)=\frac12\sum_{i,j}|V_{ij}|^2
 \int_0^1e^{-(1-s)\lambda_i-s\lambda_j}\,ds,
\]
so each weight is strictly positive and bounded between the two
extremal exponential factors. This proves the inner-product and norm
bounds.  Proposition~\ref{p3b:prop:v245-exact-finite-heat-two-jet}
with $X_{uv}=0$ makes the full finite Hessian of $S$ the Gram form of
$uB+\sum_i z_iC_i$; its affine part contributes no Hessian.
Linear independence gives $M>0$. The implicit function theorem and
$S_z(0,0)=0$ yield a stationary interior germ; differentiating that
stationarity gives $z_*'(0)=-M^{-1}g$. Substitution produces the Schur
formula. Completing the $\mathcal G_A$ square proves the minimization
identity; the Gram determinant formula is the block determinant
identity.  The minimizing residual vanishes precisely if $B$ belongs
to the interior span.  Applying the lower Gram bound before taking
the minimum proves the quantitative inequality.  Inner products and
orthogonal projections are unchanged by basis reparametrization and
simultaneous unitary conjugation.
\end{proof}

\begin{proposition}[Sharp nonlinear-contact and isospectral obstruction]
\label{p3b:prop:v286-heat-incidence-contact-obstruction}
For any $C^2$ Hermitian matrix chart $X(u)$ with $X(0)=A>0$,
$X'(0)=B$, and $X''(0)=J$, the complete second jet of its
Born-half-weight heat trace is
\begin{equation}
 \label{p3b:eq:v286-contact-obstruction}
 \left.\frac{d^2}{du^2}\frac12\operatorname{Tr}e^{-X(u)}
 \right|_{0}=\mathcal G_A(B,B)
 -\frac12\operatorname{Tr}(e^{-A}J).
\end{equation}
A nonzero $B$ need not produce a nonzero second action jet. For
$A=\operatorname{diag}(1,2)$ and any skew-Hermitian $K$ not commuting
with $A$, the positive chart $X(u)=e^{uK}Ae^{-uK}$ has
$B=[K,A]\ne0$ but $\operatorname{Tr}e^{-X(u)}$ is constant and the
whole second jet vanishes. Even in the affine model of the preceding
theorem, taking one interior tangent $C_1=B\ne0$ gives
$\chi_{\rm red}=0$.  These are mathematical countermodels to inferring
physical susceptibility from $P_4\ne0$, positivity of a finite
provider, or a nonzero first operator jet alone.  They do not remove
any previously certified actual physical action incidence.
\end{proposition}
\begin{proof}
The first identity is the full Duhamel formula of
Proposition~\ref{p3b:prop:v245-exact-finite-heat-two-jet}.
Unitary conjugation leaves the full trace invariant; its derivative
$[K,A]$ is nonzero when $K$ fails to commute with $A$. The positive
Gram term is therefore exactly canceled by the $X''$ contact.
Finally $B=C_1$ lies in the interior span, so the exact Schur minimum
of Theorem~\ref{p3b:thm:v286-finite-heat-incidence-gram} vanishes. All these assertions hold on strictly positive operator
charts, and are compatible with local operator-square
representatives $D(u)=\Lambda_g\sqrt{X(u)}$; they do not construct a
generated physical Dirac/action word.
\end{proof}


\begin{definition}[Phase-space presentation algebra]
\label{p3:def:phase-space-algebra}
Let
\[
 \mathcal A_{x,m_x}^{\rm win}:=\End_{\CC}(\mathcal K_{x,m_x}),
 \qquad
 \mathcal A_{p,1}^{\rm win}:=\End_{\CC}(\mathcal K_{p,1})
\]
for the currently constituted position window and the fixed momentum window,
respectively.  The finite-matrix/smooth-phase comparison presentation algebra is
\[
 \mathcal A_{\rm total}
 =\bigl(\mathcal A_{x,m_x}^{\rm win}\otimes C^\infty(S^1_t)\bigr)
  \otimes
  \bigl(\mathcal A_{p,1}^{\rm win}\otimes C^\infty(S^1_E)\bigr).
\]
This is a presentation algebra for the reciprocal spectral datum.  It assigns
no quantum-event or effect-algebra semantics; that use belongs to Paper~IV.
The normalized circle factors are phase variables only and do not imply an
independently metricized scalar circle.
\end{definition}

\subsubsection{The Dirac Spectrum}

\begin{proposition}[Dirac Spectrum on $S^3$]
\label{paper3b:proposition:atlas-dirac-spectrum-on-s-3-b5091ee}

The Dirac operator on $S^3(\Rmom)$ has spectrum:
\begin{equation}
\operatorname{Sp}ec(D_{S^3}) = \left\{ \pm\frac{k + 3/2}{\Rmom} : k = 0, 1, 2, \ldots \right\}
\end{equation}
with degeneracy $(k+1)(k+2)$ for each sign
\cite{CamporesiHiguchi1996}.  Thus a combined two-sign shell has total
degeneracy $2(k+1)(k+2)$; this is the both-sign convention used when comparing
to the Paper~II shell counts.
\end{proposition}
\begin{proof}
This is the spinor part of Paper~I-A,
Theorem~\ref{p1:thm:representation-harmonic-package}.  Paper~III uses the
symmetric spectrum for its dynamical shell and eta-invariant consequences; it
does not rederive the sphere geometry.
\end{proof}

\begin{corollary}[No bare round-Dirac spectral asymmetry]
\label{p3:cor:round-dirac-eta-zero}
For the untwisted round Dirac operator imported from Paper~I-A, the positive and
negative branches occur with equal multiplicity at each \(k\).  Consequently the
round Dirac eta invariant vanishes,
\[
\eta_{D_{S^3}}(0)=0.
\]
Thus the round \(S^3\) Dirac spectrum used in the phase-space spectral triple
supplies gaps and finite shell multiplicities, but it supplies no intrinsic
fermionic spectral-asymmetry source.
\end{corollary}
\begin{proof}
The displayed spectrum is symmetric under \(\lambda\mapsto-\lambda\) with matching
degeneracies \((k+1)(k+2)\).  The signed eta sum cancels term by term, in agreement
with the Paper~I-A eta computation.
\end{proof}


On an auxiliary spinorial total-weight window, the Dirac label is shifted by one half relative to the scalar diagonal label; equivalently the retained integer shells satisfy \(k+\frac12\le m\), hence \(0\le k\le\lfloor m-\tfrac12\rfloor\).

\begin{corollary}[Lowest Dirac spectral gap]
\label{paper3b:corollary:atlas-lowest-dirac-spectral-gap-b395ce9}

The lowest nonzero Dirac scale is:
\begin{equation}
|D_{S^3}|_{\min} = \frac{3}{2\Rmom} = \frac{3\Punit}{2\hbar}.
\end{equation}
On the derived compact momentum constitution, where Corollary~\ref{p3:cor:selfdual-implies-sharp-saturation} gives \(\Punit=\Pmax\), this is equivalently \(3\Pmax/(2\hbar)\).
\end{corollary}
\begin{proof}
By Proposition ``Dirac Spectrum on $S^3$'' (above), the spectrum of $D_{S^3}$ is
\[
\operatorname{Sp}ec(D_{S^3})=\left\{\pm\frac{k+3/2}{\Rmom}:k=0,1,2,\dots\right\}.
\]
Therefore the smallest absolute eigenvalue is attained at $k=0$ and equals
$|D_{S^3}|_{\min}=3/(2\Rmom)=3\Punit/(2\hbar)$ (and equals $3\Pmax/(2\hbar)$ in the equality case $\Punit=\Pmax$).
\end{proof}

\begin{theorem}[Spinorial shell-ceiling dictionary]
\label{p3:thm:spinorial-shell-ceiling-dictionary}
Let \(D_{S^3}\) be the round Dirac operator and let \(m\) be the Paper-II
total-weight truncation label.  The canonical dimensionless spinorial shell
operator is
\begin{equation}
\label{p3:eq:spinorial-total-weight-operator}
        \widehat N_D:=\Rmom|D_{S^3}|-1.
\end{equation}
On the shell with Dirac label \(k\),
\[
        \widehat N_D=k+\frac12.
\]
Consequently the Paper-II condition \(k+\frac12\le m\) is equivalently
\begin{equation}
\label{p3:eq:spinorial-shell-ceiling}
        \widehat N_D\le m
        \quad\Longleftrightarrow\quad
        |D_{S^3}|\le\Lambda_D(m),
        \qquad
        \boxed{\Lambda_D(m):=\frac{m+1}{\Rmom}.}
\end{equation}
At the constituted cell \(m_p=1\),
\[
        \Lambda_D(1)=\frac2{\Rmom},
\]
so the \(k=0\) shell is retained and the \(k=1\) shell is excluded.  The
physical momentum-unit ceiling
\(\Lambda_{\rm mom}=\Pmax/\hbar=1/\Rmom\) at self-duality is therefore not the
unshifted round-Dirac eigenvalue ceiling; the difference is the canonical
spinorial representation offset in \eqref{p3:eq:spinorial-total-weight-operator}.
\end{theorem}

\begin{proof}
The round spectrum gives
\(\Rmom|D_{S^3}|=k+\frac32\).  Subtracting the unit spinorial offset yields
\(\widehat N_D=k+\frac12\), exactly the total-weight label imported from
Paper~II.  Solving \(\widehat N_D\le m\) for \(|D_{S^3}|\) gives
\(|D_{S^3}|\le(m+1)/\Rmom\).  For \(m=1\),
\[
 \frac{3}{2\Rmom}<\frac2{\Rmom}<\frac{5}{2\Rmom},
\]
which retains exactly the lowest shell.
\end{proof}

\begin{remark}[Three cutoff roles]
\label{p3:rem:three-cutoff-roles}
The equality \(\Lambda_{\rm mom}=1/\Rmom\) fixes the vector-momentum unit at
the self-dual cell.  On the already derived seed, the auxiliary lowest-spinor shell condition is
\(\widehat N_D\le1\), equivalently \(\Lambda_D(1)=2/\Rmom\).  The common
spectral-action regulator \(\Lcut\) is a profile scale whose Einstein--Hilbert
combination is calibrated by Newton's constant.  A sharp count written with
the unshifted operator \(D_{S^3}\) uses a variable Dirac threshold and must not
identify that threshold with \(\Lambda_{\rm mom}\).
\end{remark}


\subsubsection{Protected Clifford Octet and Lowest Dirac Window}

The protected eight-channel structure is a Clifford-operator statement, not
a dimension-count identification with one chiral Dirac eigenspace.  Its
physical use in this paper requires only the imported nondegenerate
protected-coordinate stratum \((\mathrm{ND}_{\rm sp})\).  Paper~II-C proves
there that the normalized even generators are exactly \(C_iC_4\), so the
represented protected Clifford channel algebra is independent of the
remaining nonzero completion modulus.  The lowest Dirac window below remains
a separate analytic comparison object.

\begin{remark}[Relation to the protected diagonal-weight carrier]
\label{p3:rem:8mode-protected-carrier}
Paper~II derives the protected complex spinor module
\[
 \mathcal S_{\rm prot}
 =V_{\frac12,0}\oplus V_{0,\frac12},
 \qquad \dim_{\mathbb C}\mathcal S_{\rm prot}=4.
\]
On \((\mathrm{ND}_{\rm sp})\), Paper~II-C proves that the normalized
protected Euclidean vector has the exact represented even algebra
\[
 A_{\rm ev}
 \cong \Cl^0_{4,0}
 \cong \Cl(0,3),
 \qquad \dim_{\mathbb R}A_{\rm ev}=8.
\]
This is the protected Clifford octet used in the dynamical scalar-projection
arguments below.  The lowest positive round-Dirac eigenspace is a separate
analytic two-complex-dimensional window.  Equality of its endomorphism-space
dimension with eight does not define a map to \(A_{\rm ev}\).
\end{remark}

\begin{definition}[Lowest Dirac projector]
\label{p3:def:lowest_dirac_projector}
Let \(\kappa_0:=3/(2\Rmom)\).  Denote by \(P_+\) the spectral projector onto
the lowest positive Dirac eigenspace,
\[
 P_+:=\chi_{\{\kappa_0\}}(D_{S^3}),
 \qquad
 S_+:=P_+\mathcal H_{\rm spin}.
\]
By the round Dirac spectrum, \(\dim_{\mathbb C}S_+=2\).
\end{definition}

\begin{theorem}[Protected Clifford octet and independent lowest-window count]
\label{p3:thm:8mode_operator_basis}
Assume \((\mathrm{ND}_{\rm sp})\).  Then:
\begin{enumerate}
\item the imported protected even operator algebra satisfies
\[
 A_{\rm ev}\cong\Cl(0,3),
 \qquad
 \dim_{\mathbb R}A_{\rm ev}=8;
\]
\item under the standard Chevalley vector-space identification,
\[
 A_{\rm ev}
 \cong
 \Cl^{(0)}\oplus\Cl^{(1)}\oplus\Cl^{(2)}\oplus\Cl^{(3)},
 \qquad
 \dim(\Cl^{(0)},\Cl^{(1)},\Cl^{(2)},\Cl^{(3)})=(1,3,3,1),
\]
and hence, as an \(SO(3)\)-module,
\[
 \mathbf 8=\mathbf1\oplus\mathbf3\oplus\mathbf3\oplus\mathbf1;
\]
\item independently, the analytic lowest positive Dirac window satisfies
\[
 \dim_{\mathbb R}\End_{\mathbb C}(S_+)=8.
\]
The equality of the two real dimensions in (i) and (iii) is only a numerical
count.  The physical protected eight-channel algebra used below is
\(A_{\rm ev}\), inherited on \((\mathrm{ND}_{\rm sp})\) by the exact
Paper~II-C even-algebra theorem; no algebraic or
canonical identification
\(\End_{\mathbb C}(S_+)\cong\Cl(0,3)\) is asserted.
\end{enumerate}
\end{theorem}

\begin{proof}
Items (i)--(ii) follow from Paper~II-C,
Corollary~\ref{p2c:cor:nondegenerate-protected-clifford-orbit}, together with the standard
Chevalley vector-space decomposition of \(\Cl(0,3)\)\cite{lawson1989spin}.
For (iii), Definition~\ref{p3:def:lowest_dirac_projector} gives
\(\dim_{\mathbb C}S_+=2\), hence
\(\End_{\mathbb C}(S_+)\cong M_2(\mathbb C)\) and therefore
\(\dim_{\mathbb R}\End_{\mathbb C}(S_+)=8\).  Since
\(M_2(\mathbb C)\) and \(\Cl(0,3)\cong\mathbb H\oplus\mathbb H\) are not
isomorphic as real algebras, this last equality supplies no bridge between
the two typed objects.
\end{proof}

\begin{remark}[Operator channels versus Hermitian observables]
The protected octet \(A_{\rm ev}\) is a real Clifford operator algebra.  It
should not be confused with either the four-dimensional real vector space of
Hermitian endomorphisms of a two-complex-dimensional lowest eigenspace or the
full complex algebra \(\End_{\mathbb C}(S_+)\).  Later reality,
self-adjointness, or field conditions may impose further restrictions, but
those restrictions do not constitute the octet itself.
\end{remark}

\begin{remark}[Clifford sign and operator realization]
\label{p3:rem:clifford-sign-operator-realization}
Throughout the protected scalar-projection arguments,
\(\Cl(0,3)\) denotes the real algebra generated by
\(e_ie_j+e_je_i=-2\delta_{ij}\).  Thus
\[
 \Cl(0,3)\cong\mathbb H\oplus\mathbb H.
\]
By contrast, \(\End_{\mathbb C}(S_+)\cong M_2(\mathbb C)\).  These algebras
have the same real dimension but are not identified.  The physical Clifford
channel provenance is the imported \(A_{\rm ev}\cong\Cl(0,3)\) theorem;
the lowest-Dirac endomorphism count is retained only as an independent
analytic comparison.
\end{remark}

\begin{remark}[Carrier-level Clifford independence and finite-resolution dependence]
\label{p3:rem:hurwitz-link-origin-octet}
The protected octet is already inherited on \((\mathrm{ND}_{\rm sp})\)
and does not require a finite-support realization.  If a Paper-I-A
binary-tetrahedral core is also actualized, the Hurwitz-link Walsh algebra
provides a concrete model
\[
 (\mathbb R^{N(v)},\star_{\sigma_-})\cong\Cl(0,3)
\]
for any Hurwitz-unit vertex \(v\).  This is an optional realization of the
same Clifford algebra type, not the provenance of the protected octet and not
a map from the lowest Dirac endomorphism algebra.

This independence is a carrier statement only.  The intrinsic grade split of
\(\Cl(0,3)\) is exceptional-support independent, whereas a finite hierarchy
of local, interface and global resolutions of its even quaternionic unit
group may be generated by exceptional subgroups of
\(U_1(\Cl^{\rm even}(0,3))\cong\operatorname{Sp}(1)\).  Paper~III-C derives that resolution
on its explicitly declared mobile-Clifford branch; no such finite branch is
created by the Clifford algebra alone.
\end{remark}

\paragraph{Explicit protected-channel bookkeeping}

A convenient Clifford basis compatible with the grade decomposition is
\begin{center}
\begin{tabular}{cll}
\toprule
Channel & Representative & \(SO(3)\) type \\
\midrule
1 & \(1\) & scalar \\
2--4 & \(e_1,e_2,e_3\) & vector \\
5--7 & \(e_2e_3,\ e_3e_1,\ e_1e_2\) & bivector \\
8 & \(e_1e_2e_3\) & pseudoscalar \\
\bottomrule
\end{tabular}
\end{center}

\paragraph{Lowest-window mode exclusion}
Projecting with \(P_+\) removes higher positive Dirac harmonics from the
analytic lowest-window comparison.  This spectral projection is independent
of the protected Clifford octet: the latter is the imported
\(A_{\rm ev}\cong\Cl(0,3)\) algebra on \((\mathrm{ND}_{\rm sp})\).  No
claim is made that projecting to \(S_+\) generates the Clifford channels.

\subsection{The Spectral Action}
\label{p3:sec:spectralaction}


\subsubsection{Characterization of admissible finite spectral actions}
\label{p3:sec:uniqueness_spectral_action}


In the fuzzy truncation each sector Dirac operator $D_x,D_p$ is a finite-dimensional
self-adjoint matrix.  Unitary equivalence implements automorphisms of the unlabelled
spectral datum, so admissible scalar valuations are spectral.  Disjoint finite
blocks are represented by orthogonal central projectors in the generated
B-test algebra.  For an actual unmarked closed occurrence, direct-sum
additivity is a consequence of canonical closure on the licensed biproduct;
it is not attributed to GP3.

\begin{lemma}[Direct-sum additivity of the closed occurrence]
\label{p3:lem:additivity}
Let \(D_1\) and \(D_2\) be disjoint finite spectral blocks.  The declared
direct-sum constructor and Theorem~\ref{p3:thm:gp2-dcl-derivation} give
\[
        S(D_1\oplus D_2)=S(D_1)+S(D_2).
\]
Thus block additivity is an exact derived property of the closed-loop
realization of \(\Dcl\).  The later
characterization theorem shows how that clause is realized by ordinary finite
trace linearity.
\end{lemma}

\begin{proof}
Functional calculus respects block diagonals and the loop theorem closes the
result by ordinary finite trace.  Applying
\(\operatorname{Tr}(A\oplus B)=\operatorname{Tr}A+\operatorname{Tr}B\)
gives the formula.  No profile selection or GP3 premise is used.
\end{proof}

\begin{remark}[Why multi-trace functionals are excluded]
A multi-trace expression such as $(\Tr f(D))^2$ is not additive on orthogonal
alternatives: for $D=D_1\oplus D_2$ it creates the cross term
$2\Tr f(D_1)\Tr f(D_2)$.  That term is an extra correlation rule not contained
in a primitive additive loop.  Its exclusion is therefore a typing statement
about \(\Dcl\), not a GP3 theorem or an imported cluster-decomposition axiom.
\end{remark}


\begin{theorem}[Converse finite spectral-action characterization]
\label{p3:thm:spectral-characterization}
Let a finite sector datum, or a finite spectral window in the sense of Definition~\ref{p3:def:D1-finite-window}, consist only of an unlabelled finite Hilbert space
\(\mathcal H\), a self-adjoint compressed Dirac operator \(D\), and one cutoff scale \(\Lcut\).  Its
isomorphisms are unitary equivalences
\[
        (\mathcal H,D)\longmapsto(\mathcal H,UDU^{-1}).
\]
Conversely, suppose a family of finite sector actions satisfies:
\begin{enumerate}[label=(\roman*),leftmargin=2.1em]
\item \emph{unitary naturality}: \(S(UDU^{-1})=S(D)\) for every unitary \(U\);
\item \emph{block additivity}: \(S(D_1\oplus D_2)=S(D_1)+S(D_2)\) for disjoint finite
spectral blocks, as recorded in
Lemma~\ref{p3:lem:additivity};
\item \emph{Born/Dirac evenness}: \(S(D)=S(-D)\);
\item \emph{single-cutoff scaling}: the only dimensionful scale entering the
functional is \(\Lcut\), and for every \(s>0\) the exact covariance law is
\[
 S(sD;s\Lcut)=s^wS(D;\Lcut)
\]
for one declared output weight \(w\in\mathbb R\);
\item \emph{no extra labelled datum}: no fixed projector, grading insertion, channel
endomorphism, independent profile, or external field is part of the sector datum.
\end{enumerate}
Then every admissible action has single-profile trace form: there exists one
common cutoff profile \(f\), shared across the unlabelled blocks, such that
\[
        \boxed{
        S(D;\Lcut)=\Lcut^{\,w}
        \Tr_{\mathcal H} f(D^2/\Lcut^2).}
\]
The normalized dimensionless \(\Dcl\) convention is the specialization
\(w=0\).  Any further dimensionful coefficient is a separately declared
calibration or coupling, not part of this characterization.
\begin{remark}[Exact hypothesis boundary of the five conditions]
\label{p3:rem:dcl-conditions-derived}
This theorem is the abstract converse and does not need an occurrence premise.
On the positive chain, however, the conditions have distinct provenance:
\begin{enumerate}[label=(\roman*),leftmargin=2.1em]
\item unitary naturality and canonical normalization follow from the unmarked
closed-loop occurrence;
\item block additivity follows when the direct-sum constructor is licensed;
\item Dirac evenness follows because the generated positive process is applied
to \(D^2\), not from loop topology alone;
\item single-cutoff scaling follows only when the occurring process interface
exports one resolution scale and a fixed output weight \(w\);
\item absence of extra labels is exactly the unmarked-loop typing; generated
projectors, gradings, fields, or sources define marked or boundary branches.
\end{enumerate}
The profile remains the occurring process law.  Neither this converse nor the
closed-loop derivation makes it a consequence of GP1--GP3.
\end{remark}

Thus the spectral action has the displayed form under the stated converse
hypotheses and, on the positive chain, inside the loop-derived class \(\Dcl\); it is not an
absolute exclusion of all datum-enlarged actions whose additional geometric datum has been explicitly supplied.  For a single finite spectrum this determines only the values of \(f\) on that finite set; Lemma~\ref{p3:lem:finite-window-spectral-separation} shows that finite-window separation is supplied by intrinsic spectral calculus.  Choosing smooth representatives across a directed family of windows is the continuum comparison convention used to compare sectors and take semiclassical limits.
On the native unmarked reciprocal branch, Paper~II-A,
Theorem~\ref{p2a:thm:generated-born-inheritance}, supplies the Born covariance
of the generated scalar action; the trivial scalar target makes it invariant.
Together with the unlabelled single-profile constructor this gives the same
profile on the reciprocal sectors and hence the physical action of
Definition~\ref{p3:def:physical_spectral_action}.  For an auxiliary functional
without that frozen constructor origin, Born symmetry remains an explicit
class condition rather than an inferred property.
\end{theorem}

\begin{proof}
\ProofPlan{We prove the finite spectral-action form in four steps: (a) reduce the functional to the eigenvalue multiset by unitary naturality; (b) use block additivity to obtain a trace functional; (c) use evenness to pass to \(D^2\); (d) impose the single admissible cutoff scale.}
\begin{proofsteps}
\item \textbf{Unitary naturality and eigenvalue dependence.}
Fix a finite sector and write \(D\) for its Dirac operator.  By the spectral theorem,
\(D\) is unitarily diagonalizable.  Therefore unitary naturality implies that
\(S(D)\) depends only on the multiset of eigenvalues of \(D\).  Hence, for
\[
        D=\operatorname{diag}(\kappa_1,\ldots,\kappa_n),
\]
there is a symmetric function \(F_n\) such that
\[
        S(D)=F_n(\kappa_1,\ldots,\kappa_n).
\]


\item \textbf{Block additivity and trace form.}
Block additivity gives, for all finite lists \(\kappa\) and \(\lambda\),
\[
F_{n+m}(\kappa_1,\ldots,\kappa_n,\lambda_1,\ldots,\lambda_m)
 =
F_n(\kappa_1,\ldots,\kappa_n)+F_m(\lambda_1,\ldots,\lambda_m).
\]
Applying this identity to one-dimensional blocks and iterating gives
\[
        F_n(\kappa_1,\ldots,\kappa_n)=\sum_{i=1}^n F_1(\kappa_i).
\]
Define \(\phi(t):=F_1(t)=S([t])\).  Functional calculus then gives
\[
        S(D)=\sum_i \phi(\kappa_i)=\Tr\phi(D).
\]


\item \textbf{Evenness and dependence on \(D^2\).}
Evenness gives \(S(D)=S(-D)\) for every finite self-adjoint \(D\).  Applying this to
one-dimensional blocks forces \(\phi(t)=\phi(-t)\).  Hence there is a function
\(\widetilde f\) on the non-negative spectrum such that
\[
        \phi(t)=\widetilde f(t^2).
\]

\item \textbf{Cutoff-scale admissibility.}
Apply the stated covariance to a one-dimensional block and choose
\(s=\Lcut^{-1}\).  Then
\[
 \phi(t;\Lcut)
 =\Lcut^{\,w}\phi(t/\Lcut;1).
\]
Evenness makes \(u\mapsto\phi(u;1)\) a function \(f(u^2)\) on
\(u^2\ge0\).  Hence
\(\phi(t;\Lcut)=\Lcut^{\,w}f(t^2/\Lcut^2)\), and therefore
\[
        S(D;\Lcut)=\Lcut^{\,w}\Tr f(D^2/\Lcut^2).
\]
The same argument applies to \(D_x\) and \(D_p\).  Born reciprocity identifies
the two sector profiles inside the Born-symmetric uninserted class.  Setting
\(w=0\) recovers the normalized \(\Dcl\) form used elsewhere in this paper.
\end{proofsteps}
\end{proof}

\begin{remark}
This theorem is the finite-window characterization result actually needed in
Paper~III.  The theorem does not say that no other writable functional exists.  It says
that, given only the unlabelled finite spectral datum or finite spectral window and the \(\Dcl\) admissibility
requirements, the action is forced to be an uninserted single-profile trace.  The broader
Chamseddine--Connes spectral-action literature is then used not to choose the functional
class anew, but to supply the continuum heat-kernel expansion and the almost-commutative
comparison formulas that apply once the relevant finite spectral data are fixed
\cite{ChamseddineConnes1997,Connes1996,GraciaBondia2001,ChamseddineConnesvanSuijlekom2020Entropy}.
\end{remark}

\begin{corollary}[Additive-collapse test for multitrace statistics]
\label{p3:cor:no-multitrace-finite-datum}
Inside the finite-window \(\Dcl\) class, a nonlinear symmetric spectral statistic such as
\[
        (\Tr f(D^2/\Lcut^2))^2,
        \qquad
        \Phi\!\left(\Tr f_1(D^2/\Lcut^2),\ldots,\Tr f_N(D^2/\Lcut^2)\right),
\]
is an uninserted action only if the resulting function of trace coordinates is
additive on the attained direct-sum cone and vanishes at the zero block.  In
that case the spectral characterization collapses it to the already classified
single-profile trace.  Otherwise it correlates disconnected finite blocks and
violates block additivity, so it requires an additional global
coupling/source datum and lies outside \(\Dcl\).
\end{corollary}

\begin{proof}
For two disconnected finite spectral blocks, Theorem~\ref{p3:thm:spectral-characterization} requires
\[
        S(D_1\oplus D_2)=S(D_1)+S(D_2).
\]
A function \(\Phi\) of trace coordinates survives block additivity exactly
when, on the attained trace cone,
\[
 \Phi(x+y)=\Phi(x)+\Phi(y),\qquad \Phi(0)=0.
\]
If this identity holds, the resulting spectral functional is additive and the
preceding characterization theorem collapses it to the already classified
single-profile trace on the stated unlabelled domain.  If it fails, evaluating
on \(D_1\oplus D_2\) produces a cross-block term and violates the defining
additivity.  No assertion that an abstract additive function is linear or
affine is needed; such a conclusion would require an additional regularity
hypothesis.  The nonadditive case is therefore an inserted global coupling
between independent blocks.
\end{proof}


\begin{definition}[Uninserted sector functional inside \texorpdfstring{$\Dcl$}{D1}]
\label{p3:def:uninserted-sector-functional}
Let \(D\) be one finite sector Dirac operator acting on a Hilbert space \(\mathcal H\).  A
sector functional is called \emph{uninserted} when it is a natural functional of
\((\mathcal H,D,\Lcut)\) alone:
\[
        S(D)=\Tr\,\phi(D)
        \quad\text{or, in the even single-cutoff case,}\quad
        S(D)=\Tr f(D^2/\Lcut^2),
\]
and contains no additional labelled endomorphism \(B\in\End(\mathcal H)\), such as
\[
        \Tr\!\bigl(B f(D^2/\Lcut^2)\bigr),
        \qquad
        \sum_g c_g\Tr\!\bigl(P_g f_g(D^2/\Lcut^2)\bigr).
\]
Equivalently, the functional is natural under unitary equivalence of the unlabelled
finite spectral datum \((\mathcal H,D)\) without carrying any separately labelled
subbundle, grade, channel projector, source, independent profile, or background field.
\end{definition}


\begin{lemma}[Finite-window separation by functional calculus]
\label{p3:lem:finite-window-spectral-separation}
For a finite spectral datum \((\mathcal H,D,\Lcut)\), ordinary finite
functional calculus can separate every spectral cluster of \(D^2\) from its
complement by a polynomial evaluated on \(D^2/\Lcut^2\).  Whether the
occurring process family contains all such separating profiles is an explicit
analytic condition on the \(\Dcl\) datum, not a consequence of GP1.
\end{lemma}

\begin{proof}
Let the distinct eigenvalues of \(D^2/\Lcut^2\) in the finite window be
\[
        \lambda_1,\ldots,\lambda_N .
\]
For any cluster \(C\subset\{\lambda_1,\ldots,\lambda_N\}\), finite interpolation gives a polynomial \(p_C\) with
\[
        p_C(\lambda_i)=1\quad(\lambda_i\in C),
        \qquad
        p_C(\lambda_j)=0\quad(\lambda_j\notin C).
\]
By the spectral theorem, \(p_C(D^2/\Lcut^2)\) is exactly the spectral
projector onto the cluster \(C\), but it has been produced as a function of
the unlabelled operator \(D\), not by adjoining a labelled projector.  The
construction proves mathematical availability in the finite functional
calculus.  Definition~\ref{p3:def:D1}, rather than GP1, determines which
common profiles are admissible action data.
\end{proof}

\begin{definition}[Finite-window separating profile class]
\label{p3:def:finite-window-separating-profile}
Inside the loop-derived class \(\Dcl\), call a finite-window profile class
\(\mathcal F\) \emph{separating} when, for every finite spectral datum and
every finite spectral cluster \(C\) in the spectrum of \(D^2\), there exists
\(f\in\mathcal F\) whose values distinguish \(C\) from its complement.
Lemma~\ref{p3:lem:finite-window-spectral-separation} proves that finite
functional calculus supplies such functions mathematically; the
occurring process family in \(\Dcl\) is assumed to contain the separating
class used below.
This condition does not require the occurring \(f\) to be a characteristic
function or to produce an exact projector: only distinct values on the chosen
cluster and its complement are required.  Lemma~\ref{p3:lem:finite-window-spectral-separation}
records the mathematical availability of exact projectors in the full finite
functional calculus, but it does not infer that every interpolating polynomial
occurs.  Without the displayed separation property, the no-hidden-insertion
conclusion below holds only on the spectral subspaces visible to the allowed
profiles.
\end{definition}

\begin{lemma}[GP2 datum-enlargement test for inserted finite-sector terms]
\label{p3:lem:insertions-scalar}
Consider a candidate inserted finite-sector action of the form
\[
        S_B(D)=\Tr\!\left(B f(D^2/\Lcut^2)\right),
\]
or a finite sum of such terms.  Assume that the occurring common profile
family is separating in the sense of
Definition~\ref{p3:def:finite-window-separating-profile}.  Inside the
unlabelled \(\Dcl\) datum \((\mathcal H,D,\Lcut)\), exactly one of the
following happens:
\begin{enumerate}[label=(\roman*),leftmargin=2.1em]
\item \(B\) is not an independent datum and the resulting functional depends only on
\(D\).  Then Theorem~\ref{p3:thm:spectral-characterization} already forces the action to
collapse to \(\Tr\psi(D^2/\Lcut^2)\) for a single profile \(\psi\).
\item \(B\) is fixed but not transformed as part of the datum.  Then unitary naturality
forces \(B=b\mathbf 1\), so the insertion is only
\(b\Tr f(D^2/\Lcut^2)\), absorbed into normalization.
\item \(B\) is a genuine non-scalar projector, grading, channel label, internal
endomorphism, or background field transported together with the data.  Then the datum is
not \((\mathcal H,D,\Lcut)\) but \((\mathcal H,D,\Lcut,B)\), which is a datum-enlarged inserted
class rather than \(\Dcl\).
\end{enumerate}
Thus no non-scalar hidden insertion survives inside \(\Dcl\).
If the occurring profile family is not separating, the same conclusion holds
only after quotienting \(\End(\mathcal H)\) by the annihilator
\[
 \mathcal N_{\mathcal F}
 :=\left\{C:
 \Tr\!\bigl(Cf(D^2/\Lcut^2)\bigr)=0
 \text{ for every allowed }(D,f)\right\}.
\]
In that weaker domain an insertion may have an invisible representative in
\(\mathcal N_{\mathcal F}\); it cannot be claimed scalar as an operator.
\end{lemma}

\begin{proof}
If the candidate action depends only on \(D\), it is covered by the classification
Theorem~\ref{p3:thm:spectral-characterization}; unitary naturality and block additivity
make it a linear spectral statistic \(\Tr\phi(D)\), and evenness with one cutoff puts it in
the form \(\Tr\psi(D^2/\Lcut^2)\).

Now suppose \(B\) is fixed but is not part of the transformed datum.  Naturality requires
\[
        \Tr\!\left(B f(UD^2U^{-1}/\Lcut^2)\right)
        =
        \Tr\!\left(B f(D^2/\Lcut^2)\right)
\]
for every self-adjoint \(D\), every unitary \(U\), and every finite-window separating admissible profile in the common spectral-action class.  By cyclicity,
\[
        \Tr\!\left((U^{-1}BU-B)f(D^2/\Lcut^2)\right)=0.
\]
Fix an arbitrary rank-one projector \(Q\).  Choose two admissible numbers
\(u,v\ge0\) with \(u\ne v\), and choose the self-adjoint datum so that
\[
 \frac{D^2}{\Lcut^2}=uQ+v(\mathbf1-Q).
\]
By the assumed separation property there is an occurring common profile
\(f\in\mathcal F\) with \(\alpha:=f(u)\ne f(v)=:\beta\).  Functional
calculus then gives the exact two-value identity
\[
 f(D^2/\Lcut^2)=\alpha Q+\beta(\mathbf1-Q);
\]
no occurrence of an interpolating projector profile has been assumed.  Put
\(C:=U^{-1}BU-B\).  Cyclicity gives \(\Tr C=0\), and naturality therefore
implies
\[
 0=\Tr\!\left(C[\alpha Q+\beta(\mathbf1-Q)]\right)
   =(\alpha-\beta)\Tr(CQ).
\]
Thus \(\Tr(CQ)=0\) for every rank-one projector \(Q\).  Equivalently,
\(\langle v,Cv\rangle=0\) for every unit vector \(v\).  The polarization
identity gives \(U^{-1}BU-B=0\).  Since this holds for every unitary \(U\), \(B\) commutes
with the full matrix algebra \(\End(\mathcal H)\).  Therefore \(B\) lies in its center,
so \(B=b\mathbf 1\).

Finally, if \(B\) is non-scalar and is nevertheless preserved as meaningful structure,
then it is additional labelled data.  The datum has been enlarged from
\((\mathcal H,D,\Lcut)\) to \((\mathcal H,D,\Lcut,B)\).  That may define a writable
datum-enlarged inserted theory, but it is not the \(\Dcl\)-uninserted sector datum.
Without separation, the polarization argument sees only the quotient by
\(\mathcal N_{\mathcal F}\), which proves precisely the weaker terminal
statement of the lemma.
\end{proof}

\begin{corollary}[Classification of hidden channel profiles inside \texorpdfstring{$\Dcl$}{D1}]
\label{p3:cor:no-channel-profiles}
Inside the uninserted sectorwise spectral-action class of
Definition~\ref{p3:def:uninserted-sector-functional}, every admissible sector action is an
uninserted single-profile spectral trace.  A term such as
\[
        \sum_g c_g\Tr\!\left(P_g f_g(D^2/\Lcut^2)\right)
\]
has only three possibilities:
\begin{enumerate}[label=(\roman*),leftmargin=2.1em]
\item the \(P_g\) are functions of \(D\), so the whole expression collapses into a single
spectral profile;
\item the \(P_g\) are scalar, so they only renormalize the common profile or the overall
normalization;
\item the \(P_g\) are genuine channel projectors, so they are extra labelled endomorphism
data and the action lies outside \(\Dcl\).
\end{enumerate}
Thus \(\Dcl\) does not merely assume away hidden channel insertions; it classifies them.
They either collapse to the uninserted trace or require additional data.
\end{corollary}

\begin{proof}
The finite sum must first be classified as one scalar output.  If the displayed
sum is genuinely an unlabelled natural, block-additive functional of the
unmarked datum \(D\) alone, Theorem~\ref{p3:thm:spectral-characterization}
applies directly to the aggregate and gives
\[
 \sum_g c_g\Tr\!\left(P_g f_g(D^2/\Lcut^2)\right)
 =\Tr\psi(D^2/\Lcut^2)
\]
for one admissible spectral profile \(\psi\).  Aggregate naturality does
\emph{not} imply that each displayed summand is separately natural: complementary
non-scalar projectors \(P,Q\) already give
\(\Tr(Pf(D^2))+\Tr(Qf(D^2))=\Tr f(D^2)\).  If the individual \(P_g\) are
retained as meaningful marks, they enlarge the constructor datum even when
their scalar output happens to cancel.  If each summand is to be classified
separately, separate summand naturality is an additional premise, and then
Lemma~\ref{p3:lem:insertions-scalar} applies to that summand.  Hence the
unlabelled aggregate still collapses to one spectral trace, whereas genuine
retained channel marks remain outside \(\Dcl\).
\end{proof}


\begin{theorem}[Casewise GP closure; no universal selection rule]
\label{p3:thm:D1-universal-constitution}
\label{p3:thm:D1-gp2-constitution-rule}
Fix a closure problem \(C\) on the completed Born/contact geometry and let
\(\mathcal B_C\) be the class of admissible finite constitutions realizing
\(C\).  A least-element conclusion is an application of recursive GP3 only
when
\[
 \mathcal B_C\subseteq\operatorname{Adm}_R
\]
and the displayed preorder is a generated role-faithful structural
comparison between complete GP1-admissible, GP2-exact endpoints.  Under those
hypotheses a unique nonzero closed least element is retained by GP3; no
additional selector is required.  The three recurring cases have the
following exact status:
\begin{enumerate}[label=(\roman*),leftmargin=2.2em]
\item \emph{Reciprocal fixed point.}  For a positive Born-exchanged pair
\((X,Y)\) on \(XY=C_{XY}>0\), additivity and Born evenness make the reduced
valuation proportional to \(X+Y\), and AM--GM proves that its unique numerical
minimum is \(X=Y=\sqrt{C_{XY}}\).  This becomes a physical GP3 conclusion
only when a generated variational or stability constructor makes comparison
by that valuation a role-faithful structural comparison.  AM--GM alone does
not constitute such a constructor.
\item \emph{First nonzero order.}  If a sector carries a generated ordered
filtration \(\mathcal F^0\subset\mathcal F^1\subset\cdots\) and the lower-order
invariants vanish by role symmetry, trace parity, chirality,
commutator parity, collapse to a lower-dimensional face, or rephasing gauge
invariance, then
\[
        \nu_*:=\min\{n:\mathcal I_n\ne0\}
\]
is the leading nonvanishing index.  GP2 must nevertheless retain every higher
generated nonzero invariant; the leading-order theorem is not permission to
delete the rest.
\item \emph{Exact-fidelity endpoint.}  For
\(0\le r\le r_{\max}\), the endpoint \(r=r_{\max}\) follows from GP2 only
when a generated normalized relation is realized exactly there and every
interior value loses that relation or inserts an ungenerated defect channel,
as in Paper~II-A,
Theorem~\ref{p2:thm:gp2-exact-fidelity-endpoint}.  Boundedness, injectivity,
positivity, or the absence of a canonical interior number does not force
saturation.  Otherwise the interval remains an unresolved family.
\end{enumerate}
Thus every later use must cite its theorem-specific structural comparison,
variational/stability constructor, lower-order vanishing proof, or
exact-fidelity relation.  If incomparable irreducible elements remain, the
framework leaves them unresolved.
\end{theorem}

\begin{proof}
The first paragraph is exactly the recursive GP3 comparison rule of
Paper~II-A applied only after candidate-internal GP2 exactification.  The
reciprocal minimum follows from
\(X+Y\ge2\sqrt{XY}\), with equality at \(X=Y\), but a physical comparison
requires the additional generated variational/stability typing stated in the
theorem.  The filtered claim is merely the definition of the first
nonvanishing index after the lower-order vanishing theorems; GP2 no-omission
retains all later generated terms.  The endpoint claim is
Theorem~\ref{p2:thm:gp2-exact-fidelity-endpoint}.  These exhaust the asserted
casewise conclusions and introduce no universal selector.
\end{proof}

\begin{theorem}[Exact hypothesis boundary for statements internal to \texorpdfstring{$\Dcl$}{D1}]
\label{p3:thm:D1-hypothesis-boundary}
Every statement in Paper~III of the form ``\(X\) is fixed inside
\(\Dcl\)'' has the following exact logical form.  Let
\[
        \mathfrak D=(\mathcal H,D,\Lcut;\mathcal F)
\]
be the completed unlabelled finite spectral datum, together with one common
admissible separating profile class \(\mathcal F\), and let isomorphisms be
unitary conjugacies of \((\mathcal H,D)\).  Then \(X\) is forced only when the
relevant closure problem either has a unique nonzero closed least realization
under a generated role-faithful GP3 structural comparison, or satisfies a
theorem-specific generated variational, vanishing, or exact-fidelity relation
of Theorem~\ref{p3:thm:D1-universal-constitution}, using only:
\begin{enumerate}[label=(\roman*),leftmargin=2.1em]
\item the completed Born/contact geometry and the scalar Hamiltonian/Born pair already constituted before \(\Dcl\) is invoked;
\item finite-window unitary naturality, direct-sum trace additivity,
Born/Dirac evenness, and one common spectral-action profile regulator, kept
distinct from the physical carrier and Dirac-shell ceilings;
\item functional-calculus projectors generated by \(D^2\) itself, but no externally labelled projector, channel weight, endomorphism, branch selector, sector-dependent profile, or independent source coefficient;
\item the local normalization data explicitly displayed in the theorem statement when a normalization-conditional constant is being computed.
\end{enumerate}
Consequently a \(\Dcl\) conclusion is never a claim of derivability from GP
closure or topology alone.  It is internal to the occurring spectral datum and
its local normalizations.  The physical momentum seed is already fixed
upstream by GP2 exact/host-faithful constitution.  A leading-order conclusion
requires the displayed lower-order vanishings and retains all higher generated
terms; a saturated endpoint requires the GP2 exact-fidelity hypotheses.
Alternative action branches in
Definition~\ref{p3:def:dynamical-branch-atlas} remain admissible when their
own generated data and occurrences are stated.
\end{theorem}

\begin{proof}
Definitions~\ref{p3:def:D1} and~\ref{p3:def:D1-finite-window} specify the
datum and finite-window reading.  Theorem~\ref{p3:thm:spectral-characterization}
reduces every unlabelled finite action in this branch to a single-profile
trace.  Lemma~\ref{p3:lem:insertions-scalar} and
Corollary~\ref{p3:cor:no-channel-profiles} classify inserted projectors and
channel profiles as either collapsible to that trace or as datum enlargements.
Theorem~\ref{p3:thm:D1-universal-constitution} then supplies only the
casewise GP2/GP3 consequences actually proved.  Thus the remaining ambiguity
inside \(\Dcl\) is the explicitly named common-profile/normalization freedom
plus any unresolved theorem-specific family.  When a later theorem supplies
the required structural comparison or exact relation, the result is unique
under those hypotheses; otherwise it remains conditional or unresolved.
\end{proof}

\subparagraph{Framework reading of \(\Dcl\)-statements.}
Theorem~\ref{p3:thm:D1-hypothesis-boundary} fixes the paper-wide
interpretation of ``fixed'', ``unique'', and ``least/irreducible'' inside \(\Dcl\):
the completed finite spectral datum, one shared admissible profile, and displayed
local normalizations are active hypotheses; variational selection is never
silently supplied by GP3.

\subsubsection{The Physical Cutoff}

\begin{proposition}[Separated geometric and Dirac-shell cutoff scales]
\label{p3:prop:cutoff_EH_normalized}
The round momentum parent fixes the geometric momentum scale
\begin{equation}
\Lambda_{\rm mom}=\frac{1}{\Rmom}.
\label{p3:eq:bare_cutoff_radius}
\end{equation}
The lowest retained spinorial shell of the physical seed has the shifted Dirac comparison ceiling
\begin{equation}
 \Lambda_D(1)=\frac{2}{\Rmom},
\end{equation}
whereas the profile regulator \(\Lcut\) entering the heat moments is an
independent spectral-action scale.
\end{proposition}

\begin{proof}
The momentum-space radius \(\Rmom\) sets the geometric momentum unit
\[
  \Punit:=\hbar/\Rmom .
\]
Hence the associated wave-number scale is
\[
  \Lambda_{\rm mom}=\frac{\Punit}{\hbar}=\frac{1}{\Rmom}.
\]
This is the geometric momentum scale inherited from the round momentum parent.
Theorem~\ref{p3:thm:spinorial-shell-ceiling-dictionary} separately gives
\(\Lambda_D(1)=2/\Rmom\) from the shifted Dirac counting operator.  Neither
identity fixes the profile variable \(\Lcut\): that variable occurs only in
the common functional-calculus profile of the spectral action.
\end{proof}

\begin{remark}[Forward Einstein--Hilbert handoff owned by Paper~III-E]
\label{p3:rem:forward-EH-cutoff-handoff}
No Einstein--Hilbert coefficient is proved in Paper~III-B.  After importing
Paper~III-E, Definition~\ref{p3:def:f2_rescaled} and
Proposition~\ref{p3:prop:EH_from_a2} give the dependency-labelled forward
handoff
\begin{equation}
 \boxed{U:=\Lambda_{\rm EH}^2:=\frac{f_2}{6}\,\Lcut^2},
\label{p3:eq:Lambda_EH_def_forward}
\end{equation}
where \(f_2\) is the rescaled second spectral moment.  This generates the
normalized EH curvature coefficient but does not create a second independent
Planck-curvature owner. Paper~III-E separately evaluates the actual
momentum/EH comparison $U/(\rho_H\mathcal K_p)$, with
$\mathcal K_p=6/\Rmom^2$ and $\rho_H=3/\pi$. Its vanishing logarithmic
residual gives $U=18/(\pi\Rmom^2)$; the scalarizer alone does not prove
that action/geometry port. The scalar-resolution coordinate is
\begin{equation}
 \boxed{G_{\rm scal}:=\frac{c_{\rm op}^3}{2\pi\hbar U}},
\label{p3:eq:Rp_EH_normalization}
\end{equation}
whereas the raw joint coordinate is $G_{J,\rm raw}=G_{\rm scal}/96$.
A final physical Newton coupling uses the complete matched action/source
normalization. Only after a successful momentum/EH port does the scalar
coordinate additionally read
\[
 G_{\rm scal}=\frac{c_{\rm op}^3\Rmom^2}{36\hbar}.
\]  The geometric momentum scale
\(\Lambda_{\rm mom}\), shifted shell ceiling \(\Lambda_D(m)\), and
spectral-action regulator \(\Lcut\) retain distinct defining roles; in
particular no identity \(\Lcut=1/\Rmom\) or \(2/\Rmom\) is inferred.  When
Paper~V-C proves that the vacuum photon has the same undeformed principal owner,
its theoretical front speed satisfies \(c_{\rm light}^{\rm th}=c_{\rm op}\);
empirical signal identification is separate.
\end{remark}


\subsubsection{Scalar projection and protected Clifford structure}
\label{p3:sec:factor_one_eighth}

\subparagraph{Protected-Clifford hypothesis boundary.}
All statements in this subsection that physically identify the normalized
protected operator channels with a real Clifford algebra require only
\((\mathrm{ND}_{\rm sp})\), as in
Theorem~\ref{p3:thm:8mode_operator_basis}.  The stronger representative
\(\mathsf C_{\rm PI}\) is retained elsewhere only where absolute completion
coefficients or coordinate scales are used.  Purely algebraic statements
about \(\Cl(0,3)\) remain valid independently of physicalization.

The protected Clifford block supplies a finite scalar projection, a normalized real-channel readout, and the corresponding normalized complex protected Hilbert trace.  Nonlinear hierarchy formulas that require later index-locking or active-density input are forward synthesis and are deferred to Paper~III-G, Section~\ref{p3b:app:forward-synthesis}.

\begin{lemma}[Regular Clifford trace extracts the scalar coefficient]
\label{p3:lem:clifford_scalar_projection}
Let $\Cl(0,3)$ be the real Clifford algebra generated by
$e_i e_j + e_j e_i = -2\delta_{ij}I$.  Let
$L_a:\Cl(0,3)\to\Cl(0,3)$ denote left multiplication by
$a\in\Cl(0,3)$ in the real regular representation.  Writing
$\Tr_{\mathbb R}$ for the ordinary real trace,
\[
\frac{1}{8}\,\Tr_{\mathbb R}(L_a)=\langle a\rangle_0,
\]
where $\langle a\rangle_0$ is the scalar (grade-$0$) part of $a$.
\end{lemma}

\begin{proof}
Fix the standard real basis
\[
\{1,\ e_i,\ e_ie_j\ (i<j),\ e_1e_2e_3\},
\]
so $\dim_{\mathbb R}\Cl(0,3)=8$.  Write $a=a_0+a_\perp$.  Left
multiplication by $a_0$ contributes $8a_0$ to the real trace.  Left
multiplication by every non-scalar basis element permutes the basis with sign
changes and has zero diagonal.  Hence
$\Tr_{\mathbb R}(L_a)=8a_0$.
\end{proof}

\begin{theorem}[Normalized real Clifford-channel trace and complex protected trace]
\label{p3:thm:normalized-protected-octet-trace}
Assume \((\mathrm{ND}_{\rm sp})\) and the regular-module identification of
Theorem~\ref{p3:thm:regular-module-two-eights}.  Let
\[
 \mathcal K_{\rm Cl}:=\Cl(0,3)_{\rm reg}
\]
be the eight-real-dimensional regular Clifford module, and let
\(\mathcal S_{\rm prot}\) be the corresponding four-complex-dimensional
physical protected module.  Define the normalized real-channel readout
\begin{equation}
 \tau_8(A):=\frac18\Tr_{\mathbb R}(A),
 \qquad A\in\End_{\mathbb R}(\mathcal K_{\rm Cl}),
 \label{p3:eq:normalized-octet-trace}
\end{equation}
and the normalized complex Hilbert trace
\[
 \tau_{\rm prot}^{\mathbb C}(B)
 :=\frac14\Tr_{\mathbb C}(B),
 \qquad B\in\End_{\mathbb C}(\mathcal S_{\rm prot}).
\]
On the represented Clifford algebra the two normalizations agree on scalar
extraction:
\[
 \tau_8(L_a)
 =\tau_{\rm prot}^{\mathbb C}(\rho(a))
 =\langle a\rangle_0,
 \qquad a\in\Cl(0,3).
\]
The real rank-one scalar-line projector
\(P_0:\mathcal K_{\rm Cl}\to\mathbb R\mathbf1\) obeys
\(\tau_8(P_0)=1/8\).  This \(1/8\) is a real Clifford-channel dimension
fraction; it is not the normalized weight of a rank-one projector on the
four-complex-dimensional protected Hilbert module, which would be \(1/4\).
The complex trace \(\tau_{\rm prot}^{\mathbb C}\), not \(\tau_8\), is the
ordinary finite-Hilbert trace normalization compatible with \(\Dcl\).
\end{theorem}

\begin{proof}
The real regular-module identity is
Lemma~\ref{p3:lem:clifford_scalar_projection}.  Under the real module
isomorphism of Theorem~\ref{p3:thm:regular-module-two-eights}, the underlying
real trace of a complex-linear operator satisfies
\[
 \Tr_{\mathbb R}(B_{\mathbb R})
 =2\,\Re\Tr_{\mathbb C}(B).
\]
For the represented real Clifford algebra the normalized scalar trace is real,
so division by \(8\) on the underlying real module equals division by \(4\)
on the complex protected module.  Full orthogonal invariance on the real
regular module gives the unique normalized real matrix trace, while Paper~II,
Theorem~\ref{p2:thm:no-insertion-finite-trace-weight}, gives the normalized
complex unitary trace on the physical complex Hilbert module under the
tracial/full-unitary hypothesis used by \(\Dcl\).  Finally, a real rank-one
projector has real trace one, giving \(\tau_8(P_0)=1/8\).
\end{proof}

\begin{definition}[Ambient and intrinsic real scalar-channel readouts]
\label{p3:def:scalar-corner-readout}
Let \(P_0\) be the rank-one real orthogonal projector onto
\(\mathbb R\mathbf1\subset\Cl(0,3)_{\rm reg}\), and let \(T\) be a
trace-class geometric operator.  The regular and corner \emph{real-channel}
readouts of a scalar coefficient \(a_0\) are, respectively,
\[
 (\Tr_{\rm geom}\otimes\tau_8)
   (T\otimes L_{a_0\mathbf1})
 =a_0\Tr_{\rm geom}(T),
\]
and
\[
 (\Tr_{\rm geom}\otimes\tau_8)
   (T\otimes a_0P_0)
 =\frac{a_0}{8}\Tr_{\rm geom}(T).
\]
The first line is the intrinsic full-support regular scalar readout.  The
second is the ambient rank-one scalar-constituent readout.  They are not two
normalizations of one untyped operation: their support ranks and generated
provenance differ.  The tag \((\mathrm{CORN})\) denotes the theorem status that
the second readout is the scalar constituent of an occurring protected-octet
composite as specified below; it is not an independent choice between the two
displayed formulas.  Neither formula is an ordinary rank-one projector weight
on the four-complex-dimensional protected Hilbert module.
\end{definition}

\begin{theorem}[Co-realized protected scalar constituent]
\label{p3:thm:derived-scalar-corner-readout}
Assume \((\mathrm{ND}_{\rm sp})\), and let \(P\) be a G0-occurring
protected-octet composite whose frozen generated presentation contains the
canonical scalar projector \(P_0\) of
Lemma~\ref{p3:lem:o3-clifford-expectation}.  Suppose that the presentation is
formed before the scalar constituent is evaluated and retains the canonical
split
\[
 \mathcal K_{\rm Cl}=P_0\mathcal K_{\rm Cl}
 \oplus P_W\mathcal K_{\rm Cl},
 \qquad P_W=1-P_0,
\]
and evaluates its scalar constituent with the ambient normalized real trace
\(\tau_8=\Tr_{\mathbb R}/8\), without renormalizing after restriction to
\(P_0\mathcal K_{\rm Cl}\).  Then the scalar and complementary constituents
are simultaneously realized, and the scalar constituent has the forced
coefficient
\[
 \boxed{\tau_8(P_0)=\frac18.}
\]
Its signature is rank one, \(O(3)\)-fixed, scalar, Born-even, and independent
of the running protected-complement scale.  The full-support regular readout
has support rank eight and is therefore not a competing identification of the
same action constituent.  Under these antecedents \((\mathrm{CORN})\) is a
derived certification tag, not an additional readout assumption.
\end{theorem}

\begin{proof}
Proposition~\ref{p3:prop:reeb-selection} identifies the scalar line and its
orthogonal complement inside the already realized protected octet.  The two
projector-supported pieces form a finite nonredundant generated-constituent
presentation of \(P\).  Paper~II-A
Theorem~\ref{p2:thm:generated-constituent-corealization} therefore realizes
both pieces in their native support sorts; no selection between them remains.
Theorem~\ref{p3:thm:normalized-protected-octet-trace} gives
\(\tau_8(P_0)=\operatorname{rank}_{\mathbb R}(P_0)/8=1/8\).

The scalar line is the unique \(O(3)\)-fixed rank-one summand of the displayed
split.  Hence its invariant signature separates it from the rank-seven
complement and from the rank-eight left-multiplication readout.  Paper~II-A
Corollary~\ref{p2:cor:canonical-constituent-identification} makes the
identification canonical.  The no-renormalization clause is essential: tracing
intrinsically on the already isolated one-dimensional line would instead give
one and would describe the different full-support/intrinsic readout of the
definition.
\end{proof}

\begin{corollary}[Selector status of the protected scalar/complement split]
\label{p3b:cor:protected-split-corealization}
Under the antecedents of
Theorem~\ref{p3:thm:derived-scalar-corner-readout},
\[
 \boxed{
 \operatorname{CoReal}
 \left(
 P_0\mathcal K_{\rm Cl},
 P_W\mathcal K_{\rm Cl}
 \mid P
 \right).}
\]
Thus the scalar and complementary protected constituents are not competing
realizations and no selector between them is permitted.  The ambient
coefficient \(\tau_8(P_0)=1/8\) is inherited from the generated constituent
presentation; an intrinsic renormalization after isolating the scalar line is
a different readout operation.
\end{corollary}

\subparagraph{\texorpdfstring{Linear real-channel term ($1/8$).}{Linear real-channel term (1/8).}}
When the forward handoff uses an occurring protected-octet composite with the
ambient constituent presentation of
Theorem~\ref{p3:thm:derived-scalar-corner-readout}, the scalar term carries
the derived factor \(1/8\).  If the scalar remnant has already been isolated
and is represented instead by left multiplication by \(a_0\mathbf1\), the
real-channel factor is \(1\), not \(1/8\).  The forward Einstein--Hilbert and
gauge handoffs must preserve the same ambient generated-constituent
presentation.  The value \(1/8\) follows from GP2 co-realization only after
that occurring presentation and the protected split have been generated; it
does not follow from the unmarked \(\Dcl\) loop alone.

\begin{proposition}[Protected scalar-channel counting]
\label{p3:prop:kinematic-sector-weight}
Assume \((\mathrm{ND}_{\rm sp})\).  On the derived momentum seed \(m_p=1\),
the protected operator sector contains eight real channels, while the
canonical scalar projector of
Proposition~\ref{p3:prop:reeb-selection} has rank one.  Consequently, whenever
the generated active orbit family terminates at this common projector, the number of
protected microscopic real channels feeding the surviving scalar channel is
\[
\nu_{\mathrm{spinor}}=1,\qquad \nu_{\mathrm{scalar}}=8,
\]
that is,
\[
\nu_{\mathrm{scalar}}:\nu_{\mathrm{spinor}}=8:1.
\]
\end{proposition}

\begin{proof}
Theorem~\ref{p3:thm:8mode_operator_basis} gives the eight-real-channel
protected Clifford algebra on \((\mathrm{ND}_{\rm sp})\).
Proposition~\ref{p3:prop:reeb-selection} identifies its canonical scalar
projector and proves that its image is one-dimensional.  The ratio is
therefore \(8:1\).  Paper~III-C derives that this projector is the common
terminal active endpoint, so no separate infrared selector or localization template is
required.
\end{proof}

\begin{definition}[Octet scalar-channel dimension fraction]
\label{p3:def:octet-scalar-projection-weight}
The protected Clifford octet carries the normalized scalar-channel dimension
fraction
\begin{equation}
\label{p3:eq:octet-scalar-projection-weight}
\wp_8:=\frac{1}{\dim_{\RR}\Cl(0,3)}=\frac18.
\end{equation}
This symbol is first a dimension fraction, not yet a physical trace weight and
not a protected-module projector usually denoted by an octet projection
operator in later papers.  Theorem~\ref{p3:thm:normalized-protected-octet-trace}
identifies it with \(\tau_8(P_0)\); it becomes the physical scalar-channel
weight when the antecedents of
Theorem~\ref{p3:thm:derived-scalar-corner-readout} hold.  The tag
\((\mathrm{CORN})\) records that derived status.
\end{definition}

\begin{theorem}[Regular-Clifford scalar-channel dimension rule]
\label{p3:thm:octet-scalar-projection-rule}
Assume \((\mathrm{ND}_{\rm sp})\).  Let $A_{\rm Cl}\cong\Cl(0,3)$ be the protected real Clifford block.  The
normalized scalar-channel dimension fraction is
\[
w_0=\frac{\dim_{\mathbb R}\mathbb R\mathbf 1}{\dim_{\mathbb R}A_{\rm Cl}}=\frac18.
\]
For $n$ independent protected Clifford factors,
$A_{\rm Cl}^{\otimes n}$, the pure scalar-channel dimension fraction is
\[
w_0^{(n)}=\left(\frac18\right)^n=8^{-n},
\]
and the total protected octet multiplicity is
\[
\dim_{\mathbb R}A_{\rm Cl}^{\otimes n}=8^n.
\]
\end{theorem}

\begin{proof}
As a real vector space, $\Cl(0,3)$ has basis
\[
1,e_1,e_2,e_3,e_{12},e_{23},e_{31},e_{123},
\]
so $\dim_{\mathbb R}\Cl(0,3)=8$.  The scalar line $\mathbb R\mathbf 1$ has dimension $1$, hence the normalized scalar-channel fraction is $1/8$.  Taking $n$ independent tensor factors multiplies dimensions: the full block has dimension $8^n$, while the pure scalar tensor channel remains one-dimensional.  Its normalized weight is therefore $8^{-n}$.
\end{proof}

Under the occurring ambient protected-constituent presentation certified by
\((\mathrm{CORN})\),
Theorem~\ref{p3:thm:normalized-protected-octet-trace} identifies these
dimension fractions with the corresponding real protected-channel readout
weights.

\begin{remark}[Octet measure versus Riemannian volume]
The factor $1/8$ in
Theorem~\ref{p3:thm:octet-scalar-projection-rule} is first a
Clifford scalar-channel dimension fraction.  It becomes the protected
scalar-channel coefficient on the occurring ambient constituent branch of
Theorem~\ref{p3:thm:derived-scalar-corner-readout}; \((\mathrm{CORN})\) is
the certification tag for that theorem domain.
It is not, by itself, an ordinary Riemannian volume ratio on $S^3$. Any
appearance of $1/8$ as a compact volume fraction should therefore be
understood as an octet-normalized cell measure induced by the protected
Clifford decomposition unless an independent Riemannian volume computation
is supplied.
\end{remark}

\begin{theorem}[Explicit regular-module reconciliation of the protected octet]
\label{p3:thm:regular-module-two-eights}
Assume \((\mathrm{ND}_{\rm sp})\) and use the imported protected even Clifford algebra \(A_{\rm ev}\cong\Cl(0,3)\).  Choose the following explicit equivalent real module realization.  Let
\[
\mathcal S_{\rm prot}=\mathcal S_+\oplus\mathcal S_-
\]
be the protected chiral/diagonal-weight spinor carrier, with
\(\mathcal S_s\cong\mathbb C^2_{\mathbb R}\) and
\(\dim_{\mathbb R}\mathcal S_s=4\) for \(s=\pm1\).  Work in the convention
\[
 e_i e_j+e_j e_i=-2\delta_{ij},
 \qquad \omega=e_1e_2e_3,
 \qquad p_s=\frac12(1+s\omega).
\]
On \(\mathcal S_s\) define the real-linear Clifford action
\[
 \Gamma_i^{(s)}:=s\,i\sigma_i .
\]
Then
\[
 \Gamma_i^{(s)}\Gamma_j^{(s)}+\Gamma_j^{(s)}\Gamma_i^{(s)}
 =-2\delta_{ij}I,
 \qquad
 \Gamma_1^{(s)}\Gamma_2^{(s)}\Gamma_3^{(s)}=sI.
\]
For \(v_s=(1,0)^T\), the map
\begin{equation}
\label{p3:eq:explicit-clifford-intertwiner}
 \Phi_s:\Cl(0,3)p_s\longrightarrow\mathcal S_s,
 \qquad
 \Phi_s(ap_s):=\rho_s(a)v_s,
\end{equation}
is an explicit real left-\(\Cl(0,3)\)-module isomorphism.  Hence
\begin{equation}
\label{p3:eq:protected-carrier-regular-clifford-module}
 \mathcal S_{\rm prot}
 \cong
 \Cl(0,3)p_+\oplus\Cl(0,3)p_-
 \cong
 \Cl(0,3)_{\rm reg}.
\end{equation}
Let \(\mathcal S_{{\rm prot},\mathbb R}\) denote the underlying real
module.  Consequently, for every \(a\in\Cl(0,3)\),
\begin{equation}
\label{p3:eq:protected-carrier-trace-scalar-projection}
 \frac18\Tr_{\mathbb R}\!\bigl(\rho_{\mathbb R}(a)\bigr)
 =\frac14\Tr_{\mathbb C}\!\bigl(\rho(a)\bigr)
 =\frac18\Tr_{\mathbb R}(L_a)
 =\langle a\rangle_0.
\end{equation}
The first trace is taken on \(\mathcal S_{{\rm prot},\mathbb R}\), the
second on the four-complex-dimensional \(\mathcal S_{\rm prot}\), and the
third on \(\Cl(0,3)_{\rm reg}\).  The module identification is unique up to
the quaternionic commutant on each summand and therefore leaves these
normalized scalar projections unchanged.
\end{theorem}

\begin{proof}
The central volume element satisfies \(\omega^2=1\), so the \(p_s\) are
orthogonal central idempotents and
\(\Cl(0,3)=\Cl(0,3)p_+\oplus\Cl(0,3)p_-\).  Each ideal has real basis
\[
 p_s,\quad e_1p_s,\quad e_2p_s,\quad e_3p_s.
\]
The Pauli relations give the displayed Clifford identities.  Moreover,
\[
\begin{aligned}
 \Phi_s(p_s)&=(1,0)^T,\\
 \Phi_s(e_1p_s)&=(0,si)^T,\\
 \Phi_s(e_2p_s)&=(0,-s)^T,\\
 \Phi_s(e_3p_s)&=(si,0)^T.
\end{aligned}
\]
These four vectors are linearly independent over \(\mathbb R\), so \(\Phi_s\)
is bijective.  For \(a,b\in\Cl(0,3)\),
\[
 \Phi_s\bigl(b(ap_s)\bigr)=\rho_s(b)\Phi_s(ap_s),
\]
which proves equivariance.  Taking the direct sum gives
\eqref{p3:eq:protected-carrier-regular-clifford-module}; trace invariance and
Lemma~\ref{p3:lem:clifford_scalar_projection} give
\eqref{p3:eq:protected-carrier-trace-scalar-projection}.
\end{proof}

\begin{remark}[Chiral grading versus geometric Hopf charge]
The theorem identifies the protected chiral/diagonal-weight sectors with the two simple Clifford ideals.  Paper~II treats the diagonal-spin grading and the
geometric left-Hopf charge as distinct operators.  Therefore no identification
of the two Hopf-charge sectors with \(p_\pm\) is used here unless a separate
Hopf-charge intertwining theorem is imposed.
\end{remark}

\begin{remark}[Reconciled roles of the two eights]
\label{p3:rem:two-eights}
The normalization \(1/8\) in the regular trace, the corner weight
\(\tau_8(P_0)=1/8\), and the ratio \(8:1\) of
Proposition~\ref{p3:prop:kinematic-sector-weight} should not be conflated
with the complex algebra \(\End_{\CC}(S_+)\).  The first extracts a scalar
coefficient and returns \(a_0\) on \(L_{a_0\mathbf1}\); the second is the
co-realized rank-one constituent certified by \((\mathrm{CORN})\); the third counts protected channels
feeding the infrared scalar remnant.  Theorem~\ref{p3:thm:regular-module-two-eights}
shows that all three use the same regular Clifford octet without identifying
their operations.
\end{remark}


\begin{remark}[Supplementary Bott stability and resolution data]
The Bott-stable protected-octet theorem and its compatibility with Paper~I-A's
scalar-free \(56\)-dimensional resolution-data module are recorded in
Paper~III-G, Section~\ref{p3b:app:resolution-data-import}.  Neither result is used in the
native III-B action proof.
\end{remark}

\begin{lemma}[Unoriented Clifford averaging leaves only the scalar]
\label{p3:lem:o3-clifford-expectation}
Let $V\cong\RR^3$ be the real generator space of the protected Clifford algebra and identify
\[
\Cl(0,3)\cong \bigwedge\nolimits^* V
\]
as a real $\mathrm O(V)$-module by the Chevalley map.  Let
\[
\mathbb E_{\mathrm O}(a):=\int_{\mathrm O(V)} \rho(g)a\,d\mu(g)
\]
be Haar averaging over the full orthogonal group, with $\rho$ the induced action on
$\Cl(0,3)$ and $\mu$ normalized Haar measure.  Then
\begin{equation}
\label{p3:eq:o3-expectation-scalar}
\mathbb E_{\mathrm O}(a)=\langle a\rangle_0\,\mathbf 1,
\qquad
\Cl(0,3)^{\mathrm O(3)}=\RR\mathbf 1.
\end{equation}
Consequently the kernel of the infrared scalar averaging map has real dimension seven.
\end{lemma}

\begin{proof}
Under the Chevalley identification,
\[
\Cl(0,3)\cong \bigwedge^0 V\oplus\bigwedge^1 V\oplus\bigwedge^2 V\oplus\bigwedge^3 V .
\]
The $\mathrm O(V)$-fixed subspace of $\bigwedge^0 V$ is $\RR$.  The fixed subspace of
$\bigwedge^1V=V$ is zero, since no nonzero vector is fixed by all rotations.  In three
dimensions $\bigwedge^2V\cong \det(V)\otimes V$, so a fixed bivector would again determine
a fixed vector up to the determinant character; none exists.  Finally
$\bigwedge^3V\cong\det(V)$, and every orientation-reversing element of $\mathrm O(V)$ acts
on this line by $-1$, so its fixed subspace is also zero.  Thus the only fixed subspace is
$\bigwedge^0V=\RR\mathbf 1$.  Haar averaging over a compact group is the canonical projection
onto the fixed subspace, proving \eqref{p3:eq:o3-expectation-scalar}.  Since
$\dim_{\RR}\Cl(0,3)=8$, the kernel has dimension $8-1=7$.
\end{proof}

\begin{remark}[Parity-even scalar branch: \texorpdfstring{$\mathrm O(3)$}{O(3)} versus \texorpdfstring{$\SO(3)$}{SO(3)}]
\label{p3:rem:o3-not-so3}
Averaging only over the orientation-preserving rotation group leaves
\[
 \Cl(0,3)^{\SO(3)}
 =\mathbb R\mathbf1\oplus\mathbb R(e_1e_2e_3).
\]
Thus the pseudoscalar line is an oriented singlet.  The characterization
below shows that parity evenness eliminates this only second singlet.  In
Paper~III-C the rank-one terminal active projector supplies that elimination
dynamically: parity is certified by the terminal rank rather than assumed to
choose a map.  The resulting map is the full \(O(3)\)-equivariant scalar
projection.  Orientation reversal sends the pseudoscalar to its negative, and
\[
 \Cl(0,3)^{\mathrm O(3)}=\mathbb R\mathbf1.
\]
Hence the rank drop \(8\to1\), and therefore the reduced rank seven, is exact
on the parity-even infrared scalar branch.  On an orientation-preserving
\(\SO(3)\)-only counterfactual branch the invariant remnant would be
two-dimensional and the rank drop would be six.  This counterfactual is not
the native active-depth-six question: Paper~III-C,
Theorem~\ref{p3c:thm:active-depth-six-obstruction}, proves that the actual
active upper-support chain jumps from rank three at \(q=5\) to the scalar
rank-one endpoint at \(q=7\).  Its only equivariant rank-two subprojector is
the rotating plane \(u^\perp\), which excludes the scalar line and therefore
is not an intervening active carrier.
\end{remark}

\begin{proposition}[Canonical parity-even scalar projector]
\label{p3:prop:reeb-selection}
Assume \((\mathrm{ND}_{\rm sp})\).  A unital linear map
\(E:\Cl(0,3)\to\mathbb R\mathbf1\) that is \(SO(3)\)-equivariant and
parity-even is necessarily unique.  It is
\[
 P_0(a):=\langle a\rangle_0\mathbf1.
\]
Equivalently, \(P_0\) is the orthogonal projector onto
\(\mathbb R\mathbf1\), and it equals Haar averaging:
\begin{equation}
\label{p3:eq:IR-clifford-expectation}
 P_0(a)
 =\int_{\mathrm O(3)}g\cdot a\,dg
 =\langle a\rangle_0\mathbf1.
\end{equation}
Its fixed space is \(\mathbb R\mathbf1\), and its kernel is the
seven-dimensional non-scalar complement \(W\).  It is unital, idempotent,
self-adjoint for the normalized regular Clifford inner product, parity-even,
and fully \(O(3)\)-equivariant.  The label
\((\mathrm{IR}_{O(3)})\) certifies that the occurring physical
active orbit family has the common terminal projector \(P_0\); it is not an
independent extraction hypothesis.  Paper~III-C proves that orbit-invariant
terminal-generation theorem.
\end{proposition}

\begin{proof}
As an \(SO(3)\)-module, \(\Cl(0,3)\cong
\bigwedge^0V\oplus\bigwedge^1V\oplus\bigwedge^2V\oplus\bigwedge^3V\)
has exactly two trivial summands: the scalar and pseudoscalar lines.  An
\(SO(3)\)-equivariant map to the scalar line vanishes on the vector and
bivector summands.  Parity fixes the scalar and negates the pseudoscalar, so
parity evenness forces the latter coefficient to vanish.  Unitality fixes the
remaining scalar coefficient to one.  Hence
\(E(a)=\langle a\rangle_0\mathbf1\), which is precisely the
full orthogonal Haar average of
Lemma~\ref{p3:lem:o3-clifford-expectation}.  Since
\(\dim_{\mathbb R}\Cl(0,3)=8\), its kernel has real dimension seven.
Orthogonal Haar averaging is self-adjoint and idempotent, and its range is the
scalar line, proving the remaining assertions.
\end{proof}

\begin{remark}[Selector chronology of the infrared projector]
\label{p3b:rem:infrared-derived-later}
At the III-B stage, Proposition~\ref{p3:prop:reeb-selection} proves that
\(P_0\) is the unique admissible unital parity-even scalar map, but it does
not manufacture the physical occurrence of an active response terminating at
that map.  Paper~III-C independently constitutes the equivariant orbit family
of finite active responses and proves that their common terminal spectral
projector is exactly \(P_0\).
Accordingly the physical endpoint has the chronological selector status
\[
 \boxed{\begin{gathered}
 \text{III-B: unique candidate map; occurrence unresolved}\\
 \longrightarrow\ \text{III-C: }P_{\rm term}=P_0
 \quad(\mathrm{DERIVED\text{-}LATER}).
 \end{gathered}}
\]
The later result does not alter the stage-relative III-B occurrence status:
the active endpoint is not claimed to have been generated at this stage.
\end{remark}

\begin{remark}[Geometric support from collapsing Hopf fibrations]
The endpoint statement above is representation-theoretic: it identifies the
only possible parity-even scalar endpoint.  The independent geometric picture
is the Dirac spectrum of collapsing $S^1$-bundles: for projectable spin
structures, low-lying eigenvalues can survive the collapse, whereas
nonprojectable sectors are driven to $\pm\infty$ \cite{Ammann1998}.  This
supports shellwise decoupling of Reeb-charged content.  It is not used as an
occurrence postulate: Paper~III-C constructs an equivariant orbit family of
finite self-adjoint active responses whose common terminal spectral projector is exactly
\(P_0:\Cl(0,3)\to\RR\mathbf1\).
\end{remark}
\begin{remark}[Forward active-density synthesis]
Scale-dependent active-channel entropy and the nonlinear
\(8^{n(\scal)}\)-type response require a coarse-graining density not derived
by \(\Dcl\) alone.  They are therefore recorded, with their endpoint
identities, in Paper~III-G, Section~\ref{p3b:app:forward-synthesis}.  Paper~III-C derives
the equivariant orbit of discrete active-support chains and the orbit-invariant
scalar infrared endpoint.  A particular rank-three support requires the
marked diagonal and does not descend from the unmarked orbit.  At the III-B
stage a general scale-dependent density profile remains under the explicit
forward datum \((\mathrm{ACT})\); the later occurring \((\mathrm{MRESP})\)
composite in Paper~III-E independently discharges one particular member of
that family, without making the whole generic density atlas unique.
\end{remark}

\subsubsection{Heat Kernel Expansion and finite/continuum separation}

By Definition~\ref{p3:def:physical_spectral_action}, the physical action is
an exact finite sectorwise trace.  The momentum term is always evaluated on
the fixed physical seed \(\mathcal H_1\); no physical momentum large-window
limit is taken.

\begin{theorem}[Fixed-time exhaustion and the controlled joint heat regime]
\label{p3:thm:directed-spectral-window-exhaustion}
Let \(L\ge0\) be a self-adjoint operator on \(\mathcal H\) such that
\(e^{-tL}\) is trace class for every \(t>0\), and let \(P_m\) be increasing
spectral projections converging strongly to the identity.  Set
\[
        L_m=P_mLP_m\big|_{P_m\mathcal H}.
\]
Then for every \(t>0\),
\[
        \Tr_{P_m\mathcal H}(e^{-tL_m})
        \longrightarrow
        \Tr_{\mathcal H}(e^{-tL})
        \qquad(m\to\infty).
\]
This fixed-\(t\) limit does not imply convergence of strict finite-window
small-\(t\) coefficients.

Assume in addition that for some \(C,q,t_0>0\),
\[
 \Tr(e^{-tL})\le Ct^{-q}\qquad(0<t\le t_0),
\]
and let \(\mu_m\) be the least spectral value of \(L\) omitted by \(P_m\).
If
\(t_m\downarrow0\), \(m\to\infty\), and
\[
 \frac{t_m\mu_m}{|\log t_m|}\longrightarrow\infty,
\]
then for every fixed \(N>0\),
\[
 \Tr_{\mathcal H}(e^{-t_mL})
 -\Tr_{P_m\mathcal H}(e^{-t_mL_m})
 =o(t_m^N).
\]
If \(L\) also has a small-time heat expansion, any prescribed finite part of
that expansion is therefore reproduced along this
intermediate-time/large-window scaling.  The Dirac application is
\(L=D^2\); a first-order Dirac operator itself is not called Laplace type.
This statement
does not require the physical momentum label \(m_p=1\) to run; the
position-side continuum parent and its large-window presentation are
logically distinct from the fixed microscopic momentum readout.
\end{theorem}

\begin{proof}
Because \(P_m\) is a spectral projection of \(L\),
\([P_m,L]=0\) and
\[
        e^{-tL_m}=P_me^{-tL}P_m
\]
on \(P_m\mathcal H\).  The projections increase strongly to the identity and
\(e^{-tL}\) is positive trace class.  Monotone
convergence for positive trace-class operators gives the displayed limit.
It cannot be promoted to a coefficient limit by exchanging limits.  For
example,
\[
 K_m(t)=\sum_{|k|\le m}e^{-tk^2}
\]
obeys \(K_m(t)\to2m+1\) as \(t\downarrow0\) for fixed \(m\), whereas
the circle parent has \(K_\infty(t)\sim\sqrt{\pi/t}\).

For the positive joint result, every omitted spectral value \(\lambda\) obeys
\(\lambda\ge\mu_m\).  Splitting the exponential gives the exact estimate
\[
 0\le
 \Tr(e^{-tL})-\Tr(P_me^{-tL}P_m)
 \le e^{-t\mu_m/2}\Tr(e^{-(t/2)L})
 \le C2^q t^{-q}e^{-t\mu_m/2}
\]
for \(0<t\le2t_0\).
The joint hypothesis makes the right-hand side \(o(t_m^N)\) for every
fixed \(N\).  The parent Seeley--DeWitt expansion may therefore be truncated
at any prescribed finite order whenever that expansion is part of the stated
operator domain, while the omitted spectral tail is smaller
than that order.  The coefficients remain those of the continuum parent, not
strict coefficients of one fixed matrix.
\end{proof}

\begin{theorem}[Finite momentum trace and controlled parent/position heat comparison]
\label{p3:thm:hk_sectorwise}
The physical momentum contribution remains the exact finite trace
\[
 S_p^{\rm fin}:=\Tr_{\mathcal H_1}f(D_p^2/\Lcut^2).
\]
If a position-side parent operator admits the large-window regime of
Theorem~\ref{p3:thm:directed-spectral-window-exhaustion}, then any prescribed
finite part of its parent heat expansion may be recovered along that joint
large-window/intermediate-time scaling.  No analogous physical
\(m_p\to\infty\) limit is used.

A four-dimensional \(S^3\times S^1\) Seeley--DeWitt expansion additionally
requires a scalar-circle spectral metric.  Since III-A does not constitute
such a metric, III-B does not identify those four-dimensional parent
coefficients with the literal physical action.  If a later paper supplies
that additional metric datum, the standard four-dimensional expansion may be
used there as a parent comparison theorem.
\end{theorem}

\begin{remark}
The native III-B action therefore separates exact finite momentum data from
any later continuum/large-window comparison.  This avoids assigning
continuum spectral dimension or Seeley coefficients to the fixed finite
matrix \(\mathcal H_1\).
\end{remark}

\begin{remark}[Deferred four-dimensional metric comparison]
If a later paper supplies a scalar-circle spectral metric, the corresponding
commutative $S^3\times S^1$ parent admits the standard local curvature
expansion.  III-B does not use those four-dimensional coefficients as the
literal action of the fixed finite momentum seed.
\end{remark}


\subsubsection{The Fuzzy Heat Kernel}

\begin{definition}[Classical scalar diagonal heat trace]
\begin{equation}
K_\infty^{\rm diag}(t) = \sum_{k=0}^{\infty} (k+1)^2 e^{-k(k+2)t/\Rmom^2}.
\end{equation}
\end{definition}

\begin{definition}[Full fuzzy heat trace]
For the full \(\operatorname{Sp}in(4)\)-equivariant truncation \(j_L+j_R\le m\), with \(j_L,j_R\in\tfrac12\mathbb N_0\), define
\begin{equation}
K_m^{\rm full}(t)
=
\sum_{\substack{j_L,j_R\in\frac12\mathbb N_0\\ j_L+j_R\le m}}
(2j_L+1)(2j_R+1)
\exp\!\left(
-\frac{2t}{\Rmom^2}\bigl[j_L(j_L+1)+j_R(j_R+1)\bigr]
\right).
\end{equation}
\end{definition}

\begin{definition}[Scalar diagonal fuzzy heat trace]
On the scalar diagonal sector \(j_L=j_R=k/2\), define
\begin{equation}
K_m^{\rm diag}(t)
=
\sum_{k=0}^{\lfloor m\rfloor}
(k+1)^2 e^{-k(k+2)t/\Rmom^2}.
\end{equation}
\end{definition}

Both fuzzy heat traces are \textbf{finite sums}---no divergences, no regularization needed.  The notation distinguishes the full auxiliary-module trace from the scalar diagonal trace used when only scalar harmonics are being counted.

\begin{remark}[Optional hopfion comparison]
The native III-B action and heat-window theorems do not require a dynamical
weighting of Hopf sectors.  Paper~III-G, Section~\ref{p3:sec:hopfions-spectral} constructs
a common-scheme comparison functional under explicit heat-expansion,
trace-class, lower-bound, scale, and normalization hypotheses.  It supplies no
sector occurrence or preference theorem.
\end{remark}

\subsection{The complete bare metric heat coefficient}
\label{p3b:sec:v2216-bare-metric-heat}

The ordinary joint heat and the scalar-resolution response keep their
distinct coefficient owners. The volume term is retained when the
smooth-parent metric coefficient is read as a local action; its removal
would change the complete metric variation.

\begin{corollary}[Bare vacuum and Einstein coefficients in common heat units]
\label{p3b:cor:v2216-bare-vacuum-ratio}
On the admitted four-dimensional Euclidean metric heat branch of the
ordinary joint Dirac action, retain the spin--internal rank
$4\cdot48=192$, the inherited pure-metric weight $1/8$, and the
native heat parameter $\Lambda_g$. Its first two pure-metric local
terms are
\[
 S_{J,\mathrm{met}}^E
 =\frac{1}{16\pi^2}\int\sqrt g\,
       (24\Lambda_g^4-2\Lambda_g^2R_E)
       +\text{retained higher terms}.
\]
On the corresponding licensed Lorentzian local real form, write the
two-derivative metric word as $\kappa_JR-V_{0,J}$. Then
\[
 \kappa_J=\frac{\Lambda_g^2}{8\pi^2}=12U,
 \qquad V_{0,J}=12\kappa_J\Lambda_g^2,
 \qquad \Lambda_g^2=96\pi^2U.
\]
An actually generated common positive constant action transport rescales
$\kappa_J$ and $V_{0,J}$ together and preserves $V_{0,J}/\kappa_J$.
It cannot match a different vacuum intercept by changing only the
Einstein coefficient. This coefficient identity selects no stationary
metric or renormalized effective vacuum.
\end{corollary}
\begin{proof}
For the admitted Dirac square $D^2=-\nabla^2+R_E/4$, the local
Lichnerowicz heat coefficient per spin--internal component is
$R_E/6-R_E/4=-R_E/12$. Multiplication by
$(1/8)\cdot192=24$ gives the displayed volume and curvature
coefficients. The native scalar coefficient is
$U=\Lambda_g^2/(96\pi^2)$; the existing finite local real-form
map gives the printed Lorentzian signs. Multiplying the complete
action by a constant multiplies both coefficients by that constant.
The higher heat terms, density, localizer and actual boundary retain
their existing domains and are not discarded by this comparison.
\end{proof}

The consequences for constant-field stationarity and normalized positive
proper-time profiles are proved on their explicit domains in
Paper~III-F, Theorem~\ref{p3f:thm:v2216-positive-profile-vacuum-bound}.
The bare coefficient is not identified with the physical scalar-resolution
Einstein coefficient, and the ratio $96$ supplies no action/source
projection by itself.


\subsection{Coordinate-only operator ancestry and spectral-action scope}
\label{p3b:sec:v226-spectral-action-ancestry}
Theorem~\ref{p1:thm:v226-coordinate-only} strengthens the intrinsic
operator input to the coordinate-generated algebra itself. The six symmetry
generators and chiral projectors are reconstructed by
Theorems~\ref{p1:thm:v226-rotations} and~\ref{p1:thm:v226-chirality}.
This removes their separate algebra-admission premise on this branch.

\begin{proposition}[The stronger coordinate algebra still does not select the square]
\label{p3b:prop:v226-square-scope}
On a supplied smooth common-symbol Dirac branch, adding a self-adjoint
zero-order term that commutes with scalar multiplication leaves the intrinsic
coordinate class unchanged, while it may change the Dirac square and its
spectral action. Complete coordinate-algebra generation therefore does not
discharge the complete action-contact certificate by itself.
\end{proposition}
\begin{proof}
For a scalar multiplier $M_f$ and such a term $V$,
$[D+V,M_f]=[D,M_f]$. But
$(D+V)^2=D^2+DV+VD+V^2$, whose additional terms need not vanish.
The intrinsic coordinate map depends on the first commutator, whereas the
spectral action consumes the actual square, domain and trace prescription.
The finite matrix counterpart permits separate native zero-order matrices
without changing the already specified coordinate tuple.
\end{proof}
For a supplied native block, the full operator and source-jet contact can
now be tested by 137 exact coefficients
(Theorem~\ref{p1b:thm:v226-dual-coefficients}). Those coefficients must
still be evaluated on the actual complete action and boundary owner.


\subsection{The dynamical centre quotient does not quotient the action}
\label{p3b:sec:v227-action-centre}
\begin{proposition}[Even dynamics and the heat action consume different data]
\label{p3b:prop:v227-action-centre}
For a block-reducing native Hamiltonian and actual central offset
$C=\sum_ac_aP_a$, even Heisenberg dynamics is unchanged at fixed state.
The unnormalized heat trace is instead
\[
 \operatorname{Tr}e^{-t(H+C)}
 =\sum_a e^{-tc_a}\operatorname{Tr}e^{-tH_a}.
\]
Thus equality in the even-dynamics quotient does not certify equality of
the heat action, its absolute normalization or its metric/source jets.
\end{proposition}
\begin{proof}
The commutator identity and blockwise functional calculus are
Theorem~\ref{p1b:thm:v227-central-readouts} and
Proposition~\ref{p1b:prop:v227-central-boundary}. Differentiating the
actual heat word retains the derivatives of its central coefficients,
projectors and all other source-dependent factors.
\end{proof}
This refines the preceding square-selection boundary: a complete frequency
response may close its own readout without selecting an unused dark
complement, but no corresponding cancellation is inferred inside the
complete trace or gravity comparison.


\section{Fundamental Born-meron spin--Hodge--Dirac and determinant readouts}
\label{p3b:sec:v349-born-meron-hodge-determinant}
Consider only the round \emph{relative} three-sphere $S^3(r)$,
with $\theta=g^{-1}dg$, the stationary Born connection
$A_B=\theta/2$ of Theorem~\ref{p3b:thm:v346-born-meron-ym},
and the fundamental complex rank-two twist.  This is a
specified mathematical test bundle, not an independently
actualized four-dimensional fermion sector.

\begin{theorem}[Fixed Born connection equals the round spin connection]
\label{p3b:thm:v349-meron-spin-hodge}
Choose $\gamma_a=i\sigma_a$, orthonormal left-invariant frame
$e_a$, and the round radius-$r$ metric for which
$[e_a,e_b]=-2\varepsilon_{abc}e_c/r$.  In this frame,
\[
 \Omega_{\mathrm{spin}}(e_a)=\frac{\gamma_a}{2r}
 =A_B(e_a).
\]
Via the invariant fundamental-dual identification,
$S\otimes S^*\cong\operatorname{End}(S)$; under the normalized
Clifford map $\Phi(f,\alpha)=2^{-1/2}(fI+\alpha_a\gamma_a)$,
\[
 \Phi^{-1}D_B\Phi=
 \begin{pmatrix}0&d^\dagger\\d&-*d\end{pmatrix},
 \qquad D_B^2\cong\Delta_0\oplus\Delta_1.
\]
The curl sign changes with the orientation convention, but the
square and unordered spectrum do not.
\end{theorem}
\begin{proof}
Koszul's rule gives $\nabla_{e_a}e_b=-\varepsilon_{abc}e_c/r$.
With $\gamma_b\gamma_c=-\delta_{bc}I-\varepsilon_{bcd}\gamma_d$,
the spin connection term
$\frac14\sum_{b,c}\langle\nabla_{e_a}e_b,e_c\rangle\gamma_b\gamma_c$
is $\gamma_a/(2r)$, equal to $\theta(e_a)/2$.
The twist is equivalent to $S^*$; the tensor-product
connection is therefore the endomorphism connection.
Clifford multiplication of $fI+\alpha_a\gamma_a$ gives
$-(\nabla_a\alpha_a)I+
(\nabla_cf-\varepsilon_{abc}\nabla_a\alpha_b)\gamma_c$,
which yields the displayed matrix and its Hodge-Laplacian square.
\end{proof}

\begin{corollary}[Complete fundamental twisted Dirac spectrum]
\label{p3b:cor:v349-meron-dirac-spectrum}
The complete spectrum of $D_B$ is
\[
 \operatorname{Spec}D_B=\{0^{(1)}\}
 \cup\bigcup_{n=2}^{\infty}
 \left\{\left(\pm\frac{\sqrt{n^2-1}}r\right)^{(n^2)},
 \left(\pm\frac nr\right)^{(n^2-1)}\right\},
\]
where each sign has the displayed multiplicity.  The kernel is
one-dimensional, represented by the identity endomorphism;
the first nonzero absolute eigenvalue is $\sqrt3/r$, and
$\operatorname{Spec}D_B^2\setminus\{0\}\subset[3/r^2,\infty)$.
The ordinary eta-invariant is zero and the reduced APS eta
invariant is $1/2$ modulo integers.  The squared heat trace is
\[
 H_B(t)=1+2\sum_{n=2}^\infty
 \left[n^2e^{-t(n^2-1)/r^2}
 +(n^2-1)e^{-tn^2/r^2}\right].
\]
\end{corollary}
\begin{proof}
Scalar round-sphere harmonics have eigenvalues
$(n^2-1)/r^2$ of multiplicity $n^2$.  Exact one-forms
replicate the nonconstant scalar levels; coexact one-forms
have $n^2/r^2$ with multiplicity $2(n^2-1)$, split into
opposite curl helicities.  The Hodge--Dirac matrix pairs the
nonconstant scalar and exact one-form levels with opposite
signs, and the coexact helicities supply the remaining pairs.
$H^0(S^3)=\mathbb C$ and $H^1(S^3)=0$ give the kernel.
Pair symmetry gives eta zero and a reduced eta of one half.
\end{proof}

\begin{theorem}[Exact positive-spectrum zeta determinant and all local heat coefficients]
\label{p3b:thm:v349-meron-zeta-heat}
With the single zero eigenvalue omitted,
\[
 \zeta_B(s)=2r^{2s}\sum_{n=2}^\infty
 \bigl[n^2(n^2-1)^{-s}+(n^2-1)n^{-2s}\bigr]
\]
has a meromorphic continuation regular at $s=0$ with
\[
 \boxed{\zeta_B(0)=-1,\quad
 \zeta_B'(0)=-2\log r+2\log2
 -\frac{2\zeta_{\mathrm R}(3)}{\pi^2}},
\]
\[
 \boxed{\det{}'(D_B^2)=\frac{r^2}{4}
 \exp\left(\frac{2\zeta_{\mathrm R}(3)}{\pi^2}\right)}.
\]
Writing $u=t/r^2$, its entire small-time algebraic expansion is
\[
 H_B(t)\sim\frac{\sqrt\pi\,r^3}{t^{3/2}}
 \left[1-\frac{t}{2r^2}
 +\sum_{k=2}^\infty\frac{t^k}{2k!r^{2k}}\right].
\]
The omitted terms are exponentially small as $t\downarrow0$,
up to powers of $t$.  The determinant of $D_B^2/\mu^2$ is
$(\mu r)^2\exp(2\zeta_{\mathrm R}(3)/\pi^2)/4$.
\end{theorem}
\begin{proof}
For $A(s)=\sum_{n\ge2}n^2(n^2-1)^{-s}$ use the binomial
continuation
$A(s)=\sum_{k\ge0}(s)_k[\zeta_{\mathrm R}(2s+2k-2)-1]/k!$.
For $B(s)=\zeta_{\mathrm R}(2s-2)-\zeta_{\mathrm R}(2s)$,
$A(0)=-1$, $B(0)=1/2$, and the absolutely convergent tail
in $A'(0)$ evaluates to $3/2-\log\pi$ by telescoping
finite products and Stirling's formula.  With
$\zeta_{\mathrm R}'(0)=-\log(2\pi)/2$ and
$\zeta_{\mathrm R}'(-2)=-\zeta_{\mathrm R}(3)/(4\pi^2)$
the asserted derivative follows.
The exact heat representation
$H_B(t)=-(1+e^u)\Theta'(u)-\Theta(u)$,
$\Theta(u)=\sum_{n\in\mathbb Z}e^{-un^2}$,
and Poisson's theta identity give the full asymptotic series.
The analytic continuation and all-order coefficient derivation
are detailed in Paper~III-G.
\end{proof}

\begin{remark}[Spectral-action and observational scope]
\label{p3b:rem:v349-meron-spectrum-physical-guard}
These are exact geometric test-operator results on $S^3(r)$.
The stationary Born meron is a Yang--Mills \emph{saddle}, not a
vacuum.  Neither the mathematical zero mode nor the positive
Dirac gap implies an observed massless particle, an interacting
four-dimensional gauge mass gap, or a predetermined physical
log-determinant coefficient in the common action.
\end{remark}


\section{Clifford ancestry and finite spectral exports}
\label{paper3b:sec:organized-exports}
\label{p3:part:appendices}
\subsection{Spin ancestry of the generated diagonal ratio fibre}
\label{p3b:sec:diagonal-spin-ancestry-update}

The identity ratio fibre is the generated graph
\[
 F_4=\{(\iota_x^{-1}(g),\iota_p^{-1}(g)):g\in K\}.
\]
On the vector $SU(2)$ factor define
\[
 \delta_p(\iota_x^{-1}(g),\iota_p^{-1}(g)):=\iota_p^{-1}(g).
\]
Its inverse is the graph inclusion
$b\mapsto(\iota_x^{-1}(\iota_p(b)),b)$, so no independent identification is
chosen.

\begin{theorem}[Orientation-preserving diagonal-fibre spin ancestry]
\label{p3b:thm:diagonal-fibre-spin-ancestry}
The map $\delta_p$ is an orientation-preserving round-sphere ancestry map up
to the positive homothety already induced by the product metric.  Its spin
lift preserves the two simple factors of $\operatorname{Sp}in(4)$ and therefore induces
\[
 S_L\overset\sim\longrightarrow W_L,
 \qquad
 S_R\overset\sim\longrightarrow W_R.
\]
The central sign ambiguity of the spin lift multiplies both maps by the same
sign and cannot exchange the summands.
\end{theorem}

\begin{proof}
Both factor maps into the common $SU(2)$ model are generated group
isomorphisms, so the graph projection is a group isomorphism.  Every
$SU(2)$ automorphism is inner and preserves the Cartan three-form orientation.
Equivalently, on the four-dimensional oriented host an
orientation-preserving orthogonal transport commutes with Hodge star and
therefore preserves the self-dual/anti-self-dual splitting of
$\mathfrak{so}(4)$.  Hence the two spin factors are not exchanged.
\end{proof}


\subsubsection{Finite scalar windows versus geometric carrier levels}
\label{p3b:sec:v180-independent-scalar-window}

\begin{remark}[Independent scalar Fourier-window scope]
\label{p3b:rem:independent-scalar-window-scope}
In the finite window
\[
 \mathcal K_{m,N}=\mathcal K_m\otimes
 P_{|k|\le N}L^2(S^1)\otimes\mathbb C^2,
\]
the Fourier bound $N$ is supplied by the window datum.  The physical compact
momentum result $m_p=1$ does not determine $N$; neither does the growing
position resolution $m_x$ or the finite comparison level
$m_{\rm cmp}=2$.  The heat factor $e^{-\tau n^2}$ is nonzero for every
finite $n$ at $\tau>0$ and therefore is not a sharp replacement for this
selector.  These distinctions preserve the existing action and limit-order
domains.
\end{remark}

\begin{corollary}[Occurrence-conditioned relative maximal-torus axis]
\label{p3b:cor:v191-relative-occurrence-axis}
Let the generated relative readout on the vector factor contain a noncentral
occurring element
\[
 q=ab^{-1}\in SU(2),\qquad q\neq\pm1.
\]
Then $Z(q)$ is the unique maximal torus containing $q$.  Born exchange sends
$q\mapsto q^{-1}$ and therefore preserves this torus because
$Z(q^{-1})=Z(q)$.  Hence the occurrence determines a Born-even unoriented
relative axis.  This is occurrence-conditioned geometry and does not insert a
preferred axis into the unmarked round parent.
\end{corollary}

\begin{remark}[Diagonal/mixed-contraction firewall]
The relative-axis corollary does not identify the independent
$\operatorname{Sp}in(4)_x$ and $\operatorname{Sp}in(4)_p$ vector factors pointwise.  The existing
no-invariant-mixed-contraction result remains active until a licensed
diagonal/same-carrier intertwiner is supplied.
\end{remark}


\subsection{Protected Clifford owner and exceptional-branch exports}
\label{p3b:sec:v236-exceptional-independence}
\begin{proposition}[Protected octet precedes exceptional support]
\label{p3b:prop:v236-protected-octet-independent}
On \((\mathrm{ND}_{\rm sp})\), the protected represented algebra and regular
octet
\[
 A_{\rm prot}\cong\Cl(0,3),\qquad \dim_{\mathbb R}A_{\rm prot}=8
\]
are generated before any exceptional finite-support realization.  A later
same-owner \(2T\), \(2O\) or \(2I\) branch is therefore an action/support on
this already constituted owner and is not the provenance of the protected
Clifford algebra.
\end{proposition}
\begin{corollary}[Branchwise commutator exports]
\label{p3b:cor:v236-exceptional-commutator-exports}
On an occurring same-owner exceptional branch, the intrinsic defect sequence
is generated group-word data:
\[
 2T\to Q_8\to\{\pm1\}\to1,\qquad
 2O\to2T\to Q_8\to\{\pm1\}\to1,\qquad
 2I\to2I.
\]
These are branch-specific dynamical exports; none selects which exceptional
branch occurs.
\end{corollary}


\subsection{Complete quadratic M8 parent and Clifford kernel}
\label{p3b:sec:v189-complete-quadratic-parent}
Let
\[
 V:=\operatorname{im}\mathcal E_{M8},\qquad \dim V=4.
\]
The complete quadratic parent is the tensor square $V\otimes V$, not its
Clifford image alone.

\begin{theorem}[Quadratic sibling exact sequence]
\label{p3b:thm:v189-quadratic-sibling-exact-sequence}
With $c_2(v\otimes w)=vw$ the degree-two Clifford multiplication,
\[
 vw=g(v,w)+v\wedge w
\]
and therefore
\[
 \boxed{0\longrightarrow S^2_0V
 \longrightarrow V\otimes V
 \xrightarrow{\ c_2\ }
 \mathbb R\oplus\Lambda^2V
 \longrightarrow0.}
\]
Equivalently, $16=9+7$.  The symmetric-traceless branch is exactly the kernel
of Cliffordization, while the trace plus bivector branch is its image.
Consequently the later metric/shear and Clifford/curvature-hierarchy roles are
sibling projections of one $M_8$ parent and must not be reconstructed by a
late cross-descendant map.
\end{theorem}

\begin{proof}
Apply Paper~I-C, Theorem~\ref{p1b:thm:quadratic}, to the four-dimensional
relative solder module \(V\).  Its Clifford kernel and image are exactly
the two displayed sibling projections used here.
\end{proof}

\begin{corollary}[Quadratic fibre-product reconstruction]
\label{p3b:cor:v189-quadratic-fibre-product}
By Paper~I-C, Theorem~\ref{p1b:thm:quadratic}, the full tensor square is
recovered from its symmetric and Clifford-visible readouts by matching their
common trace scalar:
\[
 \boxed{V\otimes V\cong
 \operatorname{Sym}^2V\times_{\mathbb R}
 (\mathbb R\oplus\Lambda^2V).}
\]
This is a parent correspondence in the sense of Paper~II-A; it does not
identify $S^2_0V$ with any particular physical metric sector before the local
causal descent of Paper~III-A.
\end{corollary}
\begin{proof}
For a compatible pair $(S,(a,\beta))$ with $\operatorname{tr}_gS=a$, set
\[
 T=S+\iota(\beta),\qquad
 \iota(v\wedge w)=\tfrac12(v\otimes w-w\otimes v).
\]
Its symmetric projection is $S$, exterior projection is $\beta$, and trace
is $a$. Conversely these projections recover every tensor by its symmetric
and antisymmetric decomposition. This is a two-sided linear inverse, so
every compatible pair has a unique parent lift. The displayed equality is
therefore a genuine fibre-product reconstruction, not an inference from the
two comparison maps alone.
\end{proof}


\section{Finite Petersen spectral comparison and the first icosahedral moment defect}
\label{p3b:sec:v319-petersen-spectral}
On an actualized embedded $2I$ branch write $u_a$, $a=1,\ldots,10$,
for the ten unoriented unit Hopf axes, and let $W_P$ be the separate
fifteen-edge real Petersen module.  Neither is, by itself, a new
Dirac operator in the uninserted $\Dcl$ spectral-action class.

\begin{theorem}[Exact even-axis moments through the first anisotropy]
\label{p3b:thm:v319-petersen-moment-six}
For $v\in\operatorname{Im}\mathbb H\simeq\mathbb R^3$ define
$M_{2r}(v)=\sum_{a=1}^{10}\langle u_a,v\rangle^{2r}$.  Then
\[
 M_2(v)=\frac{10}{3}|v|^2,\qquad
 M_4(v)=2|v|^4,\qquad
 M_6(v)=\frac{10}{7}|v|^6+H_6(v),
\]
where $H_6$ is a nonzero homogeneous harmonic polynomial of degree
six.  In particular $\Delta H_6=0$ and, for any Hopf axis
$u_a$, $H_6(u_a)=160/1701$.  Uniform averaging over the complete
unoriented ten-axis orbit therefore agrees with rotational averaging
for quadratic and quartic directional polynomials, but not for all
sextic directional polynomials.
\end{theorem}
\begin{proof}
The ten axes are the projectivized twenty-vertex dodecahedral orbit
of the ordinary icosahedral rotation representation.  The only
invariant harmonic polynomials in degrees one through five are zero;
hence the normalized orbit measure has the same moments as the sphere
through degree five.  The spherical even moment for one direction is
$|v|^{2r}/(2r+1)$, proving the first two identities and the isotropic
sextic part.  The exceptional-axis coordinate substitution from Paper~I-B gives
$M_6(u_a)=370/243$; subtraction of $10/7$ yields $160/1701$.
An explicit polynomial computation verifies $\Delta H_6=0$ and
nonvanishing; the reproducible certificate is Paper~III-G,
Section~\ref{p3g:sec:v319-petersen-certificate}.
\end{proof}

\begin{proposition}[Complete invariant self-adjoint operator family]
\label{p3b:prop:v319-petersen-four-coefficients}
The exact edge resolution in Paper~I-B is
$W_P=\mathbf4\oplus\mathbf5\oplus\mathbf3\oplus\mathbf3'$
with mutually orthogonal projectors $P_4,P_5,P_3,P_{3'}$.
Every real self-adjoint $A_5$-equivariant operator on this specific
real module has the form
\[
 B_P=b_4P_4+b_5P_5+b_3P_3+b_{3'}P_{3'},\qquad b_r\in\mathbb R.
\]
Accordingly an auxiliary, separately licensed finite spectral trace is
\[
 \operatorname{Tr}_{W_P}f(B_P^2/\Lambda^2)
 =4f(b_4^2/\Lambda^2)+5f(b_5^2/\Lambda^2)
  +3f(b_3^2/\Lambda^2)+3f(b_{3'}^2/\Lambda^2).
\]
The four coefficients are not supplied by the Petersen geometry, and
this trace is not a contribution to the actual physical spectral action
unless the whole finite operator is obtained from the licensed source
constructor with its measure, normalization and field incidence.
\end{proposition}
\begin{proof}
The four real $A_5$ irreducibles are pairwise nonisomorphic and of
real type.  The real self-adjoint commutant is therefore scalar on
each summand.  Finite functional calculus and the multiplicities
$4,5,3,3$ give the trace.  The action-incidence qualification is
Theorem~\ref{p3:thm:no_residual_regularity_D1} and its no-insertion
condition.
\end{proof}


\section{Exact scalar Galerkin actions on finite harmonic bands}


\subsection{An exact degree-four scalar interaction and dissipative Galerkin calculus}
\label{p3b:sec:cc17-720-quartic-compiler}
\begin{theorem}[Exact finite quartic energy, derivative and gradient flow]
\label{p3b:thm:cc17-exact-scalar-quartic-galerkin}
Let $\mathcal P_4(S^3(R))$ be the real scalar harmonic band through
coordinate degree four (dimension $1+4+9+16+25=55$).
Let $S:\mathcal P_4\to\mathbb R^{720}$ be the isometric evaluation
$(Sf)_a=\sqrt{w_a}f(z_a)$ at the positive 720-node
vertex--facet rule of Theorem~\ref{p1b:thm:cc17-positive-vertex-facet-19},
and let $P_4=SS^*$ and $L_4$ denote the corresponding finite scalar
Laplacian. For fixed comparison coefficients $m^2\ge0,\lambda\ge0$
define on the 55-dimensional band
\[
 \mathcal E(f)=\tfrac12\langle f,(-\Delta+m^2)f\rangle
      +\tfrac\lambda4\int_{S^3(R)}f(q)^4\,d\sigma(q).
\]
Writing $x=Sf$ and $w=(w_a)$ gives the \emph{exact}, finite expression
\[
\boxed{\mathcal E_S(x)=\tfrac12x^T(L_4+m^2P_4)x
    +\tfrac\lambda4\sum_{a=1}^{720}\frac{x_a^4}{w_a}.}
\]
Its Hilbert-space gradient within the sample image is
\[
\boxed{\nabla\mathcal E_S(x)=
    (L_4+m^2P_4)x+\lambda P_4(x^{\odot3}\oslash w).}
\]
Therefore the 720-coordinate constrained ODE
$\dot x=-\nabla\mathcal E_S(x)$ is \emph{exactly} the degree-four
Galerkin projection of the corresponding continuous scalar gradient
flow, and it obeys
\[\boxed{\frac{d}{dt}\mathcal E_S(x(t))
       =-\|\dot x(t)\|_2^2\le0.}\]
Here $\oslash$ means componentwise division and $\odot$ denotes
componentwise multiplication. No claim that the full unprojected
continuum nonlinear flow remains band-limited is made.
\end{theorem}
\begin{proof}
For $f\in\mathcal P_4$ the quartic polynomial $f^4$ has coordinate
degree at most $16<20$; the 720-point cubature therefore integrates
it exactly. Because $x_a=\sqrt{w_a}f(z_a)$, its weighted fourth
power is $w_af(z_a)^4=x_a^4/w_a$. The scalar Laplacian part is
exact by the harmonic spectral projector theorem.
Differentiate the finite expression in the Euclidean sample
coordinates and orthogonally project the result onto $\operatorname{Ran}P_4$.
For any $h\in\mathcal P_4$, the term $hf^3$ has degree at most 16,
so the finite gradient agrees with the continuum weak Galerkin
gradient in every direction. The energy dissipation identity is
the standard finite-dimensional chain rule for a constrained
Hilbert-space gradient flow.
\end{proof}

\begin{corollary}[Higher-than-degree-two exact scalar interaction ports]
\label{p3b:cor:cc17-scalar-vertices}
The 720-point rule integrates every scalar mixed $r$-field product
with input maximum coordinate degrees $a_1,\ldots,a_r$ satisfying
$\sum_j a_j\le19$; its projected-product readout is exact whenever
an additional test-field degree is included in that sum.
In particular every quartic vertex on $\mathcal P_4$ and every
cubic vertex on $\mathcal P_6$ is exact. The 120-node uniform $2I$
rule alone gives complete scalar quartic exactness only through
$\mathcal P_2$.
\end{corollary}
\begin{proof}
The degree bound is the sum of the coordinate polynomial degrees.
Apply the positive nineteen-design identity to the product;
weak Galerkin identities follow by testing against another band field.
The prior 120-node cutoff is the source's sharp uniform-$2I$
complete-band frontier.
\end{proof}

\begin{remark}
The comparison scalar action here is independently \emph{defined} and
not identified with the native, physical full spectral/Einstein,
Higgs or QCD action. No coupling $\lambda$, mass $m$ or physical
operator owner is derived by these purely mathematical identities.
\end{remark}


\section{Higher-band scalar spectral-action comparisons and typed response transport}
\label{p3b:sec:cc20-band30-math-consumer}
\begin{corollary}[Finite positive sampling of scalar spectral and polynomial actions]
\label{p3b:cor:cc20-band30-action-comparison}
On the round scalar $S^3(R)$, take the positive sixty-thousand-node
cubature of Theorem~\ref{p1b:thm:cc112-degree61-positive}, and
its isometric degree-thirty evaluation map of
Theorem~\ref{p1c:thm:cc112-band30-spectral-projectors}.
For any function $F$ defined on the finite set
$\{n(n+2)/R^2:0\le n\le30\}$, its scalar-band trace is
\[
 \boxed{\operatorname{Tr}_{\mathcal P_{30}} F(-\Delta)
 =\sum_{n=0}^{30}(n+1)^2F(n(n+2)/R^2).}
\]
If $f\in\mathcal P_{15}(S^3)$, every quartic polynomial scalar
potential and its band-projected first variation and Hessian are
computed without quadrature error from these nodes, since each
associated test integrand has polynomial degree at most sixty.
The conclusion concerns only the chosen scalar finite-dimensional
Galerkin family; no Dirac, gauge, QCD or complete physical action is
identified with it.
\end{corollary}
\begin{proof}
The degree-sixty-one rule gives $S^*S=I$ on the degree-thirty
harmonic band and exactness for products through degree sixty-one.
The round scalar Laplacian eigenvalue $n(n+2)/R^2$ has multiplicity
$(n+1)^2$.  The quartic potential and its first two derivatives
involve at most four factors of degree at most fifteen, so the
cubature evaluates their Galerkin integrals exactly.
\end{proof}


\subsection{Same-owner finite-response transport and the Hurwitz comparison firewall}
\label{p3b:subsec:s3-owner-transport-v1190}

\begin{proposition}[Exact two-jet under a typed configuration transport]
\label{p3b:prop:s3-full-owner-pullback-v1190}
Let $\mathcal U_H$ be a finite Hurwitz-link configuration chart and
$\mathcal Q_x$ the position-action configuration domain of the existing
Paper~III-B finite window.  Suppose an actual $C^2$ chart map
$T:\mathcal U_H\to\mathcal Q_x$ and a $C^2$ generated action
$S_{\rm Born}:\mathcal Q_x\to\mathbb R$ are supplied on the \emph{same}
action owner, including the source/localizer, cutoff and trace density.
Then for link tangents $u,v$,
\begin{equation}\label{p3b:eq:s3-full-owner-twojet-v1190}
\boxed{D^2(S_{\rm Born}\circ T)[u,v]
=D^2S_{\rm Born}[DTu,DTv]
 +DS_{\rm Born}[D^2T[u,v]].}
\end{equation}
If the finite position window and the momentum seed are held fixed,
its spectral part is computed by the exact Duhamel and contact formula of
Proposition~\ref{p3b:prop:v245-exact-finite-heat-two-jet}.
If any carrier, localizer, cutoff, source, or measure moves, its variation
must instead be included in the \emph{complete} occurring action.
\end{proposition}
\begin{proof}
Apply the second Fr\'echet chain rule to the actual $C^2$ maps.
The fixed-window spectral assertion is precisely the scope of
Proposition~\ref{p3b:prop:v245-exact-finite-heat-two-jet}; it makes no
claim about an unspecified position-to-momentum map.
\end{proof}

\begin{remark}[Comparison Hessian versus generated physical response]
\label{p3b:rem:s3-comparison-owner-firewall-v1190}
The exact rank-eight quadratic-blind sector of the massless
24-cell momentum-shell spin \emph{comparison} is documented with finite
integer replay in the S$^3$--DAG companion and Paper~VII-D.
It is not identified with a physical Hessian of the generated
position-sector operator $D_x$ merely by shared S$^3$ incidence.
A physical claim needs an actual same-owner transport satisfying
\eqref{p3b:eq:s3-full-owner-twojet-v1190}, or an independently generated
physical action on that link chart.  Vanishing of a comparison Hessian
is not evidence that an operator first variation or nonlinear response
vanishes.  Nothing here modifies the upstream constitutional or
sectorwise-additivity results.
\end{remark}

\section{Conclusion and scope of the common spectral action}
\label{paper3b:sec:organized-conclusion}
\label{p3b:sec:conclusion}
Paper~III-B begins from the already constituted momentum seed
\[
 \mathcal H_p\cong\mathcal H_1,
 \qquad
 m_p=1,
\]
and from the reciprocal role geometry of Paper~III-A, derived by the
premise-minimal Paper-II-C/GP2 reciprocal-role objectification and
relation-free composition rather than by metric-curry or independent BR/IND
assumptions.  It does not reselect the compact seed.  Under
\((\mathrm{CAL}_{xp})\) the physical configuration radius is
\(\Rcs=\alpha_{xp}\Rmom\); under \((\mathrm{CAL}_{\hbar})\) the primitive
scalar phase class has its \(2\pi\hbar\) representative; and
\((\mathrm{CAL}_{t})\) supplies the physical duration \(T_t\) of one already
generated positive primitive turn.  These are calibrations, not structural
selectors.  A separate scalar-circle spectral metric is not introduced in
the constitutional core.

The relative and seven-dimensional readouts are generated by conservative
graph lifts on the complete relation-free product.  Their kernel pairs are
exactly the stated diagonal orbit relations, but the product and all factor
tests remain present.  Thus no duplicate occurrence assumption is needed for
the readout roles, while a local base, nonzero connection, fibre metric,
reduction measure, generator normalization, and \((\mathrm{GRED}_7)\) remain
outside the theorem.

Inside the additive Born central trace class, the reciprocal response is
\[
 \frac12(\chi+\chi^{-1})-1=\cosh(\log\chi)-1,
\]
with the physical basepoint \(\chi_{\rm phys}=1\) fixed upstream.  The
spectral-action class \(\Dcl\) is the derived form of an explicit closed
single-process occurrence; its spectral-orbit and no-channel-insertion
results classify that form.  Its occurring profile is not retyped as the
complete physical reciprocal-position test frontier; the later
\(\mathsf{SF}_\Lambda\) occurrence is a separate operational interface.
Momentum dynamics remains an exact finite \(\mathcal H_1\) trace.
Position-side large-window asymptotics may be used only under the separately
constituted position interfaces and the controlled exhaustion hypotheses.
All such physical occurrences, trajectories, realized action profiles, and
calibrations are typed by G0; G0 does not choose the process law or replace
the analytic hypotheses.

Absolute Paper-I-A coordinate values and completion-dependent Clifford scales
still specialize to \(\mathsf C_{\rm PI}\), while the normalized protected
even algebra requires only \((\mathrm{ND}_{\rm sp})\).  Finite-support
results use an actualized \(\mathsf{Supp}_{2T}\) role.  The local
Born--Liouville identity does not identify semiclassical phase-volume counts
with the representation dimension of \(\mathcal H_1\).  The canonical scalar
map \(P_0\) is derived here as the unique allowed parity-even scalar projector;
its physical terminal occurrence is
\(\mathrm{DERIVED\text{-}LATER}\) in Paper~III-C.  The generic active-density
family remains forward synthesis at the III-B stage, while the particular
hierarchy-active state on the independently occurring \((\mathrm{MRESP})\)
branch is \(\mathrm{DERIVED\text{-}LATER}\) in Paper~III-E.  Its separate
state-to-metric-Hessian coupling remains an explicit later response-class
law.  Other nonlinear hierarchy, index-locking, and gravitational formulas
retain their own stated antecedents.  There is no universal selection rule: a
unique nonzero closed least element is retained only under a generated
role-faithful GP3 comparison, and endpoint saturation only under the
Paper-II-A exact-fidelity criterion.  Later discharge never rewrites an
earlier unresolved family as though it had been unique at the III-B stage.

Finally, the generated diagonal ratio fibre has an orientation-preserving
round-sphere ancestry map.  Its spin lift preserves the two \(\operatorname{Sp}in(4)\)
simple factors and induces \(S_L\simeq W_L\), \(S_R\simeq W_R\); the common
central sign cannot exchange them.  Thus the relative even chirality ancestry
exported to Papers~III-E and V-A is theorem-level, while an odd bridge would
require an additional one-sided orientation reversal.


\printbibliography
\end{document}
